% Options for packages loaded elsewhere
\PassOptionsToPackage{unicode}{hyperref}
\PassOptionsToPackage{hyphens}{url}
%
\documentclass[
]{article}
\usepackage{amsmath,amssymb}
\usepackage{iftex}
\ifPDFTeX
  \usepackage[T1]{fontenc}
  \usepackage[utf8]{inputenc}
  \usepackage{textcomp} % provide euro and other symbols
\else % if luatex or xetex
  \usepackage{unicode-math} % this also loads fontspec
  \defaultfontfeatures{Scale=MatchLowercase}
  \defaultfontfeatures[\rmfamily]{Ligatures=TeX,Scale=1}
\fi
\usepackage{lmodern}
\ifPDFTeX\else
  % xetex/luatex font selection
\fi
% Use upquote if available, for straight quotes in verbatim environments
\IfFileExists{upquote.sty}{\usepackage{upquote}}{}
\IfFileExists{microtype.sty}{% use microtype if available
  \usepackage[]{microtype}
  \UseMicrotypeSet[protrusion]{basicmath} % disable protrusion for tt fonts
}{}
\makeatletter
\@ifundefined{KOMAClassName}{% if non-KOMA class
  \IfFileExists{parskip.sty}{%
    \usepackage{parskip}
  }{% else
    \setlength{\parindent}{0pt}
    \setlength{\parskip}{6pt plus 2pt minus 1pt}}
}{% if KOMA class
  \KOMAoptions{parskip=half}}
\makeatother
\usepackage{xcolor}
\usepackage{longtable,booktabs,array}
\usepackage{calc} % for calculating minipage widths
% Correct order of tables after \paragraph or \subparagraph
\usepackage{etoolbox}
\makeatletter
\patchcmd\longtable{\par}{\if@noskipsec\mbox{}\fi\par}{}{}
\makeatother
% Allow footnotes in longtable head/foot
\IfFileExists{footnotehyper.sty}{\usepackage{footnotehyper}}{\usepackage{footnote}}
\makesavenoteenv{longtable}
\setlength{\emergencystretch}{3em} % prevent overfull lines
\providecommand{\tightlist}{%
  \setlength{\itemsep}{0pt}\setlength{\parskip}{0pt}}
\setcounter{secnumdepth}{-\maxdimen} % remove section numbering
% unicode-fix.tex — render raw-Unicode symbols that the Latin Modern *text* font
% lacks, using Computer Modern math/text equivalents. This keeps the established
% Computer Modern look while filling the glyphs xelatex would otherwise drop.
% Included via: pandoc ... --include-in-header=unicode-fix.tex
\usepackage{newunicodechar}
\usepackage{pifont}
\usepackage{graphicx} % unicode-math's smallsetminus uses rotatebox
% Use the Computer Modern-compatible set-difference construction; the large
% Unicode glyph selected by TeX Live 2022 is absent from Latin Modern Math.
\AtBeginDocument{\let\setminus\smallsetminus}
% --- Greek letters -> math (Latin Modern Math; consistent with CM body text) ---
\newunicodechar{α}{\ensuremath{\alpha}}
\newunicodechar{β}{\ensuremath{\beta}}
\newunicodechar{γ}{\ensuremath{\gamma}}
\newunicodechar{δ}{\ensuremath{\delta}}
\newunicodechar{ε}{\ensuremath{\varepsilon}}
\newunicodechar{θ}{\ensuremath{\theta}}
\newunicodechar{μ}{\ensuremath{\mu}}
\newunicodechar{ν}{\ensuremath{\nu}}
\newunicodechar{ρ}{\ensuremath{\rho}}
\newunicodechar{σ}{\ensuremath{\sigma}}
\newunicodechar{τ}{\ensuremath{\tau}}
\newunicodechar{φ}{\ensuremath{\varphi}}
\newunicodechar{χ}{\ensuremath{\chi}}
\newunicodechar{ψ}{\ensuremath{\psi}}
\newunicodechar{η}{\ensuremath{\eta}}
\newunicodechar{κ}{\ensuremath{\kappa}}
\newunicodechar{λ}{\ensuremath{\lambda}}
\newunicodechar{π}{\ensuremath{\pi}}
\newunicodechar{ω}{\ensuremath{\omega}}
\newunicodechar{Ḣ}{\ensuremath{\dot H}}
\newunicodechar{≪}{\ensuremath{\ll}}
\newunicodechar{≫}{\ensuremath{\gg}}
\newunicodechar{⟺}{\ensuremath{\Longleftrightarrow}}
\newunicodechar{⟹}{\ensuremath{\Longrightarrow}}
\newunicodechar{⇔}{\ensuremath{\Leftrightarrow}}
\newunicodechar{⟨}{\ensuremath{\langle}}
\newunicodechar{⟩}{\ensuremath{\rangle}}
\newunicodechar{⁰}{\textsuperscript{0}}
\newunicodechar{¹}{\textsuperscript{1}}
\newunicodechar{²}{\textsuperscript{2}}
\newunicodechar{³}{\textsuperscript{3}}
\newunicodechar{⁴}{\textsuperscript{4}}
\newunicodechar{⁵}{\textsuperscript{5}}
\newunicodechar{⁶}{\textsuperscript{6}}
\newunicodechar{⁷}{\textsuperscript{7}}
\newunicodechar{⁸}{\textsuperscript{8}}
\newunicodechar{⁹}{\textsuperscript{9}}
% --- Relations / operators -> math ---
\newunicodechar{↔}{\ensuremath{\leftrightarrow}}
\newunicodechar{⇐}{\ensuremath{\Leftarrow}}
\newunicodechar{⇒}{\ensuremath{\Rightarrow}}
\newunicodechar{∂}{\ensuremath{\partial}}
\newunicodechar{∅}{\ensuremath{\varnothing}}
\newunicodechar{∈}{\ensuremath{\in}}
\newunicodechar{∼}{\ensuremath{\sim}}
\newunicodechar{≈}{\ensuremath{\approx}}
\newunicodechar{≠}{\ensuremath{\neq}}
\newunicodechar{≤}{\ensuremath{\leq}}
\newunicodechar{≥}{\ensuremath{\geq}}
\newunicodechar{⊂}{\ensuremath{\subset}}
\newunicodechar{⊃}{\ensuremath{\supset}}
% --- Special ---
\newunicodechar{ℏ}{\ensuremath{\hbar}}
\newunicodechar{ℤ}{\ensuremath{\mathbb{Z}}}
% --- Superscript modifier letters (footnote-style markers) -> text superscript ---
\newunicodechar{ʰ}{\textsuperscript{h}}
\newunicodechar{ʲ}{\textsuperscript{j}}
\newunicodechar{ᵃ}{\textsuperscript{a}}
\newunicodechar{ᵇ}{\textsuperscript{b}}
\newunicodechar{ᵈ}{\textsuperscript{d}}
\newunicodechar{ᵉ}{\textsuperscript{e}}
\newunicodechar{ᵍ}{\textsuperscript{g}}
\newunicodechar{ᶠ}{\textsuperscript{f}}
\newunicodechar{⁻}{\textsuperscript{\ensuremath{-}}}
% --- Subscript digits -> text subscript ---
\newunicodechar{₀}{\textsubscript{0}}
\newunicodechar{₁}{\textsubscript{1}}
\newunicodechar{₂}{\textsubscript{2}}
\newunicodechar{₃}{\textsubscript{3}}
\newunicodechar{₄}{\textsubscript{4}}
\newunicodechar{₈}{\textsubscript{8}}
% --- Checkmarks ---
\newunicodechar{✓}{\ding{51}}
\newunicodechar{✗}{\ding{55}}

% --- added 2026-08-02 after an independent glyph audit ---
\newunicodechar{ℓ}{\ensuremath{\ell}}
\newunicodechar{⊕}{\ensuremath{\oplus}}
\newunicodechar{⊗}{\ensuremath{\otimes}}
\newunicodechar{∠}{\ensuremath{\angle}}
\newunicodechar{≡}{\ensuremath{\equiv}}
\newunicodechar{∝}{\ensuremath{\propto}}
\newunicodechar{ℝ}{\ensuremath{\mathbb{R}}}
\newunicodechar{₅}{\ensuremath{_5}}
\newunicodechar{₋}{\ensuremath{_-}}
\newunicodechar{⁺}{\ensuremath{^+}}
\newunicodechar{⁽}{\ensuremath{^(}}
\newunicodechar{⁾}{\ensuremath{^)}}
% U+239B..U+23A0: multi-line parenthesis assembly pieces emitted by some
% converters. There is no sane inline equivalent; map to ordinary delimiters so
% the notation survives rather than vanishing silently.
\newunicodechar{⎛}{\ensuremath{\big(}}
\newunicodechar{⎜}{\ensuremath{\big\vert}}
\newunicodechar{⎝}{\ensuremath{\big(}}
\newunicodechar{⎞}{\ensuremath{\big)}}
\newunicodechar{⎟}{\ensuremath{\big\vert}}
\newunicodechar{⎠}{\ensuremath{\big)}}
\ifLuaTeX
  \usepackage{selnolig}  % disable illegal ligatures
\fi
\IfFileExists{bookmark.sty}{\usepackage{bookmark}}{\usepackage{hyperref}}
\IfFileExists{xurl.sty}{\usepackage{xurl}}{} % add URL line breaks if available
\urlstyle{same}
\hypersetup{
  hidelinks,
  pdfcreator={LaTeX via pandoc}}

\author{}
\date{}

\begin{document}

\hypertarget{ux210f-the-bekenstein-hawking-entropy-and-the-running-vacuum-form-of-dark-energy-from-the-cosmological-horizon}{%
\section{ℏ, the Bekenstein-Hawking Entropy, and the Running-Vacuum Form
of Dark Energy from the Cosmological
Horizon}\label{ux210f-the-bekenstein-hawking-entropy-and-the-running-vacuum-form-of-dark-energy-from-the-cosmological-horizon}}

\textbf{Author:} Alex Maybaum\\
\textbf{Date:} April 2026\\
\textbf{Status:} DRAFT PRE-PRINT\\
\textbf{Classification:} Theoretical Physics / Gravitation / Cosmology

\begin{center}\rule{0.5\linewidth}{0.5pt}\end{center}

\hypertarget{abstract}{%
\subsection{Abstract}\label{abstract}}

The exact de Sitter horizon supplies a candidate causal partition for
the embedded-observation framework of {[}Main{]}. Under H-count, H-slope
on identified transitions, the horizon and frame conditions of §§2 and
8.5, local boost structure and symmetric detailed-balance rates, the
thermal calibration gives \(\hbar=c^3\epsilon^2/(4G\sigma)\) and
\(\epsilon=2\sqrt{\sigma}\,l_p\), with \(\sigma\) the entropy per
counted boundary mode. The calibrated count \(\sigma A/\epsilon^2\)
gives \(S=A/(4l_p^2)\) through the normalization ratio \(2\pi/8\pi\).
The lattice application uses the spherical or equal-density scope of §5.
The Kerr-area inference from GW250114 supports the classical area
theorem {[}5{]}. An entropy measurement would address the separate
\(1/4\) coefficient. For the designated trace-out, the emergent vacuum
energy and the supplied Friedmann-scale density belong to different
descriptions. The covariant field equations, sourcing, Bianchi structure
and continuum limit remain open at G3. The Type II running-vacuum
structural reading assumes conserved matter and running \(G\), with a
near-null estimate conditional on the source split. The Hubble-volume
average of the horizon thermal energy equals the supplied critical
density in three spatial dimensions. Its interpretation as a source for
the nonbaryonic budget requires a matter coupling and the baryon
fraction. With the de Sitter temperature, the open heat input
\(Q=M_Bc^2\), and the factor \(1/6\) imported from Verlinde, entropy
displacement gives the conditional scale \(a_0=cH/6\) and the deep-MOND
Tully-Fisher relation. The outer declines reported by Jiao et al.~and
Genzel et al.~are in tension with this law, with a quantitative
comparison requiring the systems' baryon profiles and kinematics (§7.3).

\begin{center}\rule{0.5\linewidth}{0.5pt}\end{center}

\hypertarget{introduction}{%
\subsection{1. Introduction}\label{introduction}}

The cosmological observables that most resist explanation within
standard physics --- the origin of the action scale \(\hbar\), the
Bekenstein-Hawking area law, the \(10^{122}\) cosmological constant
problem, the accelerated expansion, the \textasciitilde95\% dark sector,
and the MOND-like behavior of galactic rotation curves --- are
considered using the cosmological event horizon as a candidate
visible--hidden partition. The thermodynamic comparisons and the
physical application of that partition carry the conditions below.

\textbf{Prerequisites.} This paper takes as established the following
results from a companion work {[}Main{]}: (i) an observer embedded in a
deterministic bijection on a finite configuration space, coupled to a
hidden sector that is non-trivial (C1), record-persistent (C2), of
sufficient memory capacity (C3 --- with \(N_H \gg N_V\) the comfortable
regime), and read back (C4), has a reduced visible-sector description
that is --- in the memory-bearing sector that C1, C3 and C4 characterize
--- non-Markovian, and that universally admits a unitarily evolving
fixed-basis quantum representation --- the representation universal over
finite laws (\(S \iff D \iff Q_{\mathrm{fb}}\)), the necessity holding
in the memory currency, with accessible non-Markovianity --- the
framework's finite-horizon visible-law memory notion, not generic
quantum non-Markovianity --- obtaining exactly when, for some accessible
order \(k\), every finite accessible horizon \(K \geq k+1\) admits a
deterministic realization with C1, C3 and C4, the readback at \(k\) (a
realization with C4 forces the memory by the readback lemma of {[}Main
§2.3{]}; accessible memory yields such a realization by the
finite-horizon process dilation of {[}Main §3.4{]}; both directions
proved in prose), C2 lying outside that equivalence and the slow-bath
form of its necessity for physical hidden sectors conditional on the
hidden-sector mixing hypothesis, which the eigenstate thermalization
hypothesis motivates, its persistence form holding without it, and full
operational quantum theory approached at every finite test with its
unique selection open; (ii) under the canonical measure of {[}Main{]}
Lemma 3 --- the invariant measure selected by maximal entropy --- the
emergent description is uniquely determined by the partition
(partition-relativity); (iii) the emergent transition probabilities
\(T_{ij}(t)\) determine the Bohr-frequency set of the emergent
Hamiltonian and, under the genericity hypotheses (G1)--(G4) of the
phase-locking lemma or, without (G4), conditional on the classification
hypothesis H-BG of §3.3, the emergent Hamiltonian itself up to a
diagonal-unitary D-gauge, an energy origin, and one global antiunitary
conjugation --- the two-branch form of §3.3 --- via the phase-locking
lemma applied to continuous-time sampling. The technical bridge from
P-indivisible classical stochastic processes on finite configuration
spaces to unitary quantum evolution on finite Hilbert spaces is
Barandes' stochastic-quantum correspondence {[}35, 36{]}, used by the
companion paper for compression and interpretation --- the fixed-basis
representation itself is supplied internally there
(\(S \iff D \iff Q_{\mathrm{fb}}\)). We treat these results as inputs
and apply them to the cosmological horizon. For the cosmological
application the Einstein constraints motivate a candidate C1 coupling,
the Hubble-time mixing scale is an input for C2, and the boundary mode
count gives a C3 capacity estimate for the chosen laboratory demand. The
physical horizon application requires the realization of accessible
coupling, persistence and capacity, and remains conditional on C4
(§2.2).

The present paper reproves none of {[}Main{]}'s foundational results.
Where a theorem below invokes a specific input --- partition-relativity
(§3.1), the emergent-description determination by transition
probabilities (§3.2--§3.3), or the P-indivisibility/recurrence argument
(§3.2 deep-sector lemma; §8 theorem requirements) --- the invocation is
stated as a known result and the reader is directed to {[}Main{]} for
the proof.

The action-scale calibration matches the classical shell and emergent
frequency slopes of the same operational transitions on exact de Sitter
(§§2.3, 3.2). Under H-count, H-slope, the horizon and frame conditions,
local boost structure and the symmetric detailed-balance ratio, it gives
\(\hbar=c^3\epsilon^2/(4G\sigma)\). The Planck-length relation then
fixes \(\epsilon=2\sqrt{\sigma}\,l_p\) and the calibrated area-law
coefficient is \(2\pi/8\pi=1/4\). The boundary count uses a specified
mode density (§5). GW250114 provides observational evidence for the
classical area increase inferred from the waveform {[}5{]}. Section 6
distinguishes the emergent vacuum energy from the supplied cosmological
energy scale. Dynamical vacuum-offset decoupling requires the coupling
open at G3 (§6.2).

Section 7 gives empirical comparisons for the gravitational model. The
Type II running-vacuum channel is a structural reading of the boundary
bookkeeping, with conserved matter and running \(G\) in that
implementation. The smooth local expansion is shared by other models.
The emergent-QFT estimate \(|\nu|\sim10^{-32}\) is conditional on its
source attribution and the stated split. The separate Bertini et al.~fit
{[}15{]} is compatible with zero. The Hubble-volume mean of the boundary
thermal energy supplies a normalization. The proposed boundary
interpretation of \(1-\Omega_B\) needs the measured baryon fraction and
the open matter coupling. Conditional on the de Sitter temperature and
\(Q=M_Bc^2\), with the Verlinde factor \(1/6\), the deep-modified regime
has \(a_0=cH/6\) and the crossover scale \(r_M=\sqrt{6GM_B/(cH)}\). The
Milky Way decline reported by Jiao et al.~over \(19\)--\(26.5\) kpc lies
beyond the quoted \(9\)--\(16\) kpc crossover. The declining outer
velocities Genzel et al.~report at \(z=0.9\)--\(2.4\) also differ from
the trend at fixed baryonic profile as \(H(z)\) increases. A
quantitative comparison requires the galaxies' baryon profiles and
pressure support (§7.3).

Three epistemic tiers are worth distinguishing before the technical
development begins. \textbf{Tier 1} comprises the normalization
relations on the stated classical and counting inputs: the classical
horizon temperature of the boundary count, the finite horizon counting,
and the structural two-level distinction between the emergent vacuum
energy and the effective cosmological constant. The results built on
these carry named conditions: \(\hbar = c^3\epsilon^2/(4G\sigma)\),
\(\epsilon = 2\sqrt{\sigma}\,l_p\) and the \(1/4\) coefficient are
conditional on H-slope together with H-count, the horizon and frame
conditions (§§2, 5), local boost structure at the horizon and the
symmetric detailed-balance rate ratio (§8.5), and the quoted value
\(\epsilon = 2\,l_p\) additionally on the assumption \(\sigma = 1\) of
one nat per boundary mode (§3.1), of which the \(1/4\) coefficient is
independent; the descriptive-level distinction behind the
cosmological-constant dissolution is established while the dynamical
dissolution is conditional on the G3 coupling theorem (§6.2); and the
Type II running-vacuum functional form is not presently derived, its
local expansion being shared by every smooth model (§7.1). \textbf{Tier
2} comprises the dark-sector phenomenology: the \textasciitilde95\%
invisible matter fraction, the MOND acceleration scale, and the
rotation-curve crossover. The proposed budget interpretation requires
the baryon fraction and a physical coupling. The acceleration scale and
crossover use the Clausius relation and require the heat-to-mass
identification \(Q=M_Bc^2\); they inherit the conditions on the de
Sitter temperature of §3.2, and the dimensional factor that sets the
acceleration scale is imported from Verlinde's emergent-gravity analysis
(§7.3). \textbf{Tier 3} comprises solution-specific quantities --- the
transition steepness between the MOND and Keplerian regimes, the
greybody factor for Hawking radiation, and the specific microscopic
realization of the substratum --- which depend on the particular
bijection \(\varphi\) and are addressed in the companion paper {[}SM{]}.

The tier system maps to the framework's broader hierarchical structure
articulated in {[}Structure §2{]}. \emph{Tier 1} --- the classical
horizon temperature of the boundary count, the finite horizon counting,
and the two-level distinction --- operates at Level C × Level G4: these
relations use the classical first-law and horizon-counting setup and
admit different microscopic representatives. They are
universality-class-level structural results, and the conditional results
built on them (\(\hbar\), \(\epsilon = 2\sqrt{\sigma}\,l_p\), the
\(1/4\) coefficient) are universality-class-level in the same sense once
their named conditions are granted, since those conditions concern the
horizon, the hidden sector's density of states and the emergent dynamics
rather than the lattice (§8.5). \emph{Tier 2} --- the dark-sector
phenomenology --- sits at Level D × Levels G1--G2, with physical
sourcing and the heat input separately required in the cosmological
application. \emph{Tier 3} --- solution-specific quantities --- is at
the bijection-specific level below Level D, depending on the particular
\(\varphi\) within \([(S, \varphi)]/\mathcal{G}_{\text{sub}}\). This
stratification clarifies what a conditional universality-class result
claims: the same calibration formula applies in any realization meeting
the horizon, counting and dynamical conditions, with their physical
fulfillment open for the cosmological cut.

The paper is organized as follows. §2 identifies the cosmological
horizon as a causal partition and states the inputs for its physical
realization. §3 derives the emergent action scale \(\hbar\) from the
classical horizon temperature by thermal self-consistency, the
calibration \(\hbar=\tau_K/\beta_E\) of §2.3; boundary-only dependence
of the transition data is consistent with it, and dimensional
determination checks its form. §4 fixes the discreteness scale
\(\epsilon = 2\sqrt{\sigma}\,l_p\). §5 derives the Bekenstein-Hawking
entropy and discusses the GW250114 classical area comparison. §6 gives
the conditional dissolution of the cosmological constant problem. §7
presents the quantitative predictions for dark energy, the dark sector,
and dark matter, with a summary table in §7.4. §8 discusses scope, the
status of the discreteness scale, vacuum energy and the Higgs potential,
and the framework's falsifiability conditions. §9 concludes. The
cubic-lattice realization of the substratum --- which fixes the Standard
Model gauge structure via the wave equation --- is developed in a
companion paper {[}SM{]} and is independent of the GR results presented
here. The reconstruction theorem and the synthesis claim that quantum
mechanics and general relativity descend as projections of the same
\((S, \varphi)\), with a third projection bearing on the arrow of time,
are developed in {[}Substratum{]}.

\begin{center}\rule{0.5\linewidth}{0.5pt}\end{center}

\hypertarget{the-cosmological-horizon-as-causal-partition}{%
\subsection{2. The Cosmological Horizon as Causal
Partition}\label{the-cosmological-horizon-as-causal-partition}}

\hypertarget{the-partition}{%
\subsubsection{2.1 The partition}\label{the-partition}}

The cosmological event horizon, when present, defines the candidate
causal partition of Lemma 2. For the observer's worldline, \(\Gamma_V\)
is the region in causal reach and \(\Gamma_H\) its complement. This uses
the supplied causal geometry. A physical realization of the
memory-bearing law requires the additional conditions of §2.2.

Different observer worldlines can have different causal pasts. In the
calibrated formula \(\hbar=c^3\epsilon^2/(4G\sigma)\) the area and
surface gravity cancel. Calibrated horizons that share \(c\), \(G\), and
the boundary counting ratio \(\epsilon^2/\sigma\) therefore share the
calibrated action scale (Appendix A.4).

\hypertarget{conditions-of-the-horizon-application}{%
\subsubsection{2.2 Conditions of the horizon
application}\label{conditions-of-the-horizon-application}}

\textbf{(C1)} The Hamiltonian and momentum constraints correlate
interior and exterior initial data on a Cauchy surface. In the ADM
description a horizon boundary term can depend on data on both sides.
This motivates a candidate coupling. C1 requires hidden-to-visible
influence on an accessible conditional {[}Main §1.3{]}, which the
constraint correlations alone do not establish.

\textbf{(C2)} The slow-bath example takes a hidden mixing time of order
the Hubble time, \(\tau_B\sim H^{-1}\sim10^{17}\) s, as a physical
input. For a laboratory timescale \(\tau_S\sim10^{-15}\) s the ratio is
of order \(10^{-32}\) under that assumption.

\textbf{(C3)} With one mode per area \(\epsilon^2\), the boundary at
\(A\sim10^{52}\) m\(^2\) has order \(10^{122}\) modes at the one-nat
calibration. This gives a large capacity relative to a chosen laboratory
demand. The estimate uses the mode count and the laboratory register
under comparison.

\textbf{(C4)} Cosmological history readback would require a
visible-history record written to the hidden sector to be routed back
into future visible conditionals within the accessible window {[}Main
§1.3{]}. \textbf{Not presently discharged.} The action-scale calibration
is carried by H-slope together with the horizon and frame conditions and
H-count, without a readback premise.

\hypertarget{application}{%
\subsubsection{2.3 Application}\label{application}}

\textbf{Which horizon, and why the calibration is stated on exact de
Sitter.} Two distinct surfaces are in play and they coincide only in
exact de Sitter. The causal definition above is the cosmological event
horizon, \(R_{\rm event}=a(t)\int_t^\infty c\,dt'/a(t')\); the area
\(A=4\pi c^2/H^2\) used below is that of the Hubble/apparent horizon
\(R_A=c/H\). Generic FLRW has \(\dot H\neq0\), no timelike Killing
vector and hence no bifurcate Killing horizon, so exact Euclidean KMS
periodicity --- a stationary result --- is not available there. For the
flat-FLRW apparent horizon the Hayward--Kodama surface gravity gives
\(T_A=\frac{\hbar H}{2\pi k_B}\big|1+\dot H/(2H^2)\big|\), and for flat
\(\Lambda\)CDM with \(\dot H/H^2=-\tfrac32\Omega_m\) this is
\(\approx0.775\,\hbar H/(2\pi k_B)\) today --- a \(22\%\) effect, not
Planck suppressed, and larger than the lattice-dispersion correction
\(O((\epsilon H/c)^2)\sim10^{-122}\) at \(\sigma = 1\) by 121 orders of
magnitude. The two must not be conflated. The gap equation below is
therefore stated as a calibration on an exact de Sitter horizon, where
\(R_{\rm event}=R_A=c/H\), the Killing structure exists, and the
comoving frame is freely falling. Because \(\kappa\) cancels between
\(c^2\epsilon^2\kappa/(8\pi G\sigma)\) and \(\hbar\kappa/(2\pi c)\), the
resulting \(\hbar=c^3\epsilon^2/(4G\sigma)\) does not depend on which de
Sitter \(H\) calibrates it. Extending the argument to a genuinely
dynamical FLRW horizon would require a Kodama--Hayward treatment showing
the same factor enters both the classical and emergent thermal
descriptions and cancels; that is not established here.

\textbf{The classical horizon temperature as a counting derivative,
under H-count.} Under A1, restrict the finite total system to a
conserved energy shell and let visible configuration \(i\) carry energy
\(e_i\). If the boundary interaction is negligible at the shell scale,
uniform counting gives \(P(i)=\Omega_H(E-e_i)/\sum_j\Omega_H(E-e_j)\),
so with \(s_H=\ln\Omega_H\) and \(\beta=s_H'(E)\),
\(P(i)\propto\exp[-\beta e_i+\tfrac12 s_H''(E)e_i^2+O(e_i^3)]\). Gibbs
is the leading finite-bath marginal of uniform counting, classically.
Under the hidden-sector condition \textbf{H-count} --- that the total
system is finite (A1) and that the energy-dependent entropy of the whole
hidden sector at the shell scale is the boundary count,
\(s_H=\sigma A/\epsilon^2\) up to an energy-independent constant, the
shell energy \(E\) being the energy of the horizon first law and the
deep hidden sector contributing no energy-dependent multiplicity and, by
the choice of energy zero, no energy; S1--S4 do not supply it, their
hidden capacity \(\sim S_{\rm dS}\) fixing no entropy--energy slope ---
feeding in \(dE=c^2\kappa\,dA/(8\pi G)\) gives
\(\partial s_H/\partial E=8\pi G\sigma/(c^2\kappa\epsilon^2)\) and hence
\(k_BT_{\rm micro}=(\partial s_H/\partial E)^{-1}=c^2\kappa\epsilon^2/(8\pi G\sigma)=k_BT_{\rm cl}\)
--- with no quantum mechanics, no KMS condition, no vacuum choice and no
infinite bath. On the exact de Sitter calibration \(E=c^4R/G\) and
\(s_H=4\pi\sigma G^2E^2/(c^8\epsilon^2)\), so \(s_H''=\beta/E\) and the
shell marginal is \(P(i)\propto\exp[-\beta e_i+\beta e_i^2/(2E)]\): not
exactly canonical, carrying the finite-reservoir correction
\(\delta_{\rm bath}\sim1/(4s_H)\sim10^{-123}\) at \(s_H\sim10^{122}\).
This is a second correction of the same spectacular smallness as the
lattice-dispersion term \((\epsilon H/c)^2\), arising for an unrelated
reason; the two should not be conflated.

\textbf{What this leaves.} The chain above is classical and gives a
Gibbs marginal in the classical subsystem energy. The horizon statement
needs \(\rho_{\rm obs}\propto e^{-\beta\hat H_{\rm eff}}\) in the
emergent time-translation generator, which {[}Main{]} distinguishes from
\(H_{\rm tot}\) and does not prove survives the trace-out. H-state
therefore splits into three narrower conditions: \textbf{H-shell} (the
finite microcanonical factorization --- \(P(i)\propto\Omega_H(E-e_i)\)
assumes weak boundary coupling, whereas C1 requires nonzero
cross-boundary coupling and strong coupling would instead give a
Hamiltonian of mean force), \textbf{H-energy} (the thermodynamic energy
intertwines with \(\hat H_{\rm eff}\)), and \textbf{H-spectrum}
(freely-falling UV regularity across the horizon). Only the third is a
genuinely spectral condition.

\textbf{H-balance: the minimal load-bearing state condition.} Full
H-spectrum is stronger than the calibration needs. For an accessible
horizon probe with positive frequency gap \(\omega\), require only \[
\frac{\Gamma_\uparrow(\omega)}{\Gamma_\downarrow(\omega)}=e^{-\tau_K\omega},\qquad \tau_K=\frac{2\pi c}{\kappa},
\] over the finite frequency window the observer can probe. KMS implies
this ratio. At detailed balance, \(p_mW_{m\to n}=p_nW_{n\to m}\), so the
excitation-to-reverse ratio is
\(W_{m\to n}/W_{n\to m}=p_n/p_m=e^{-\tau(\omega_n-\omega_m)}\). On a
connected transition graph this fixes the stationary populations to the
normalized Gibbs weights. The Gibbs density matrix also requires
diagonality in the energy basis, or a condition eliminating the
coherences. An edgewise log-ratio error \(\eta\) accumulates to at most
\(L\eta\) along a path of \(L\) accessible transitions. H-slope retains
the infrared part of this condition for the calibration.

\textbf{Precision on microreversibility.} With time-reversal involution
\(\Theta\), \(\Theta\varphi\Theta=\varphi^{-1}\) and
\(\Theta M_i=M_{\theta i}\), the exact trajectory-reversal map
\(x\mapsto\Theta\varphi(x)\) gives \(N_{ij}=N_{\theta j,\theta i}\), the
time-reversed detailed-balance relation. It reduces to
\(N_{ij}=N_{ji}\), and hence \(N_iP_{ij}=N_jP_{ji}\), only when the
coarse states are time-reversal even (\(\theta i=i\)). ``Spatial
partition'' must therefore not be read as merely ``\(\Theta\) maps
visible variables to visible variables'': if visible configurations
carry T-odd data such as momenta, only the \((\theta j,\theta i)\) form
holds, and the symmetric form used downstream is unavailable.

\textbf{Classical microreversibility already gives the classical half.}
With \(M_i\) the shell microstates compatible with visible macrostate
\(i\), \(N_i=|M_i|\) and \(N_{ij}=|\{x\in M_i:\varphi(x)\in M_j\}|\),
time-reversal invariance and time-reversal-even coarse states give a
bijection between \(i\to j\) and \(j\to i\) trajectories, so
\(N_{ij}=N_{ji}\) and hence \(N_iP_{ij}=N_jP_{ji}\) exactly. With
\(N_i=\Omega_H(E-e_i)\) this yields
\(\ln(P_{ij}/P_{ji})=-\beta_E(e_j-e_i)+\tfrac12 s_H''(E)(e_j^2-e_i^2)+\cdots\),
with \(\beta_E=s_H'(E)\) taking the boundary-count value above under
H-count --- the thermal rate ratio from finite deterministic reversible
counting, with no quantum mechanics. Equating the two descriptions of
the same observable rate ratio gives
\(\beta_E\Delta E=\tau_K\Delta\omega\) up to the two controlled
corrections, hence \(dE/d\omega=\tau_K/\beta_E\) in the infrared, and on
a connected accessible sector \(E_n=\hbar\omega_n+E_0\). The relation
naive H-energy tried to assume emerges instead from thermal rate
consistency. The remaining qualification is an identification of
observables --- that the classical shell transitions and the
frequency-resolved detector transitions are the same operational
transitions --- which, if it requires arbitrary detector couplings,
touches the coherent operational completion, which §3.3 characterizes,
within complex matrix kinematics taken as given, by completion
conditions whose physical sourcing is open.

\textbf{The infrared limit does not remove state dependence.} Dispersive
corrections vanish as \(k\to0\), but de Sitter admits non-Bunch-Davies
\(\alpha\)-vacua that are de Sitter invariant while failing detector
detailed balance. With \(r=e^\alpha\), the excitation-to-de-excitation
ratio
\(R_\alpha(\omega)=e^{-2\pi\omega/H}\big((1+re^{\pi\omega/H})/(1+re^{-\pi\omega/H})\big)^2\)
has zero-frequency slope
\(-d\ln R_\alpha/d\omega|_0=(2\pi/H)(1-r)/(1+r)\), which returns
\(2\pi/H\) only as \(r\to0\). So H-balance cannot be discharged by
taking \(\omega\to0\).

\textbf{H-slope: the minimal condition the calibration needs, and its
status.} Only the infrared slope is load-bearing. Writing the classical
shell rate ratio \(R_E\) and the observer frequency response
\(R_\omega\) for the same operational transition with
\(\Delta E=E(\omega)\), \(E(0)=0\), and
\(-\partial_{\Delta E}\ln R_E|_0=\beta_E\), the single requirement is
\(-\partial_\omega\ln R_\omega|_0=\tau_K\) (\textbf{H-slope}), whence,
with \(\beta_E\) at its H-count value,
\(E'(0)=\tau_K/\beta_E=c^3\epsilon^2/(4G\sigma)=\hbar\). Expanding
\(\ln R_E=-\beta_EE+BE^2+\cdots\),
\(\ln R_\omega=-\tau_K\omega+C\omega^2+\cdots\) and
\(E(\omega)=a\omega+d\omega^2+\cdots\) gives \(\beta_Ea=\tau_K\) at
first order exactly: the finite-bath and lattice-dispersion corrections
alter \(d\) and higher coefficients, not the linear IR action scale.
Full H-balance is still required for the global statement
\(E_n=\hbar\omega_n+E_0\) across a connected spectrum.

\textbf{H-balance is a selection condition, not a representation
theorem.} {[}Main{]}'s intervention-dilation theorem realizes any finite
action-labelled conditional family with one finite reversible dynamics,
a fixed prior and permutation interventions. Take a two-gap probe with
\(\omega_2=2\omega_1\) and \(R(\omega_1)=R(\omega_2)=\tfrac12\):
thermality at gap one forces \(\tau=\ln2/\omega_1\) and hence
\(R(\omega_2)=\tfrac14\), so the family is nonthermal for every common
inverse temperature --- yet it is OI-realizable. So finite realization,
reversibility and finite intervention statistics do not give H-balance,
and by the \(\alpha\)-vacuum slope above the failure already occurs at
H-slope. Note also that a bare visible-space transition law read off a
closed unitary would be bistochastic --- for two levels the forward and
reverse moduli coincide, giving ratio \(1\), the infinite-temperature
value --- which is one more reason that form is unavailable in general
(Main §3.1); so H-slope requires the ancilla-carrying \(Q_{\rm fb}\)
form together with a detector gap, an interaction and an identification
of excitation and de-excitation channels. Those are exactly the
structures the composite and coherent instrument lift leaves open.
\textbf{H-slope is therefore not an independent bridge to gravity; it is
a horizon-specific selection test on that open lift.} Complete passivity
is the natural candidate selector --- it forces every transition to
carry one common temperature and singles out Gibbs/KMS states --- but it
quantifies over composites, so invoking it relocates H-slope into the
same open layer rather than discharging it.

\textbf{The gap equation in its cleanest form:
\(\hbar=\tau_K/\beta_E\).} Write the emergent evolution as
\(U_t=e^{-it\Omega_{\rm eff}}\) with
\([\Omega_{\rm eff}]=\text{time}^{-1}\). Before the action scale is
fixed the physical energy generator can only be
\(H_{\rm phys}=a_*\Omega_{\rm eff}\) for some universal \(a_*\) with
units of action; determining \(a_*\) is what this section does, so it
must not be set to \(\hbar\) in advance. Classical horizon counting
supplies an inverse energy,
\(\beta_E=\partial s_H/\partial E=8\pi G\sigma/(c^2\kappa\epsilon^2)\)
under H-count. The emergent horizon state, conditional on H-spectrum,
supplies an inverse time, \(\tau_K=2\pi c/\kappa\). For a finite system
a KMS state for evolution generated by \(\Omega\) is Gibbs with exponent
\(-\tau_K\Omega\), so
\(e^{-\tau_K\Omega_{\rm eff}}=e^{-\beta_Ea_*\Omega_{\rm eff}}\) up to
normalization, giving \(\tau_K=\beta_Ea_*\) and \[
 a_*=\frac{\tau_K}{\beta_E}=\frac{2\pi c/\kappa}{8\pi G\sigma/(c^2\kappa\epsilon^2)}=\frac{c^3\epsilon^2}{4G\sigma}=\hbar .
\] The action scale is exactly the conversion between an inverse-energy
temperature obtained classically under H-count and an inverse-time KMS
period obtained independently of it. No quantum mechanics enters the
classical side and no \(\hbar\) is inserted into \(\tau_K\) beforehand.

\textbf{Why the classical shell state must not be identified with
\(e^{-\beta\hat H_{\rm eff}}\).} Suppose the configuration-diagonal
shell state \(\rho_{\rm shell}=\sum_ip_i|i\rangle\langle i|\) were
required to be Gibbs for the emergent generator. Since
\(\rho_{\rm shell}>0\) the matrix logarithm is unique, so
\(\Omega_{\rm eff}=-\tau^{-1}\log\rho_{\rm shell}+\text{const}\,I\) is
diagonal in the configuration basis, whence
\(|\langle j|e^{-it\Omega_{\rm eff}}|i\rangle|^2=\delta_{ij}\) --- no
configuration dynamics at all. That contradicts the generic
phase-locking sector of {[}Main §3.2{]}, where the energy eigenbasis
overlaps every configuration state and \(V=I\) is recorded as a
nongeneric degeneracy. The diagonal embedding
\(p(x)\mapsto\sum_xp(x)|x\rangle\langle x|\) is an embedding of
classical probability, not a quantization map \(E(x)\mapsto\hat H\), and
using it as one is a no-go rather than an unproved bridge. A second,
independent obstruction: \(U_\varphi=e^{-i\Omega}\) has a non-unique
Hermitian logarithm, \(\omega_a=\theta_a+2\pi n_a\) with \(n_a\)
selectable per eigenmode, so the bare permutation does not determine
continuous frequency gaps, while a conserved classical energy is merely
constant on cycles.

\textbf{Consequences for the named conditions.} H-shell is therefore not
load-bearing here: the calibration needs only the scalar
\(\beta_E=\partial s_H/\partial E\), which follows from the horizon
density of states and the first law under H-count, not a visible/hidden
product Gibbs marginal. That matters independently, since the shell
conditional \(\mu_H(h\mid i)\propto\mathbf 1[E_H(h)=E-e_i]\) depends on
\(i\), whereas {[}Main{]}'s channel uses a fixed hidden prior
\(\mu_H(h)\); substituting one for the other would replace a product
preparation by a correlated ensemble and a different projection theorem.
H-shell is retained as an auxiliary result on classical probe
equilibration and finite-bath corrections. H-energy in its naive form
--- configuration energies as eigenvalues of \(\hat H_{\rm eff}\) --- is
false by the no-go above; its meaningful form is modular, that the
static thermodynamic generator (the Hamiltonian of mean force at
non-negligible coupling) agrees with the observer's dynamical generator
in the infrared up to one universal action scale and an additive
constant. That relation follows automatically once H-spectrum supplies
the KMS property, so H-energy is redundant rather than open.
\textbf{H-slope is the one load-bearing state condition: H-spectrum
\(\Rightarrow\) H-balance \(\Rightarrow\) H-slope, and only the last is
needed for the calibration.}

\textbf{Status of \(\hbar\) and the \(1/4\).} The exact de Sitter
calibration requires H-count, including A1 and the conserved energy
shell (§2.3), H-slope on identified transitions, H-frame, the relevant
H-Hawking horizon and frame conditions (§8.5), local boost structure and
the symmetric detailed-balance rate ratio. The freely falling
adiabatic/de Sitter vacuum class (H-state with H-spectrum) suffices for
the thermal slope in this setting. The calibration uses H-slope itself.
The finite representation admits different fixed hidden priors for the
same update and partition. Selection of the physical slope remains an
additional condition.

With the one-mode-per-area count and a positive counting scale
\(\epsilon\), the finite horizon area \(A=4\pi c^2/H^2\) gives finitely
many boundary modes. Lemma 1 uses finite distinguishability of visible
configurations as a separate premise. For a finite invariant
representative with a non-permutation rooted map, the rooted description
is P-indivisible across its recurrence cycle and, by either route of
{[}Main, §3{]}, admits a fixed-basis unitary quantum representation,
with \(\hbar\) fixed by the calibration of §3.2, conditional on its
named conditions.

\begin{center}\rule{0.5\linewidth}{0.5pt}\end{center}

\hypertarget{the-emergent-action-scale-hbar}{%
\subsection{\texorpdfstring{3. The Emergent Action Scale
\(\hbar\)}{3. The Emergent Action Scale \textbackslash hbar}}\label{the-emergent-action-scale-hbar}}

\hypertarget{the-classical-horizon-temperature}{%
\subsubsection{3.1 The classical horizon
temperature}\label{the-classical-horizon-temperature}}

Jacobson {[}1{]} showed \(dE = T\,dS\) applied to local causal horizons
yields Einstein's equations, with \(dE = (c^2 \kappa / 8\pi G)\,dA\),
where \(\kappa\) is the surface gravity of the horizon. The entropy
density is \(\eta = \sigma/\epsilon^2\), where \(\epsilon\), the
discreteness scale of the boundary count, is defined by
\(\epsilon^2 = A/N_{\rm modes}\): the horizon area per boundary mode
that couples across the partition, each mode carrying the dimensionless
entropy \(\sigma\), in nats. Its relation to the lattice spacing
\(\epsilon_{\rm lat}\) is fixed by the link count below. The count does
not fix \(\sigma\): a mode with \(q\) internal states carries at most
\(\ln q\), and the alphabet-size hypothesis of {[}SM §2.7{]} (two
systems with different \(q\) and the same coupling structure produce
identical transition probabilities) does not fix it either. The
Bekenstein-Hawking coefficient does not depend on \(\sigma\), which
cancels from \(S = \eta A\) exactly as \(\kappa_m\) does below;
\(\sigma\) enters only the relation between \(\epsilon\) and \(l_p\)
(§4). Thus \(dS = \sigma\,dA/\epsilon^2\). Equating:

\[T_{\text{cl}} = \frac{c^2 \epsilon^2 \kappa}{8\pi G k_B \sigma}\]

For de Sitter (\(\kappa = cH\)):
\(k_B T_{\text{cl}} = c^3 \epsilon^2 H / (8\pi G \sigma)\). No \(\hbar\)
appears. Every value quoted here and in the companion papers assumes
\(\sigma = 1\), one nat per boundary mode. That is an assumption about
the entropy of a mode, not a choice of units, since entropy is already
measured in nats; under the bound \(\sigma \leq \ln q\) it is
unavailable for a single mode with \(q = 2\).

\emph{Reconciliation with the internal component count.} The counting
\(\eta = \sigma/\epsilon^2\) is one boundary mode per area
\(\epsilon^2\), whereas the substratum carries \(K = 2d = 6\) internal
components per site ({[}SM §4.5{]}). Under H-link with the single-link
reading --- component \(a\) assigned to link direction
\(\hat e_a \in \{\pm\hat e_1, \ldots, \pm\hat e_d\}\) and coupling only
across that one link (the factorization \(\delta = 1\), which {[}SM
§4.5, Lemma 6a{]} records as not a consequence of {[}SM §4.4{]}) --- the
components whose coupling crosses the horizon are one per
nearest-neighbour link crossing it. On the cubic lattice of spacing
\(\epsilon_{\rm lat}\) a planar patch with unit normal \(\mathbf n\) is
crossed by
\((\lvert n_1\rvert+\lvert n_2\rvert+\lvert n_3\rvert)/\epsilon_{\rm lat}^2\)
links per unit area: \(1/\epsilon_{\rm lat}^2\) for a patch orthogonal
to a lattice axis, \(\sqrt3/\epsilon_{\rm lat}^2\) for one orthogonal to
a body diagonal. For the de Sitter horizon, a sphere of radius
\(R \gg \epsilon_{\rm lat}\) in the lattice frame with the sites inside
it on the visible side, the area average of \(\lvert n_j\rvert\) is
\(\tfrac12\) for every orientation of the lattice, and, for a sphere
centred on a site, the number of crossing links is exactly
\(6N_2(R^2/\epsilon_{\rm lat}^2)\), with \(N_2(r^2)\) the number of
integer points in a disc of radius \(r\), that is
\(\tfrac32 A/\epsilon_{\rm lat}^2 + O(R/\epsilon_{\rm lat})\). Writing
\(\eta = \kappa_m\sigma/\epsilon_{\rm lat}^2\), the spherical horizon
has \(\kappa_m = \tfrac32\) and
\(\epsilon^2 = \epsilon_{\rm lat}^2/\kappa_m = \tfrac23\epsilon_{\rm lat}^2\).
The Bekenstein-Hawking coefficient is independent of \(\kappa_m\): the
gap equation gives \(\hbar = c^3/(4G\eta)\) and
\(S = \eta A = A/(4l_p^2)\) with \(\kappa_m\) and \(\sigma\) cancelling,
and the mode count is \(N_{\rm modes} = A/\epsilon^2 = S/\sigma\). The
component count enters only the lattice spacing: with
\(\epsilon = 2\sqrt{\sigma}\,l_p\) (§4),
\(\epsilon_{\rm lat} = 2\sqrt{\kappa_m\sigma}\,l_p = \sqrt{6\sigma}\,l_p\),
which is \(\sqrt6\,l_p \approx 2.45\,l_p\) at \(\sigma = 1\). The
density depends on orientation: a local patch with normal \(\mathbf n\)
has
\(\eta(\mathbf n) = \sigma(\lvert n_1\rvert+\lvert n_2\rvert+\lvert n_3\rvert)/\epsilon_{\rm lat}^2\),
equal to the spherical value \(\sigma/\epsilon^2\) only where
\(\lvert n_1\rvert+\lvert n_2\rvert+\lvert n_3\rvert = \tfrac32\);
whether the emergent description carries an orientation-independent
density is open. (The entropy per mode, \(\sigma\), is a separate
question: one mode per area \(\epsilon^2\) does not establish one nat
per mode.) The matching scale \(\mu = M_{\text{Pl}}\) of {[}SM §6.2{]}
is a convention into which this factor does not enter; matching at
\(1/\epsilon_{\rm lat}\) instead shifts the fitted threshold parameters
there ({[}SM §6.2{]}).

\hypertarget{the-emergent-action-scale}{%
\subsubsection{3.2 The emergent action
scale}\label{the-emergent-action-scale}}

\textbf{Step 1: Uniqueness.} \(\hbar\) is derived, not free, conditional
on the calibration \(a_*=\tau_K/\beta_E\) (§2.3; Step 4) and its named
conditions.

\textbf{Step 2: The transition data are structural, not volumetric.}

\textbf{Lemma (Boundary-only dependence).} \emph{Let the bijective
update \(\varphi\), with update step \(\Delta t\), act on a finite
product of site configuration spaces with coupling graph \(G_\varphi\)
({[}Main §3.3{]}, coupling-graph causal cone), and let the hidden-sector
prior be uniform. (a) For \(n \geq 1\) and \(t = n\,\Delta t\),
decompose the hidden sector as
\(\mathcal{C}_H = \mathcal{C}_{B_n} \times \mathcal{C}_{D_n}\), where
\(B_n = B_{G_\varphi}(V, n) \setminus V\), the visible sector's causal
cone over the elapsed time, is the set of hidden sites within graph
distance \(n\) of \(V\), and \(D_n\) is the remainder. Then, for every
\(d^* \in \mathcal{C}_{D_n}\), exactly,}

\[T_{ij}(t) = T^{(B_n)}_{ij}(t) \equiv \frac{1}{|\mathcal{C}_{B_n}|} \sum_{b \in \mathcal{C}_{B_n}} \delta_{x_j}[\pi_V(\varphi^n(x_i, b, d^*))].\]

\emph{(b) Let \(\mathcal{C}_B\) denote the boundary modes, within the
interaction range of \(H_{\text{int}}\), and \(\mathcal{C}_D\) the
remainder of the hidden sector. Under H-BD, for \(t = n\,\Delta t\),}

\[\sum_j \big|T_{ij}(t) - T^{(B)}_{ij}(t)\big| \leq \frac{2t}{\tau_{BD}}, \qquad T^{(B)}_{ij}(t) = \frac{1}{|\mathcal{C}_B|} \sum_{b \in \mathcal{C}_B} \delta_{x_j}[\pi_V((\varphi^{VB})^n(x_i, b))],\]

\emph{where \(T^{(B)}_{ij}\) depends only on the \(V\)-\(B\) dynamics.}

\textbf{H-BD} (deep-sector influence bound; posited). \emph{There is a
map \(\varphi^{VB}\) of \(\mathcal{C}_V \times \mathcal{C}_B\) to itself
such that, for every \((v, b)\), the \(V \times B\) component of
\(\varphi(v, b, d)\) differs from \(\varphi^{VB}(v, b)\) for at most a
fraction \(\Delta t/\tau_{BD}\) of \(d \in \mathcal{C}_D\).} For
\(\Delta t/\tau_{BD} < 1/2\) two such maps agree with the update on more
than half of \(\mathcal{C}_D\) in every fiber, so \(\varphi^{VB}\) is
the fiberwise majority value and is fixed by \(\varphi\).

\emph{Proof.} (a) By the coupling-graph causal cone ({[}Main §3.3{]}),
after \(n\) updates each visible component depends only on initial
components in \(B_{G_\varphi}(V, n)\), so
\(\pi_V(\varphi^n(x_i, b, d))\) does not depend on
\(d \in \mathcal{C}_{D_n}\). Under the uniform prior the sum over
\(\mathcal{C}_{D_n}\) in \(T_{ij}(t)\) contributes
\(|\mathcal{C}_{D_n}|\) identical terms, for every \(d^*\).

\begin{enumerate}
\def\labelenumi{(\alph{enumi})}
\setcounter{enumi}{1}
\tightlist
\item
  Fix \(x_i\) and \(b\), and follow the reference path
  \((v_s, b_s) = (\varphi^{VB})^s(x_i, b)\). Let
  \(S_s \subseteq \mathcal{C}_D\) be the set of initial \(d\) whose
  trajectory under \(\varphi\) has \(V \times B\) component
  \((v_{s'}, b_{s'})\) at every step \(s' \leq s\). Elements of \(S_s\)
  share their \(V \times B\) components through step \(s\), so
  injectivity of \(\varphi\) makes their \(D\) components at step \(s\)
  distinct, and by H-BD at most a fraction \(\Delta t/\tau_{BD}\) of
  \(\mathcal{C}_D\) leaves \(S_s\) at step \(s+1\). Hence
  \(|\mathcal{C}_D \setminus S_n| \leq (t/\tau_{BD})\,|\mathcal{C}_D|\),
  and on \(S_n\),
  \(\pi_V(\varphi^n(x_i, b, d)) = \pi_V((\varphi^{VB})^n(x_i, b))\). The
  average over \(d\) therefore lies within \(\ell_1\) distance
  \(2t/\tau_{BD}\) of \(\delta_{x_j}[\pi_V((\varphi^{VB})^n(x_i, b))]\),
  and the average over \(b\) keeps the bound. \(\square\)
\end{enumerate}

Spatial locality of the classical Hamiltonian gives the coupling chain
\(V \leftrightarrow B \leftrightarrow D\), with no direct \(V\)-\(D\)
coupling, but a deep-sector configuration that stays fixed still enters
the boundary update, so a bound on the motion of \(D\) does not bound
its effect on \(V\). For
\(\varphi(v, b, d) = (v \oplus b, b \oplus d, d)\) on bits, \(d\) never
changes, yet every entry of \(T(2\Delta t)\) is \(\tfrac12\), while the
\(V \times B\) update with \(d\) clamped at \(0\) or \(1\) gives the
identity or the bit flip. S1--S4 and {[}Main{]} fix no value of
\(\tau_{BD}\). For a linear update ({[}Main §4.5{]}, ledger item (vii))
on a finite abelian group with \(\mathcal{C}_D\) a direct summand, the
\(V \times B\) component at fixed \((v, b)\) is \(A(v, b) + Ld\) with
\(L: \mathcal{C}_D \to \mathcal{C}_V \oplus \mathcal{C}_B\) a
homomorphism, so the set of \(d\) on which it takes a given value is
empty or a coset of \(\ker L\). The fraction in H-BD is then \(0\) if
\(L = 0\), and at least \(1 - 1/|\operatorname{im} L| \geq 1/2\)
otherwise. Under (vii), H-BD with \(\tau_{BD} > 2\Delta t\) therefore
holds if \(L = 0\), that is, if the deep sector does not enter the
\(V \times B\) update, and fails if \(L \neq 0\). The update displayed
in {[}Main §3.2{]} has \(L \neq 0\) at every cut where a boundary site
has a deep neighbour, since each such site reads that neighbour's \(u\)
with coefficient \(1\).

Under H-BD the correction is at most \(2t/\tau_{BD}\), which is of order
\(10^{-32}\) for laboratory processes when \(\tau_{BD} \gtrsim 1/H\), a
value S1--S4 do not fix. Since \(T_{ij}(t)\) determines the emergent
quantum description up to its action scale (Prerequisites; Barandes
correspondence {[}35, 36{]}), and the transition data on
\([0, n\,\Delta t]\) depend only on the update of \(V\) and of the
hidden sites within \(n\) update steps of \(V\) and on the prior there,
the emergent generator \(\Omega_{\rm eff}\), as far as those data fix
it, depends only on that local update and its causal cone; under H-BD
those data agree with the one-layer law \(T^{(B)}_{ij}(t)\) to within
\(2t/\tau_{BD}\). That the local update is itself fixed by the boundary
geometry (\(A\), \(\epsilon\), \(\kappa\)) and the constants \(c\),
\(G\) of the classical Hamiltonian is a further condition: two local
bijective updates on the same geometry can give different transition
probabilities. The action scale is not fixed by these data: every
\(a_*\) in \(H_{\rm phys}=a_*\Omega_{\rm eff}\) (§2.3) gives the same
\(T_{ij}(t)\), so the lemma does not bear on \(\hbar\), which the
calibration \(a_*=\tau_K/\beta_E\) of §2.3 and Step 4 fixes. Steps 2 and
3 are consistency checks, Step 2 under the lemma's hypotheses, and no
result of §8.5(b) uses them. If the traced-out hidden state moves at a
rate of order \(1/\tau_B\), the slow-bath regime of {[}Main §3.2{]}, the
departure of the reduced visible dynamics from its frozen-bath form
enters at first order as a Hermitian drift of the emergent effective
Hamiltonian \(\hat{H}_{\text{eff}}\) of order
\((\tau_S/\tau_B)\,\|V_{\text{int}}\|\) ({[}Main §3.2{]}), with
\(\tau_S/\tau_B \sim 10^{-32}\) when \(\tau_B \sim 1/H\); this mechanism
is separate from the deep-sector influence that H-BD bounds. Whether
that drift is common-mode --- a uniform rescaling of
\(\hat{H}_{\text{eff}}\) through the constraint coupling of (C1),
cancelling in dimensionless ratios up to a composition-differential
remainder of order \((L_{\text{lab}}/R_H)^2\), \(R_H\) the horizon
radius --- is open.

\textbf{Corollary (deep-sector cardinality is gauge).} \emph{No quantity
determined by the visible transition data on \([0, n\,\Delta t]\)
depends on the cardinality \(|\mathcal{C}_{D_n}|\) of the hidden sector
beyond the causal cone \(B_n\): neither those transition probabilities,
nor the set of emergent Hamiltonians compatible with them, nor any
quantity that takes the same value on every member of that set. Under
H-BD the transition probabilities \(T_{ij}(t)\) do not depend on the
cardinality \(|\mathcal{C}_D|\) of the deep hidden sector to within
\(\mathcal{O}(t/\tau_{BD})\); the same accuracy for quantities fitted to
them, such as emergent Hamiltonians or their spectra, requires a
stability bound not derived here. Systems with the same update on the
sites of \(V \cup B_{n-1}\) produce the same visible transition data on
\([0, n\,\Delta t]\), and under H-BD systems with the same
\(\mathcal{C}_V \times \mathcal{C}_B\) dynamics produce the same data to
within \(\mathcal{O}(t/\tau_{BD})\), for every finite deep sector. A
second system has the same update on the sites of \(V \cup B_{n-1}\)
when those sites and the sites on which their updates depend are
identified with sites of the second system by an injection \(e\) that is
the identity on \(V\) and preserves alphabets, such that for each site
\(i\) of \(V \cup B_{n-1}\) the update at \(e(i)\) depends exactly on
the images of the sites on which the update at \(i\) depends, through
the same local rule. The statements that do not use H-BD hold also for a
deep sector of any cardinality under hypotheses (D1)--(D5) of Appendix
B: \(V\) is a finite set of sites; the update is defined on every
configuration, and each component's update is a function of finitely
many components; the coupling graph is locally finite; the alphabets on
\(V \cup B_n\) are finite; and the hidden prior is the product of the
uniform distribution on \(\mathcal{C}_{B_n}\) and a normalized prior on
\(\mathcal{C}_{D_n}\).}

\emph{Proof.} \(T^{(B)}_{ij}(t)\) is defined by a sum over
\(\mathcal{C}_B\) alone, and under H-BD \(\mathcal{C}_D\) enters
\(T_{ij}(t)\) only through a correction of at most \(2t/\tau_{BD}\)
(part (b) of the lemma). The statements that do not use H-BD are proved
in Appendix B under (D1)--(D5), which a finite system with the uniform
hidden prior satisfies, the uniform distribution on
\(\mathcal{C}_{B_n} \times \mathcal{C}_{D_n}\) being the product of the
uniform distributions: for \(m \leq n\),
\(T_{ij}(m\,\Delta t) = T^{(B_n)}_{ij}(m\,\Delta t)\) is a count over
the configurations of the dependency patch \(W_n \subseteq B_n\) of the
visible sector, which involves neither \(\mathcal{C}_{D_n}\) nor the
prior on it and is the same for two systems with the same update on the
sites of \(V \cup B_{n-1}\). The set of emergent Hamiltonians compatible
with \(\{T_{ij}(m\,\Delta t) : m \leq n\}\) is fixed by those data
alone, so it inherits their exact independence of \(\mathcal{C}_{D_n}\),
and so does every quantity that takes the same value on all its members.
\(\square\)

A quantity that differs between members of the compatible set is not
fixed by the visible data, and the corollary does not extend to it. An
individual eigenvalue is such a quantity: \(H \mapsto H + E_0\)
multiplies \(U(t)\) by the phase \(e^{-iE_0 t}\), leaving every
\(T_{ij}(t)\) unchanged while shifting each eigenvalue by \(E_0\). So is
a raw amplitude \(U_{ij}(t)\), whose phase changes under
\(H \mapsto DHD^\dagger\) and which is complex-conjugated under
\(H \mapsto -\bar{H}\), while both maps leave \(|U_{ij}(t)|^2\)
unchanged.

The cardinality of the deep hidden sector is gauge in the scope of this
corollary: two systems with the same update on \(V \cup B_{n-1}\) and
different \(|\mathcal{C}_{D_n}|\) share \(T_{ij}(t)\) exactly for
\(t \leq n\,\Delta t\); under H-BD, two systems with the same
\(\mathcal{C}_V \times \mathcal{C}_B\) dynamics and different
\(|\mathcal{C}_D|\) share the leading-order transition law
\(T^{(B)}_{ij}(t)\), so no statement expressible through
\(\{T^{(B)}_{ij}(t) : t \ll \tau_{BD}\}\) distinguishes them, while a
statement about \(T_{ij}(t)\) itself is shared only if it is insensitive
to \(\mathcal{O}(t/\tau_{BD})\) corrections ({[}Substratum §4{]},
Corollary (effective finiteness and gauge-class transfer)). Exact
invariance of every \(T_{ij}(t)\), the gauge enlargement of
{[}Substratum §4{]}, Theorem 24(iii), holds for a deep sector that no
dynamics couples to the boundary layer. The recurrence proof ({[}Main,
§2.3{]}) requires \(S\) finite, which Lemma 1 does not supply;
recurrence-scale content lies outside that scope: P-indivisibility at
the recurrence scale is proved on a finite representative, while for an
infinite deep sector the readback lemma ({[}Main, §2.3{]}) gives the
weaker accessible-window property, conditional memory, and the
near-return corollary ({[}Main, §2.3{]}) gives P-indivisibility for a
measurable update from a non-permutation rooted map under its
near-return condition (NR), assumed in place of finiteness as a
hypothesis and not adopted. Within the visible transition data on
\([0, n\,\Delta t]\), a finite deep sector beyond \(B_n\) and an
infinite one with the prior and dynamics above are indistinguishable.

\textbf{Step 3: Dimensional determination.} The calibration of §2.3 and
Step 4 expresses \(\hbar\) through boundary quantities, \(\beta_E\)
under H-count; the dimensional argument below checks the form of that
result. The boundary carries both \emph{local} geometric data
(\(\epsilon\), \(\kappa\), and the constants \(c\), \(G\) of the
classical Hamiltonian) and a \emph{global} quantity: the total area
\(A\), which forms the dimensionless ratio
\(A/\epsilon^2 = S_{\text{dS}}/\sigma\). The area and the surface
gravity can vary between horizons. The calibrated relation in Step 4 is
independent of them at fixed counting density and common \(c\) and
\(G\). Restricting the dimensional analysis to the surviving dimensional
quantities gives the unique combination of \(c\), \(G\), \(\epsilon\)
with dimensions of action:

\[\hbar = \beta \frac{c^3 \epsilon^2}{G}\]

\textbf{Step 4: Thermal self-consistency.} The classical substratum
assigns the partition boundary a temperature
\(k_B T_{\text{cl}} = c^2 \epsilon^2 \kappa / (8\pi G \sigma)\) (§3.1),
containing no \(\hbar\). The emergent QFT of Part I lives on this
classical background. On the exact de Sitter calibration of §2 that
background has a bifurcate Killing horizon with surface gravity
\(\kappa\), and regularity of the Wick-rotated metric requires Euclidean
period \(\beta = 2\pi c/\kappa\), giving a KMS state at
\(T_Q = \hbar\kappa/(2\pi c k_B)\) with \(\hbar\) unknown. Two scope
limits apply and are load-bearing. First, periodicity is not inherited
by every QFT on this background: a preferred-frame dispersive lattice
theory satisfies it only under the H-Hawking conditions of §8.5, and the
dispersive-Hawking literature supplies counterexamples where they fail,
so periodicity does not hold for every QFT on this background. Second,
the KMS temperature parameter is fixed kinematically by \(\kappa\);
whether the emergent state belongs to that KMS class is state-dependent,
and which class it occupies is H-spectrum, which §2 records as
underivable from S1--S4. The equality below is therefore a conditional
infrared calibration, not a theorem within the derived QFT. The two
temperatures are computed independently --- \(T_{\text{cl}}\) from the
classical substratum alone (no QM), \(T_Q\) from the emergent QFT alone
(no classical substratum details) --- but they refer to the \textbf{same
physical degrees of freedom}: the boundary modes \(V \times B\) across
which \(H_{\text{int}}\) couples the visible and hidden sectors. Under
H-BD, Step 2's boundary-only dependence shows that the visible
transition data from which the emergent QFT is reconstructed are fixed,
up to \(\mathcal{O}(t/\tau_{BD})\), by these same \(V \times B\)
dynamics; the matching does not use it, the identification of the two
descriptions' transitions being part of H-slope (§2.3). H-slope fixes
the frequency slope for the same operational transitions as those
counted under H-count. The slopes satisfy \(\beta_E\hbar=\tau_K\). In
temperature units, this is \(T_{\text{cl}}=T_Q\) for the calibration:

\[\frac{c^2 \epsilon^2 \kappa}{8\pi G\sigma} = \frac{\hbar \kappa}{2\pi c}\]

The common surface gravity \(\kappa\) cancels. Solving gives:

\[\boxed{\hbar = \frac{c^3 \epsilon^2}{4G\sigma}}\]

which is \(\hbar = c^3 \epsilon^2/(4G)\) at \(\sigma = 1\) (§3.1).

The gap equation determines the action scale from the boundary count and
the matched transition slopes. H-count gives the scalar inverse-energy
slope \(\beta_E\). H-slope supplies \(\tau_K\) for the identified
transitions. The finite-horizon fixed-basis representation of {[}Main{]}
by itself fixes neither the physical calibration nor a coherent
operational completion (§3.3).

\textbf{One boundary mode per area \(\epsilon^2\) fixes \(\epsilon\).}
The boundary count fixes \(\epsilon\) through
\(\epsilon^2 = A/N_{\rm modes}\) (§3.1); the substratum enters through
the relation of \(\epsilon\) to the lattice spacing,
\(\epsilon^2 = \tfrac23\epsilon_{\rm lat}^2\) on the spherical horizon
under the single-link reading (§3.1). The number of internal states per
mode bounds its entropy, \(\sigma \leq \ln q\), and does not change the
mode count; \(\sigma\) is the entropy per mode of §3.1, which neither
the count nor the alphabet-size hypothesis of {[}SM §2.7{]} fixes. The
mode density \(1/\epsilon^2\) in \(\eta = \sigma/\epsilon^2\) is
therefore not an independent physical ansatz but a consequence of the
definition of \(\epsilon\); \(\sigma\) is a separate input, and the
Bekenstein-Hawking coefficient is independent of it (§3.1).

The specific relationship \(\hbar = c^3\epsilon^2/(4G\sigma)\) enters
three later results: the Bekenstein-Hawking formula with the factor
\(1/4\) (§5), the compression ratio \(S_{\text{dS}}\) of the conditional
CC dissolution (§6.3), and the Type II running-vacuum reading of the
dark energy (§7.1). The three relations share the calibration; they are
consequences of it, not independent checks of it. The situation is
analogous to deriving the Schwarzschild metric with \(M\) as a free
parameter: the derivation has genuine content (the functional form) even
though one input is not determined from within.

Jacobson's original derivation {[}1{]} uses \(\hbar\)-containing forms;
this paper does not --- it uses the classical identity
\(dE = (c^2 \kappa / 8\pi G)\,dA\) and the classical entropy density
\(\eta = \sigma/\epsilon^2\). The logical ordering differs from
Jacobson's: his argument derives \(G\) from \(\eta\) and \(\hbar\); the
present argument takes \(G\) as a parameter of the classical Hamiltonian
(Lemma 3) and derives \(\hbar\) from \(G\) and \(\epsilon\). The two are
consistent on the spherical horizon --- substituting
\(\eta = \sigma/\epsilon^2\) and \(\hbar = c^3\epsilon^2/(4G\sigma)\)
into Jacobson's relation \(G_{\text{eff}} = c^3/(4\eta\hbar)\) recovers
\(G_{\text{eff}} = G\), with \(\sigma\) cancelling --- but the roles of
input and output are reversed. On a planar patch with unit normal
\(\mathbf n\) the single-link density of §3.1 gives
\(G_{\text{eff}}(\mathbf n) = 3G/\big(2(\lvert n_1\rvert+\lvert n_2\rvert+\lvert n_3\rvert)\big)\),
between \(\tfrac{\sqrt3}{2}G\) and \(\tfrac32G\), so the relation holds
for the area-averaged density of the spherical horizon (and on patches
with \(\lvert n_1\rvert+\lvert n_2\rvert+\lvert n_3\rvert = \tfrac32\)),
not patch by patch. The Gibbons-Hawking form {[}2{]}
\(k_B T_{\text{GH}} = \hbar \kappa / (2\pi c)\) enters as the KMS
temperature of Step 4; the calibration relates the \(\hbar\) in it to
\(\epsilon\) and \(\sigma\) through
\(\hbar = c^3\epsilon^2/(4G\sigma)\). The KMS condition used above is
derived for continuum QFT on bifurcate Killing horizons; for the
lattice-regularized theory emerging from Part I, the relevant result is
that thermal periodicity of the Euclidean Green's functions is robust
against UV modifications of the dispersion relation, with corrections at
\(\mathcal{O}((\epsilon\kappa/c^2)^2)\) {[}3, 4{]}. For the cosmological
horizon, \(\epsilon\kappa/c^2 = \epsilon H / c \sim 10^{-61}\) at
\(\sigma = 1\), so lattice corrections are \(\sim 10^{-122}\).

\textbf{An independent necessity statement.} Einstein--Hilbert dynamics
follows, given the single-cone condition stated below, from an external
theorem with a disjoint axiom base. For matter field equations whose
principal polynomial is that of standard-model fields on a Lorentzian
metric --- the class in which the emergent sector of Part I falls if its
single-cone structure holds, that structure being supported by the
rotational-protection results of {[}SM §3.1{]}, conditional there on
loop-level \(O_h\) invariance and an \(O\)-equivariant condensate under
H-link, and by the symmetric-point universality of the block cones under
H-blind ({[}SM §4.4{]}) --- the gravitational closure theorem of Düll,
Schuller, Stritzelberger and Wolz {[}Phys. Rev.~D 97, 084036 (2018){]}
determines the gravitational action to be Einstein--Hilbert with
cosmological constant, from canonical co-evolution and diffeomorphism
invariance alone, with no thermodynamic input. The closure route leaves
\(G\) and \(\Lambda\) as undetermined constants, which is where the
present relation \(\hbar = c^3\epsilon^2/(4G\sigma)\) and the analysis
of §6--7 take over. The closure theorem, moreover, is one member of a
family of independent necessity statements the emergent sector inherits
once its inputs stand: Lovelock's uniqueness theorem {[}J. Math. Phys.
12, 498 (1971){]} --- second-order, metric-only, generally covariant
dynamics in four dimensions is Einstein--Hilbert with cosmological
constant --- and the spin-2 self-coupling and soft-graviton results
{[}Weinberg, Phys. Rev.~138, B988 (1965); Deser, Gen.~Rel. Grav. 1, 9
(1970){]}, by which a massless spin-two mode with a complete matter
sector is forced into universal coupling and the Einstein--Hilbert form
at low energy. Each route would deliver Einstein--Hilbert dynamics from
hypotheses the framework has not established --- the single cone, the
massless spin-two bound state, the absence of extra light gravitational
modes --- and the perturbative members of the family carry their own
known subtleties in the literature; the routes bear on the standard
analog-gravity objection, that emergent spacetimes deliver kinematics
but not Einstein dynamics, only once those hypotheses are established.
The present derivation uses the Einstein equations through Jacobson's
relation (§3.1) and does not derive them; the covariant field equations
remain the Level-G3 target (§8.7). The finite-\(\epsilon\) deformation
of the principal polynomial, and its comparison against the closure of
weakly birefringent electrodynamics {[}arXiv:1708.03870{]}, is an open
consistency check.

\hypertarget{gauge-fixing-and-the-dimensional-obstruction}{%
\subsubsection{3.3 Gauge-fixing and the dimensional
obstruction}\label{gauge-fixing-and-the-dimensional-obstruction}}

Step 4 fixes \(\hbar\) in terms of \(\epsilon\) via thermal matching. A
prior question must be addressed: does the transition-probability data
\(\{T_{ij}(t)\}\) determine the emergent Hamiltonian uniquely, or does
residual gauge freedom leave \(\hbar\) underdetermined? The following
theorem settles what the argument needs: the transition data determine
the full dynamical frequency content, and the dimensional obstruction
then makes the determination of \(\hbar\) a question no amount of such
data can settle. The sharper question --- how far the data pin the
Hamiltonian itself --- is answered by the two-branch theorem below,
conditional on its classification hypothesis H-BG.

\textbf{Theorem (Bohr-frequency completeness).} \emph{Let
\(U(t) = e^{-iHt}\) with non-degenerate eigenvalues and non-vanishing
configuration-basis overlaps. Then the transition probabilities
determine the Bohr-frequency set \(\{E_a - E_b\}\): each return
probability expands as
\(|U_{ii}(t)|^2 = \sum_{a,b} |V_{ia}|^2 |V_{ib}|^2\, e^{-i(E_a - E_b)t}\),
in which every gap value carries a strictly positive coefficient; every
\(|U_{ij}(t)|^2\) is a combination of the same frequencies; and
consequently any \(U'\) of the same form with
\(|U'_{ij}(t)|^2 = |U_{ij}(t)|^2\) for all \(i, j, t\) has
\(\{E'_a - E'_b\} = \{E_a - E_b\}\).}

\emph{Proof.} Expanding
\(U_{ij}(t) = \sum_a V_{ia}\, e^{-iE_a t}\, V_{ja}^*\) gives
\(|U_{ij}(t)|^2 = \sum_{a,b} V_{ia} V_{ja}^* V_{ib}^* V_{jb}\, e^{-i(E_a - E_b)t}\);
at \(i = j\) the coefficients are \(|V_{ia}|^2 |V_{ib}|^2 > 0\), and
grouping terms by frequency value only adds positive quantities ---
degenerate gaps merge coefficients but cannot cancel them. Distinct
characters \(t \mapsto e^{-i\omega t}\) of \((\mathbb{R}, +)\) are
linearly independent (Dedekind), so equality of two probability families
forces equality of the per-frequency coefficients over the union of the
two gap sets, and positivity makes the supports --- the gap sets ---
coincide. \(\square\)

\textbf{The dimensional obstruction.} This is what §3's argument
consumes, and it is deliberately less than a reconstruction of \(H\).
The transition probabilities are dimensionless functions of \(t\) and
determine temporal frequencies; converting a frequency to an energy
requires the action scale \(\hbar\) itself, so no amount of transition
data fixes it. The residual ambiguities in passing from frequencies to a
Hamiltonian carry no scale either: the energy origin
\(H \mapsto H + E_0\) contributes only the global phase \(e^{-iE_0 t}\),
the antiunitary reflection \(H \mapsto -\bar{H}\) satisfies
\(e^{-i(-\bar{H})t} = \overline{e^{-iHt}}\) entrywise and so preserves
every \(|U_{ij}(t)|^2\) exactly, and the two-branch classification below
excludes further ambiguities conditional on its stated classification
hypothesis. Step 4's thermal matching is therefore the mathematically
obligatory entry point for the dimensionful input.

\textbf{Theorem (D-gauge completeness, two-branch form).} \emph{Assume
H-BG: every non-congruent pair of strictly ordered integer rulers with
distinct positive differences and equal difference sets has six marks
and, after relabeling, realizes either \(\mu\) or its global reversal,
where \(\mu\) sends each ascending mark pair of \(\{0,1,4,10,12,17\}\)
to the ascending mark pair of \(\{0,1,8,11,13,17\}\) with the same
difference. The orientation is chosen once for the whole pair. Let
\(H,H'\) act on the same finite configuration space of dimension
\(m\ge3\), with distinct nonzero energy gaps and non-vanishing
configuration-basis overlaps for both Hamiltonians. If
\(|U'_{ij}(t)|^2 = |U_{ij}(t)|^2\) for all \(i,j,t\), then
\(H' = DHD^\dagger + E_0\) or \(H' = -D\bar{H}D^\dagger + E_0\), for a
time-independent diagonal unitary \(D\) and a real \(E_0\).}

\emph{Conditional proof.} Equality of the probability families forces
equality of the per-frequency coefficients of
\(|U_{ij}(t)|^2 = \sum_{a,b} V_{ia} V_{ja}^* V_{ib}^* V_{jb}\, e^{-i(E_a-E_b)t}\).
Under the stated overlap and gap hypotheses, the spectra have equal
difference sets and the coefficient matching induces a bijection of
ordered mode pairs preserving directed gaps. A fixed correspondence can
be transferred from real to integer spectra while preserving its gap
identities, strict order, distinct differences and non-congruence: its
realization constraints are homogeneous rational-linear equalities,
strict inequalities and inequations, and rational approximation in the
equality subspace followed by denominator clearing preserves the finite
list of strict conditions. H-BG then supplies a congruent pair or one of
the two globally oriented exceptional correspondences. For congruent
spectra, diagonal-line rigidity fixes the moduli and the
phase-coboundary relation fixes the phases, giving
\(W_{i,\tau a}=d_i\beta_a V_{ia}\) with \(d,\beta\) unimodular. The
column phases cancel in the reconstructed Hamiltonian. Reflection
congruence applies this argument to \((\bar V,-E)\) and gives the second
branch. For the positive exceptional correspondence \(\mu\), the
coefficient identities force flat moduli and constrain the phase data to
a graph of maximum clique three, excluding the six orthogonal rows
required by unitarity. For the reversed correspondence, reversing the
source mode pair conjugates its coefficient tensor; applying the same
obstruction to the conjugated source eigenbasis excludes that
orientation as well. Thus the two Hamiltonian branches follow
conditional on H-BG. \(\square\)

The complete comparison of H-BG with the published integer
classification of Bekir and Golomb {[}A. Bekir and S. W. Golomb, IEEE
Trans. Inform. Theory 53 (2007) 2864--2867{]}, including all
normalizations and exceptional orientations, remains open. The unsigned
formal premise \texttt{BGIntegerClassification} is false: the strictly
ordered rulers \((0,1,4,10,12,17)\) and \((0,4,6,9,16,17)\) have
distinct positive differences and equal difference sets, are neither
translates nor reflections of one another, and satisfy the reversed
exceptional correspondence while admitting no relabelings with the
required positive orientation. The kernel implication
\texttt{twoBranch\_of\_BGClassification} therefore does not establish
reconstruction from the physical hypotheses. The example contradicts
that formal premise, not the published classification. Bohr-frequency
completeness, the real-to-integer transfer (kernel:
\texttt{real\_to\_integer\_surrogate} in
\texttt{OIBridge/TurnpikeScopeTransfer.lean}), congruent reconstruction
and the positive six-mode obstruction remain kernel-proved independently
of it. The exclusion of the reversed orientation and the assembly of the
two branches from H-BG are kernel-checked
(\texttt{twoBranch\_of\_signed\_spectral\_classification},
\texttt{twoBranch\_of\_BGSignedClassification} in
\texttt{OIBridge/PairIdentification.lean}), so the proof from H-BG is
kernel-checked and H-BG is the theorem's only unproved premise. Under
the triangle-genericity condition (G4) of the phase-locking lemma
({[}Main §3.1{]}), with which spectra that have distinct gaps and equal
gap sets are congruent, the two branches follow without H-BG (kernel:
\texttt{twoBranch\_of\_triangleGeneric}). The listed ambiguities are
exactly invisible: \(H+E_0\) contributes only a global phase, while
\(-\bar H\) conjugates the propagator entrywise. The gap condition is
used where coefficient matching induces a bijection of ordered mode
pairs: with degenerate gaps one frequency's coefficient is a sum over
several mode pairs. The \(C_4\) ring of {[}SM{]} Proposition 20 has
distinct eigenvalues, uniform overlaps and degenerate gaps; its
conjugate \(\bar H\) reproduces every \(|U_{ij}(t)|^2\) and conjugates
the loop product \(H_{01}H_{12}H_{23}H_{30}\), which no diagonal
conjugation changes. The ring therefore does not exhibit a failure,
without the gap condition, of the two-branch form conditional on H-BG:
\(\bar H\) lies in the reflection branch,
\(\bar H = -D_s\bar H D_s^\dagger\) with
\(D_s = \mathrm{diag}(1,-1,1,-1)\). The raw frequency set alone cannot
distinguish the non-congruent homometric rulers; reconstruction uses
coefficient-labelled frequency data.

\textbf{Corollary (operational antiunitary invariance).}
\emph{Transposing the preparation, every operation, and the effect
simultaneously --- \(\rho \mapsto \rho^T\),
\(\mathcal{G} \mapsto T \circ \mathcal{G} \circ T\), \(E \mapsto E^T\)
--- preserves every finite circuit probability, for arbitrary operations
and arbitrary finite sequences, hence for every fixed outcome branch of
any adaptive instrument strategy; the transformation carries Kraus
operations to Kraus operations with conjugated operators, preserves the
trace normalization, and on unitary evolution is exactly the branch map
\(H \mapsto -\bar{H}\).}

Enlarging the data from transition probabilities to arbitrary
preparations, instruments, adaptive finite circuits, and effects
therefore cannot remove the second branch: at the purely operational
level the two branches are not two empirically distinguishable theories
but two antiunitarily related reconstructions of the same complete
circuit statistics, and they become physically distinct only once an
oriented structure is fixed. (Kernel: \texttt{circuit\_invariance},
\texttt{string\_invariance}, \texttt{transposeMap\_kraus},
\texttt{kraus\_normalization}, \texttt{unitary\_channel\_transpose} in
\texttt{OIBridge/AntiunitaryInvariance.lean}.)

\textbf{Corollary (thermodynamic orientation).} \emph{For a reflected
Hamiltonian pair with distinct energies, the reflection branch inverts
thermal orientation: the antiunitary image of the positive-temperature
Gibbs state of \(H\) is the negative-temperature Gibbs state of
\(H' = -D\bar{H}D^\dagger + E_0\), and no population profile that is
non-increasing in the energies without being uniform remains so under
the reflected spectrum. One passive, non-maximally-mixed state fixed
with respect to the emergent generator therefore excludes the
antiunitary branch; positive-temperature Gibbs/KMS states are strictly
passive, hence sufficient.}

Conditional on H-BG and the reconstruction hypotheses, transition data
determine the dynamics up to ordinary gauge and one global antiunitary
orientation; adding one passive oriented state determines it up to
ordinary gauge only, the excluded alternative collapsing the theorem's
disjunction to \(H' = DHD^\dagger + E_0\). Whether the substratum
\emph{derives} such a state is not claimed here; it is the H-orientation
transport question, and it is narrower than the H-balance/KMS conditions
of §3.2 --- single-state passivity, not complete passivity and not the
Gibbs law. Its classical layer is settled: exact finite-bath counting
with \(\beta_E = \partial S_H/\partial E > 0\) makes the
counting-selected profile passive for the classical configuration
energies, with no Gibbs approximation and no thermodynamic limit
(kernel: \texttt{counting\_passive}). The open steps are the
stationarity of that state under the reconstructed evolution and the
transport of the classical energy ordering to one of the two
orientations of the reconstructed Bohr spectrum --- a sign
correspondence only, far weaker than H-spectrum. (Kernel:
\texttt{transported\_gibbs}, \texttt{gibbs\_orientation},
\texttt{passivity\_selector\_nonuniform},
\texttt{orientation\_excludes\_reflection},
\texttt{gibbs\_strictlyPassive} in
\texttt{OIBridge/ThermalOrientation.lean}.)

\textbf{The coherent-completion classification.} The completion problem
separates into three logically distinct questions --- existence,
classification, orientation --- with three different statuses, recorded
here exactly.

\emph{Existence.} The classical fixed-basis operational theory admits a
conservative quantum realization, but extension to a coherent,
visible-local, multi-time representation is not automatic. For a
stationary preparation, feasibility reduces to a convex condition
determined by the reconstructed eigenvectors and the visible readout:
\texttt{ShellRepresentationConsistency} holds precisely when the
classical marginal is a mixture of the eigenvector population columns
(kernel: \texttt{shell\_representation\_from\_comb},
\texttt{comb\_mixture\_of\_shell\_representation}, with
\texttt{uniform\_overlap\_obstruction} an explicit failure).
Visible-local interventions impose further invariants and
complete-positivity constraints, and an explicit finite carrier ---
preparation-feasible, with a globally reachable intervention table ---
admits no visible-local CPTP coherent lift at all (kernel:
\texttt{twoByTwo\_affine\_rigidity}, \texttt{twoByTwo\_nonCP},
\texttt{twoByTwo\_no\_local\_lift}). Not every nominal carrier belongs
to the coherent-completion class, and every statement below is about the
completions that exist.

\emph{Classification.} Conditional on existence, the remaining freedom
is rigid. Two-slot operational contexts separate operators --- one-slot
contexts do not, even where the accessible algebra is already the full
matrix algebra (kernel: \texttt{operational\_separation},
\texttt{sameData\_unique\_state}) --- so identical complete operational
data induce a well-defined correspondence between two completions'
accessible algebras, and the correspondence is bijective, unital,
star-preserving, and positive in both directions by construction from
the data, positivity transferring both ways through the self-duality of
the positive cone (kernel: \texttt{sameData\_orderIso},
\texttt{psd\_iff\_trace\_nonneg}). It is consequently a Jordan
\(*\)-isomorphism (kernel: \texttt{orderIso\_jordan}, proved
internally), and the finite-dimensional matrix-unit classification
leaves exactly two possibilities.

\textbf{Theorem (coherent-completion classification).} \emph{Let two
OI-compatible coherent completions on the same finite visible dimension
carry Hermitian accessible menus \(\{G^{(1)}_i\}\), \(\{G^{(2)}_i\}\),
each spanning the full matrix algebra and containing the unit at a
common index, and positive-semidefinite state families
\(\{\sigma^{(1)}_k\}\), \(\{\sigma^{(2)}_k\}\), each separating
(\(\mathrm{Tr}(M\sigma_k) = 0\) for all \(k\) forces \(M = 0\)) and each
generating the full positive cone by nonnegative finite combinations. If
the two completions carry the same complete operational data,
\(\mathrm{Tr}(G^{(1)}_i\,\sigma^{(1)}_k) = \mathrm{Tr}(G^{(2)}_i\,\sigma^{(2)}_k)\)
for all \(i, k\), then there are a linear map \(\Phi\) with
\(\Phi(G^{(1)}_i) = G^{(2)}_i\) for every \(i\) and a unitary \(W\) with
\(\Phi(X) = W X W^\dagger\) for all \(X\), or
\(\Phi(X) = W X^T W^\dagger\) for all \(X\).} (Kernel:
\texttt{sameData\_unitary\_or\_transpose} in
\texttt{OIBridge/JordanClassification.lean}, chaining
\texttt{sameData\_orderIso}, \texttt{orderIso\_jordan}, and
\texttt{matrixJordan\_unitary\_or\_transpose}; the matrix-unit chain
\texttt{orthogonal\_resolution\_rank\_one}, \texttt{corner\_form},
\texttt{corner\_nilpotent}, \texttt{corner\_unimodular},
\texttt{corner\_cocycle}, \texttt{orientation\_dichotomy} runs the same
local-freedom, triple-consistency, phase-coboundary architecture as the
Hamiltonian reconstruction of the theorem above.)

Every two existing OI-compatible coherent completions with identical
complete operational data are therefore equivalent either by ordinary
unitary gauge or by one global antiunitary orientation reversal. No
larger coherent fibre survives these conditions.

\emph{Orientation.} The residual transpose branch cannot be selected by
the unoriented operational data themselves.

\textbf{Theorem (operational orientation no-go).} \emph{Transposition
carries every admissible completion to an admissible completion ---
Hermiticity, positivity, span, separation, the unit, and the full cone
all transport --- with exactly the same complete operational data, and
realizes the transpose branch of the classification with
\(W = \mathbb{1}\) while being the conjugation action of no unitary.
Consequently every selector that factors only through the operational
data assigns the same value to the two members of an orientation pair;
and, in general, no inhabited class of completions closed under
transposition satisfies throughout a property \(P\) with
\(P(R) \Rightarrow \neg P(\Theta R)\).} (Kernel:
\texttt{transpose\_data\_eq},
\texttt{transpose\_completion\_admissible},
\texttt{transpose\_realizes\_second\_branch},
\texttt{transpose\_not\_inner}, \texttt{operational\_orientation\_noGo}
in \texttt{OIBridge/OrientationSelection.lean};
\texttt{no\_universal\_oriented\_property} in
\texttt{OIBridge/OrientationClosure.lean}.)

The two named premises are accordingly not symmetric obligations.
\texttt{ShellRepresentationConsistency} is an existence condition: it
holds for a model precisely when it holds for the reflected partner
\((\bar V, -E)\) (kernel:
\texttt{shellRepresentation\_transpose\_stable}), so it cannot select an
orientation. Orientation requires additional alignment ---
\texttt{OperationalTransitionIdentification}, or an
\texttt{OrientedShellRepresentation} in which the represented spectral
populations are non-uniform and passive with respect to the derived
classical energy ordering. Those conditions are orientation-sensitive
(kernel: \texttt{transitionIdentification\_orientation\_sensitive},\\
\texttt{orientedShellRepresentation\_orientation\_sensitive}), by the
no-go they are not consequences of transpose-symmetric operational data
alone, and either one, combined with the classification, excludes the
transpose branch (kernel:
\texttt{sameData\_unitary\_of\_orientedReference},
\texttt{sameData\_unitary\_of\_transitionIdentification},
\texttt{sameData\_unitary\_of\_shellRepresentation}). The boundary is
therefore

\[\boxed{\;\mathrm{OI} \;+\; \text{OI-compatible coherent-completion conditions} \;\Longrightarrow\; \mathrm{QM}/\mathbb{Z}_2^{\mathrm{anti}}\;}\]

for the coherent completions that exist, subject to the stated carrier,
locality, genericity, and operational-completeness hypotheses; and if
the positive classical thermodynamic orientation is additionally
identified with the coherent energy ordering --- through
\texttt{OperationalTransitionIdentification} or the oriented shell
condition ---

\[\boxed{\;\mathrm{OI} \;+\; \text{coherent completion} \;+\; \text{positive thermodynamic orientation} \;\Longrightarrow\; \mathrm{QM}\;}\]

up to ordinary unitary gauge. Two substantive questions accordingly
remain: whether the actual OI substratum lies in the class admitting an
OI-compatible coherent completion, and whether its classical
thermodynamic orientation --- the energy ordering under which exact
finite-bath counting with \(\beta_E = \partial S_H/\partial E > 0\)
makes the classical profile passive (kernel: \texttt{counting\_passive},
from a monotone bath count), an orientation of energy rather than a
direction of time --- enters that completion through an
orientation-sensitive bridge. The classification of alternative coherent
theories is otherwise closed: existing completions are ordinary complex
quantum mechanics modulo one global antiunitary \(\mathbb{Z}_2\), while
some nominal carriers admit no coherent completion at all.

\textbf{The operational-completion characterization.} The classification
above is a uniqueness statement conditional on existence. The existence
side --- which operational theories generated on a finite carrier are
exactly quantum mechanics at the level of instruments and composites ---
is settled by a characterization with no premise beyond its stated
setting. Fix a nonempty finite visible system with configuration set
\(A\). The setting takes complex matrix kinematics as given: an
operation on the system is a complex-linear map on the complex matrices
\(M_A(\mathbb{C})\) indexed by \(A\), and an operation at ancilla level
\(n\) is a complex-linear map on
\(M_{A \times [n]}(\mathbb{C}) \cong M_A(\mathbb{C}) \otimes M_n(\mathbb{C})\),
\([n] = \{1, \ldots, n\}\), the matrices indexed by a configuration
paired with one of \(n\) ancilla states. A derivation of this kinematics
from operational premises that presuppose neither a matrix
representation nor the complex field is open. An operational theory on
the system specifies which finite outcome families of such maps are
available on the system and at every finite ancilla level, and which
preparation maps \(M_A(\mathbb{C}) \to M_{A \times [n]}(\mathbb{C})\)
are available at each level, subject to six rules (kernel:
\texttt{FiniteOperationalTheory} in
\texttt{OIBridge/OperationalAssembly.lean}): (a) the identity is
available on the system; (b) coarse-graining the outcome labels of an
available family gives an available family, on the system and at every
level; (c) at every ancilla level, an available family followed by
available families chosen by its outcome is available, with the pair of
outcomes as the joint outcome --- a composition rule stated at the
ancilla levels and not on the system; (d) attaching a uniformly mixed
ancilla, \(\rho \mapsto \rho \otimes \mathbb{1}_n/n\), is an available
preparation for every \(n \geq 1\), and an available preparation
followed by an available single-outcome operation at its level is an
available preparation; (e) at every level \(n\) an available
\(n\)-outcome family reads the ancilla, its \(k\)-th branch sending
every product \(X \otimes Y\) to \(X \otimes P_k Y P_k\), with \(P_k\)
the projector on the \(k\)-th ancilla state, a spectator-independence
requirement from which the branch's form on every matrix follows
(kernel: \texttt{readout\_is\_localLuders}); (f) an available
preparation, followed by an available family at its level and then by
discarding the ancilla through the partial trace, induces an available
family on the system. No rule extends a system operation to a composite,
and no rule imposes positivity or trace preservation; those enter
through the conditions. The OI⁺, primitive-source, substratum-source,
layer-flow and typed forms below are stated over the same kinematics,
the typed form over the complex matrices of every finite carrier. Five
conditions are in play. (i) \emph{Valid probabilities}: at every ancilla
level, every available family sends positive states to positive states
and its outcomes' traces sum to the input trace. (ii)
\emph{Trivial-ancilla consistency}: an operation available on the system
is available on the system with a one-state ancilla adjoined. (iii)
\emph{Inert spectators}: an untouched independent system can be adjoined
without changing an intervention. (iv) \emph{Full reversible control}:
every finite unitary on every composite is available. (v) \emph{Iterated
composition}: a composite may itself serve as the working system of a
larger experiment --- an operation obtained by attaching a fresh
uniformly prepared ancilla to a composite, acting on the enlarged
carrier, and discarding the ancilla is available on the composite.

\textbf{Theorem (operational-completion characterization).} \emph{For
every nonempty finite visible system, the available outcome families on
the system and on every positive ancilla level are exactly the
normalized finite Kraus instruments if and only if conditions (i)--(v)
hold.} (Kernel: \texttt{exactAll\_iff\_physical\_general} and
\texttt{general\_characterization} in
\texttt{OIBridge/GeneralCarrier.lean}; necessity
\texttt{physical\_of\_exactAll}; sufficiency through
\texttt{krausSoundExt\_of\_validity\_inert},
\texttt{compositeCompleteness}, and \texttt{exactAll\_of\_levelOne}. The
one extension fact the sufficiency direction uses --- every isometry
into a composite extends to a unitary agreeing with it on the seed ---
is proved for every finite carrier from Mathlib's orthonormal-basis
extension theorem, \texttt{finiteIsometryExtensionSF\_discharged} in
\texttt{OIBridge/IsometryExtension.lean}, so no external premise enters
either direction.)

Conditions (i) and (ii) are well-formedness requirements of any
operational theory; (iii)--(v) are the substantive selection principles,
and for well-formed theories the characterization is by those three
alone (kernel: \texttt{exactAll\_iff\_substantive}). Each of the five is
independent of the other four and of the observation process itself: for
each condition there is a qubit theory that realizes the sealed OI core
--- a finite C1--C4 embedded-observation process together with its
actual visible readout --- satisfies the other four, and fails that one
(kernel: \texttt{five\_way\_minimality} in
\texttt{OIBridge/RankGapTheory.lean}, with the witnesses
\texttt{everywhereAvailable}, \texttt{countermodel},
\texttt{diagTheory}, \texttt{gapTheory}, \texttt{systemLoose}). In
particular a theory that never creates coherence realizes the OI core
and is not quantum (kernel: \texttt{oi\_alone\_not\_qm}). Hence

\[\boxed{\;\text{bare finite OI} \;\not\Rightarrow\; \mathrm{QM}\;}\]

while

\[\boxed{\;\text{OI-compatible operational theory} \;+\; \text{(i)–(v)} \;\iff\; \text{finite operational QM}\;}\]

for every nonempty finite system (kernel: \texttt{main\_result},
\texttt{oi\_compatible\_classification} in
\texttt{OIBridge/GeneralCarrier.lean}). The OI clause is not what does
the work --- every theory with full reversible control realizes the
sealed core (kernel: \texttt{realizesSealedOICore\_of\_control}) --- so
the theorem is a classification of the completions compatible with
embedded observation, not a derivation of quantum mechanics from it. The
substantive physical question it isolates is whether the OI substratum
supplies (iii)--(v). Under its stated interventions, finite bijective
read-write dynamics and the phase intervention, it does not supply (iv)
(kernel: \texttt{substratumGen\_not\_control} in
\texttt{OIBridge/StructuralClosure.lean}, on the two-state carrier;
substratum-source form below); whether an extended OI mechanism supplies
(iv), and whether the substratum supplies (iii) and (v), are open. The
action-scale calibration uses the fixed-basis representation with
H-count and H-slope on the identified operational transitions. Any
realization requiring arbitrary detector couplings inherits the open
sourcing of the operational completion (§2.3).

\textbf{Quantum-complete OI (OI⁺).} The characterization suggests a
layered extension rather than a redefinition of observation
incompleteness. Let the OI core denote the original sealed finite
observation-incompleteness condition, unchanged (kernel:
\texttt{OICore}, which is \texttt{RealizesSealedOICore}, in
\texttt{OIBridge/CompletedOI.lean}). Define OI⁺ on a nonempty finite
carrier by adjoining well-formedness --- valid probabilities and
consistency under a trivial ancilla, conditions (i) and (ii) ---
together with three substantive physical principles (kernel:
\texttt{OIPlus} in \texttt{OIBridge/CarrierGeneralOIPlus.lean}; on the
qubit carrier the OI-core conjunct is retained for its genealogical
role, \texttt{OIHierarchy.OIPlus}, and the two definitions agree there,
\texttt{oiPlus\_qubit\_iff}):

\begin{enumerate}
\def\labelenumi{\arabic{enumi}.}
\item
  \emph{Observational independence}: an available operation remains the
  same operation when an untouched system is adjoined. Equivalently,
  inert spectators do not change the operational possibilities; in
  particular, independently available operations on disjoint factors can
  be performed jointly (kernel:
  \texttt{observationalIndependence\_iff\_inert},
  \texttt{parallel\_of\_observationalIndependence}).
\item
  \emph{Reversible richness}: available reversible transformations admit
  operational inverses, and at every finite composite level there exists
  a passive drift and a finitely generated control family whose
  dynamical Lie algebra contains the special-unitary algebra. The finite
  reachability theorem then supplies exact unitary control (kernel:
  \texttt{control\_of\_reversibleRichness}; converse
  \texttt{reversibleRichness\_of\_control}), and the Lie-rank clause
  alone already does so, with no inverse of a control consumed (kernel:
  \texttt{control\_of\_lieRank} in
  \texttt{OIBridge/MicroscopicReversibility.lean}, from
  \texttt{universalReachability\_of\_lieRank\_positive} in
  \texttt{OIBridge/PositiveReachability.lean}). Thus this is a
  certificate of sufficient reversible richness, rather than a direct
  postulate that every unitary transformation is available.
\item
  \emph{Observer recursion}: a composite observable system is itself an
  admissible observable system, with its operational families inherited
  from the corresponding composite level. Equivalently, under the
  identity and relative-readout operations used by the construction,
  composites may themselves be used as systems in still larger
  experiments (kernel: \texttt{closure\_of\_observerRecursion},
  \texttt{observerRecursion\_of\_closure}).
\end{enumerate}

With these definitions, the kernel proves

\[\boxed{\;\mathrm{OI}^{+} \;\iff\; \text{exact finite endomorphic operational QM}\;}\]

on every nonempty finite carrier (kernel:
\texttt{carrier\_general\_oiPlus}). Equivalently, OI⁺ is the physically
rephrased form of completed OI, the core together with (i)--(v) (kernel:
\texttt{CompletedOI}, \texttt{oiPlus\_iff\_completedOI}): observational
independence corresponds to inert-spectator compositionality, reversible
richness to sufficient composite unitary control, and observer recursion
to iterated composition.

This equivalence is an axiomatic completion theorem, not a derivation of
the additional principles from bare OI. It also does not prove that bare
OI selects a unique relative quantum evolution: it characterizes the
completed theory, while which unitary evolution relates two times, given
a visible family admitting several coherent lifts, is a separate
question that the completion principles do not address and that remains
open ({[}Main §3.4{]}). None of observational independence, reversible
richness, or observer recursion follows from the OI core: each can fail
while the core, well-formedness, and the other two remain satisfied
(kernel: \texttt{oiPlus\_independence}). Nor is it claimed that these
are the unique natural principles from which quantum mechanics may be
selected. Indeed, because sufficient reversible control already realizes
the sealed OI core, the OI-core conjunct is logically redundant once the
full OI⁺ conditions hold (kernel: \texttt{completedOI\_iff\_physical}).
Its role in the definition is conceptual and genealogical: OI⁺
identifies a quantum-complete extension of the original OI framework
while preserving the distinction between what observation incompleteness
itself entails and what additional physical structure selects the
quantum sector. The five-condition characterization above and the OI⁺
equivalence hold at the same scope, every nonempty finite system. The
sealed OI core itself is a qubit process, so on a general carrier OI⁺
consists of well-formedness and the three principles, and by the
redundancy just noted nothing is lost by omitting the core conjunct
there.

\textbf{Primitive-source form of OI⁺.} The three physical principles of
OI⁺ admit a formulation one level deeper, in terms of physical
implementations and embedded observers rather than of the available
operations directly. On every nonempty finite carrier, exact finite
endomorphic operational quantum mechanics is equivalent to three such
principles (kernel: \texttt{carrier\_general\_oiPlusMin},
\texttt{oiPlusMin\_iff\_qm} in \texttt{OIBridge/MinimalRepertoire.lean};
with the elementary repertoire in place of the second principle,
\texttt{carrier\_general\_oiPlusPos}, \texttt{oiPlusPos\_iff\_qm} in
\texttt{OIBridge/PositivePackage.lean}):

\begin{enumerate}
\def\labelenumi{\arabic{enumi}.}
\item
  \emph{Implementation locality.} Operational interventions are those
  realized by one protocol built from admissible implementation
  operators: an admissible isometric step, the readout of a register of
  the carrier whose projectors are admissible, classical coarse-graining
  of the outcome, sequential composition with outcome-dependent
  continuation, and the discard of a uniformly attached ancilla; the
  branches of an intervention arise together as the outcomes of one
  protocol, and admissibility is invariant under relabelling and under
  adjoining an uncoupled spectator. Probability normalization in
  aggregate is a consequence of that realization rather than a clause of
  it, and the structure supplies complete positivity and observational
  independence (kernel: \texttt{instAvail\_trace},
  \texttt{observationalIndependence\_of\_implementationLocality},
  \texttt{validity\_of\_implementationLocality} in
  \texttt{OIBridge/ImplementationLocality.lean}). No closure of the
  implementations under the adjoint is assumed, and no inverse of a
  control is assumed available: the Lie-rank condition supplied by the
  second principle gives full composite unitary control on its own,
  every unitary being, up to a global scalar phase that acts trivially
  on its conjugation channel, a positive word in the flows and the
  controls (kernel: \texttt{control\_of\_lieRank} in
  \texttt{OIBridge/MicroscopicReversibility.lean},
  \texttt{universalReachability\_of\_lieRank\_positive} in
  \texttt{OIBridge/PositiveReachability.lean}; the scalar phases are
  factors of the positive monoid there, not an operation of the
  repertoire), and on a well-formed theory full control yields
  operational inverse accessibility (kernel:
  \texttt{inverseAccessibility\_of\_lieRank}).
\item
  \emph{Phase-free richness.} At every finite composite level with two
  or more states, some pair of distinguishable states is continuously
  drivable, and every exchange of two distinguishable states is
  available. No phase operation enters. On three or more states one
  driven transition and the permutations generate the full
  special-unitary Lie algebra --- permutation conjugation transports the
  drive to every pair, two drives sharing a state bracket to the
  imaginary transition on the third pair, and the real and imaginary
  transitions on a pair bracket to the population difference (kernel:
  \texttt{bracket\_XX}, \texttt{hControl\_perm} in
  \texttt{OIBridge/MinimalRepertoire.lean}) --- and thereby yield the
  Lie-rank condition and hence exact finite unitary reachability at that
  level (kernel: \texttt{control\_at\_level}, through
  \texttt{control\_of\_lieRank}). On a two-state carrier the drive and
  the exchange commute and the generated algebra reaches no population
  difference (kernel: \texttt{not\_hControl\_two}); the levels with two
  or fewer states inherit control instead from the level with three
  times as many states, through the ancilla attachment and discard that
  embedded observation supplies (kernel: \texttt{descend},
  \texttt{control\_of\_phaseFree}). The exchanges themselves are
  generated by one full cycle and one adjacent exchange (kernel:
  \texttt{phaseFree\_of\_cyclic}). Elementary transition richness ---
  every pair of distinguishable states continuously drivable, every
  exchange available, and a quarter phase on every state, the quarter
  phase supplying the imaginary partner of a transition directly --- is
  a stronger form of the same principle (kernel:
  \texttt{phaseFree\_of\_elementary}; \texttt{lieRank\_of\_elementary},
  \texttt{hControl\_star}, \texttt{bracket\_XY}), and the two forms
  characterize the same theory (kernel:
  \texttt{oiPlusMin\_iff\_oiPlusPos}).
\item
  \emph{Embedded observation.} One relabelling- and regrouping-invariant
  family of finite operational theories on all finite carriers has the
  given theory as its ambient member. This principle supplies observer
  recursion, iterated composition, and consistency under a trivial
  ancilla (kernel: \texttt{observerRecursion\_of\_embeddedObservation},
  \texttt{systemToLevelOne\_of\_embeddedObservation} in
  \texttt{OIBridge/EmbeddedObservation.lean}).
\end{enumerate}

The equivalence

\[\boxed{\begin{aligned}&\text{implementation locality} + \text{phase-free richness}\\ &\qquad + \text{embedded observation} \;\iff\; \text{exact finite endomorphic operational QM}\end{aligned}}\]

holds on every nonempty finite carrier. It refines the OI⁺
characterization, and it is not a claim that bare observation
incompleteness entails these principles: the normalization half of
operational validity is a consequence of implementation generation, an
intervention realized by one protocol preserving probability in
aggregate, and neither principle is supplied by the others, the
witnesses being stated for the Lie-rank condition that phase-free
richness supplies (kernel: \texttt{redundancy\_fails},
\texttt{lieRank\_not\_redundant}). The nonclassical repertoire the
equivalence consumes is one continuously driven transition together with
the exchanges of distinguishable states; the quarter phase is
dispensable at every level. The exchange clause is not reducible to the
powers of a single cyclic permutation: on an even carrier one driven
transition with one full cycle generates a conjugate of the real
antisymmetric matrices, which reaches no population difference (kernel:
\texttt{not\_hControl\_evenCycle}), so one full cycle together with one
adjacent exchange is the two-element discrete form the kernel
formalizes, sufficient for every exchange (kernel:
\texttt{phaseFree\_of\_cyclic}), with no claim that a smaller discrete
resource is excluded, and no stronger minimality of the one driven
transition is claimed. Whether that driven transition can be supplied by
a discrete resource is the question of the substratum-source form below.
Inverse accessibility is not a hypothesis of the equivalence: Lie-rank
richness gives full composite unitary control unconditionally, and on a
well-formed theory full control yields inverse accessibility,
well-formedness being supplied by implementation locality with embedded
observation (kernel: \texttt{inverseAccessibility\_of\_lieRank}). An
implementation class closed under the adjoint remains a property of
exact quantum mechanics (kernel:
\texttt{reversibleImplementationLocality\_of\_qm} in
\texttt{OIBridge/MicroscopicReversibility.lean}), and, with the other
two principles, the class with that closure and the class without it
characterize the same theory (kernel:
\texttt{oiPlusPos\_iff\_oiPlusElem}).

\textbf{Substratum-source form.} The three primitive-source principles
are jointly supplied by one object: an implementation architecture --- a
class of admissible implementation operators closed under the operations
of a finite operational theory --- that is stable under an uncoupled
spectator, invariant under relabelling of the carrier, closed under the
adjoint, and drives the elementary transitions. A theory generated by
such an architecture is exact finite endomorphic operational quantum
mechanics on every nonempty finite carrier, and that quantum mechanics
is generated by one (kernel:
\texttt{genTheory\_qm\_of\_quantumArchitecture},
\texttt{qm\_generated\_by\_quantumArchitecture} in
\texttt{OIBridge/SubstratumSource.lean}). The concrete OI substratum
under A1 --- finite states and bijective microscopic dynamics ---
together with one further intervention principle, that a relative phase
on a single state is a selectable intervention at every level, supplies
the structural part in full. The phase intervention enters as an
assumption on the substratum class, not as a consequence of the finite
states with the bijective read-write dynamics: every isometry that
dynamics and the stated observer access supply carries the all-ones
vector to a multiple of itself; on every finite carrier with two or more
states, the theory generated by an architecture whose admissible
isometries do so fails the requirement of a quarter phase on every state
at every level, the failure being exhibited at level one, with no claim
that the quarter phase is unavailable at each composite level; and the
quarter phase carries that vector off its ray (kernel:
\texttt{permClass\_onesFixing},
\texttt{onesFixing\_not\_phasesAvailable},
\texttt{substratumClass\_not\_onesFixing} in
\texttt{OIBridge/PhaseSource.lean}). An operator carrying the all-ones
vector off its ray is therefore necessary to any access that supplies
the phases, and is not shown sufficient; the access assumed is the
quarter-phase intervention itself, a member of the class as stated
(kernel: \texttt{phase\_monomial} in
\texttt{OIBridge/SubstratumInterface.lean}). The interventions of the
substratum so stated have monomial observable operators, a permutation
of a diagonal (kernel: \texttt{bijectiveOperator\_monomial},
\texttt{phaseOperator\_monomial} in
\texttt{OIBridge/SubstratumInterface.lean}), and that class, with
nothing added, is closed under composition, coarse-graining, and ancilla
blocks, stable under an uncoupled spectator, invariant under
relabelling, and closed under the adjoint --- the reversal of a
bijective intervention is the inverse bijection and the reversal of a
phase is its conjugate --- with the exchanges of distinguishable states
and the phases among its members (kernel:
\texttt{substratumClass\_structurallyClosed},
\texttt{bijectiveOperator\_conjTranspose},
\texttt{phaseOperator\_conjTranspose} in
\texttt{OIBridge/StructuralClosure.lean}). The stated monomial
implementation class fails elementary drivability. Its generated theory
on the two-state carrier lacks composite unitary control. A monomial
conjugation preserves diagonal states, whereas a continuously driven
transition passes through operators with two nonzero entries in one row.
The specified reversible read-write family induces permutation operators
at every parameter value. A strict interpolation toward the exchange is
not monomial. In the nontrivial memory-swap example, the selected states
exchange and the operator remains monomial (kernel:
\texttt{readWriteOperator\_eq\_perm},
\texttt{offDiagonal\_interp\_not\_monomial},
\texttt{readWriteControl\_independent} in
\texttt{OIBridge/ReadWriteControl.lean}; \texttt{substratum\_residual}
in \texttt{OIBridge/StructuralClosure.lean}). A structurally closed
extension of the substratum class is a quantum architecture exactly when
it drives the elementary transitions. An extension meeting both
conditions generates exact finite endomorphic operational quantum
mechanics on every nonempty finite carrier, and quantum mechanics is
generated by such an extension (kernel:
\texttt{substratum\_extension\_quantum\_iff\_drives},
\texttt{substratum\_plus\_control\_qm},
\texttt{qm\_generated\_by\_substratum\_extension}):

\[\boxed{\begin{aligned}&\text{current OI substratum} + \text{continuous off-diagonal controllability}\\ &\qquad \iff\; \text{exact finite endomorphic operational QM}\end{aligned}}\]

on every nonempty finite carrier. This is not a derivation of quantum
mechanics from bare observation incompleteness. The controllability
resource is not entailed by A1--A6; it is an empirical extension of the
current substratum, entering as a hypothesis on the extended
architecture, and the phase intervention enters as an assumption on the
substratum class itself, the left side of the box being that class with
the phase intervention stated. The scope is finite and endomorphic, on
nonempty finite carriers, and the elementary control repertoire that
enters is the one stated above, with no claim that it is mathematically
minimal. Observation incompleteness alone admits many theories. The
control equivalence concerns structurally closed extensions with the
stated phase assumption. Physically sourcing the elementary drive
remains open.

\textbf{Layer-flow form.} The controllability resource of the
substratum-source form can be taken smaller. Grant the structural
architecture the substratum supplies with the phase intervention stated:
an implementation class closed under the adjoint generating the
interventions, embedded observation, and at every level the exchanges of
distinguishable states, the quarter phases and the read-write operators.
Within that closure, exact finite endomorphic operational quantum
mechanics holds exactly when the continuous layer flow of one involution
of the configuration space that moves some configuration, the ancilla a
spectator, is an available operation at every level and every
intermediate time:

\[\boxed{\begin{aligned}&\text{within the substratum's closure with the phase intervention stated,}\\ &\qquad \text{exact finite endomorphic operational QM} \;\iff\; \text{one nontrivial layer flow executable}\end{aligned}}\]

for each such involution (kernel:
\texttt{derivedOI\_qm\_iff\_layerFlowExecutable\textquotesingle{}} in
\texttt{OIBridge/DerivedQ3.lean}). The executability is necessary
because full unitary control makes every unitary available, the layer
flow included (kernel: \texttt{layerFlowExecutable\_of\_control} in
\texttt{OIBridge/LiftAudit.lean}). It is sufficient because the square
of the quarter phase on the image of a moved configuration is a sign
flip, and the layer flow, times its time reversal conjugated by that
sign flip, is the continuous transition flow on the moved pair at every
angle; with the exchanges, that is phase-free richness, and phase-free
richness within the closure is quantum mechanics (kernel:
\texttt{phaseGate\_mul\_self}, \texttt{gateFlow\_isolation\_flip},
\texttt{phaseFree\_of\_derivedOI\_layerFlowExecutable},
\texttt{qm\_of\_derivedOI\_layerFlowExecutable}). For the concrete
substratum factorization, the swap-layer flow on a nonempty region with
a nontrivial alphabet is the one-parameter unitary group of one of its
gates and is an instance of this equivalence (kernel:
\texttt{derivedOI\_qm\_iff\_layerFlowExecutable\_swap\textquotesingle{}});
the equivalence itself is stated for each involution with a moved
configuration, whether or not it is a gate of that factorization. The
layer flow reaches the driven elementary transitions of the
substratum-source form through the closure, so the two boxes have the
same right side and the resource named here is the smaller one. Two
qualifications are part of the statement. The phase intervention is an
assumption on the substratum class, as stated above, and not a
consequence of the finite states with the bijective read-write dynamics.
The executability of the layer flow at intermediate times is likewise
not derived: the theory generated by any class of monomial operators,
the substratum's included, makes no layer flow of an involution with a
moved configuration available at time one half, where the flow does not
preserve diagonal states (kernel:
\texttt{gateFlow\_half\_not\_preservesDiag},
\texttt{configurationLevel\_not\_layerFlowExecutable},
\texttt{substratumTheory\_not\_layerFlowExecutable} in
\texttt{OIBridge/LiftAudit.lean}). Nor is that executability sourced by
the substratum's own dynamics with the observer's stated access: the
theory the read and write access generates on the configuration space
executes no layer flow of an involution with a moved configuration, its
own shear and swap layers included, the layers being available at time
one and not at intermediate times, and any access that does execute one
supplies a non-monomial admissible operator at level one (kernel:
\texttt{obs\_not\_layerFlowExecutable},
\texttt{waveSubstratum\_not\_layerFlowExecutable\_swap},
\texttt{exists\_nonMonomial\_of\_layerFlowExecutable} in
\texttt{OIBridge/ExecSource.lean}). The observer-level lift of {[}SM
§4.1{]}, in every formulation stated that lands in the operational
interface, derives neither such an operator nor the layer flow: the
reversible coherent lift of a permutation is monomial, every coherent
completion, reversible or not, preserves diagonal states, the explicit
Kraus realization of the projected observer operator is
configuration-level, and a generator whose time-one exponential is
available does not make the exponential available at intermediate times
(kernel: \texttt{reversibleExtension\_conj\_monomial},
\texttt{correlationExtension\_preservesDiag},
\texttt{coherentLiftClass\_eq\_substratumClass},
\texttt{stochasticChannel\_kraus\_monomial},
\texttt{avail\_one\_not\_layerFlowExecutable} in
\texttt{OIBridge/LiftSource.lean}); whether every realization of those
channels is configuration-level is not decided, and the site-space
wave-operator formulation is unresolved rather than excluded, no
operational map from it to the configuration carrier being stated.
Fractional-time access to one layer flow is therefore an additional
physical intervention assumption, as the phase intervention is:

\[\boxed{\text{stated substratum dynamics with observer access} \;\not\Rightarrow\; \text{one nontrivial layer flow executable}}\]

The equivalence displayed above is exact; the physical sourcing of this
one of its hypotheses is not. No uniqueness or minimality is claimed for
the pair of assumptions; a different pair could be sufficient as well.

\textbf{Typed form.} The statements above are endomorphic: every
operation acts from a finite carrier back to the same carrier, the
levels being \(A \times \mathrm{Fin}\,n\). Finite-dimensional quantum
mechanics also has operations between different carriers ---
preparations, isometric embeddings, partial traces, and instruments
whose input and output dimensions differ. A typed finite operational
theory has an availability predicate on finite outcome families of maps
\(M_S \to M_{S'}\) between any two finite carriers, with the closure
rules already used above stated at their carrier-general type: identity,
classical coarse-graining, feed-forward composition across carriers,
relabelling along carrier bijections, attachment of a uniformly mixed
fresh factor, discard of a factor, and a spectator-independent factor
readout whose form is derived; no clause asserts a dilation (kernel:
\texttt{TypedOperationalTheory} in
\texttt{OIBridge/TypedCompletion.lean}). Restricting a typed theory to
maps from a carrier to itself yields a finite operational theory on
every carrier, its endomorphic shadow, which is an embedded-observation
family without further assumption (kernel:
\texttt{shadow\_embeddedObservation}). The shadow is exact finite
endomorphic operational quantum mechanics on every nonempty finite
carrier --- the conclusion above, carrier by carrier --- exactly when,
between any two nonempty finite carriers, an outcome family is
typed-available precisely if it is a finite typed Kraus instrument,
rectangular Kraus operators normalized on the input carrier (kernel:
\texttt{typed\_determined\_iff},
\texttt{typed\_determined\_of\_oiPlusMin},
\texttt{typed\_determined\_of\_oiPlusPos} in
\texttt{OIBridge/TypedPositive.lean}). Soundness wraps a typed operation
into the register formed by the input and output carriers and compresses
the square Kraus operators supplied by exactness; completeness realizes
a typed Kraus instrument by attaching a uniform ancilla, relabelling,
running a register instrument that places each Kraus operator on the
output factor and records its index on the ancilla, and discarding; the
converse restricts to one carrier. The interface itself carries no
quantum content: the typed theory whose branches preserve diagonal
matrices satisfies every rule and is not quantum mechanics (kernel:
\texttt{typed\_interface\_not\_quantum}). Hence

\[\boxed{\begin{aligned}&\text{under the carrier-general typed operational interface,}\\ &\text{current OI substratum} + \text{continuous off-diagonal controllability}\\ &\qquad \iff\; \text{exact finite-dimensional typed operational QM}\end{aligned}}\]

for nonempty finite input and output carriers. Within the natural
carrier-general extension of the operational rules already used above,
``endomorphic'' is a typing artifact rather than an additional physical
restriction. The qualification stands beside the statement: uniform
attachment and discard are not derived from the endomorphic formulation,
which cannot express a map between different carriers; they are the
carrier-general counterparts of the ancilla attachment and discard rules
already present in the finite operational framework, and the
determination theorem shows that adding them introduces no specifically
quantum content. No pure-state preparation, no carrier-coherence
condition, and no availability of all completely positive maps is
assumed; the quantum content lies entirely in the endomorphic conclusion
above. The scope of this statement is finite-dimensional.

\textbf{Infinite-region completion.} The finite theory admits a
canonical infinite-region completion without a continuum limit and
without a Hilbert-space representation. The substratum holds its
fixed-spacing lattice fundamental, so the infinite-system construction
it supplies is the directed system of finite spatial regions. For a
finite region \(\Lambda\) the observable algebra is the full matrix
algebra on the region's finite configuration space; inclusion of regions
gives compatible injective unital \(*\)-homomorphisms, and observables
supported on disjoint regions commute (kernel: \texttt{inclObs\_mul},
\texttt{inclObs\_injective} in \texttt{OIBridge/QuasilocalAlgebra.lean},
\texttt{emb\_comm\_of\_disjoint} in\\
\texttt{OIBridge/QuasilocalCharacterization.lean}). The local algebra is
the algebra of equivalence classes of finite-region observables, two of
them agreeing exactly when they agree after inclusion into a common
region (kernel: \texttt{emb\_eq\_iff}); the inclusions are isometric for
the operator norm, and the norm completion of the local algebra is a
quasilocal \(C^{*}\)-algebra that is the closure of the union of the
finite stages (kernel: \texttt{norm\_inclObs},
\texttt{instCStarAlgebraQuasilocal}, \texttt{closure\_iUnion\_stage}).
Every consistent family of finite-region density matrices extends
uniquely to a state of that algebra, and the reversible finite-range
substratum update extends uniquely to an isometric \(*\)-automorphism
under which an observable of a region is an observable of an explicit
finite region after each step (kernel: \texttt{quasiState\_unique},
\texttt{quasiState\_nonneg}, \texttt{heisQ\_mul}, \texttt{norm\_heisQ},
\texttt{heis\_iterate\_emb}).

The completion is characterized independently of the construction. A
quasilocal system is a \(C^{*}\)-algebra carrying the same finite-region
matrix algebras through compatible injective unital \(*\)-embeddings,
with disjoint-region commutation and dense local stages; the definition
names no completion, no state, and no representation (kernel:
\texttt{QuasilocalSystem} in
\texttt{OIBridge/QuasilocalCharacterization.lean}). Any such system is
canonically and uniquely \(*\)-isomorphic to the region completion, by
the universal property of the local algebra and of its completion, and a
system carrying the same substratum-induced discrete dynamics receives a
canonical isomorphism intertwining the automorphisms (kernel:
\texttt{canonEquiv}, \texttt{canon\_unique},
\texttt{systemEquiv\_unique}, \texttt{canon\_dyn},
\texttt{systemEquiv\_dyn}). Hence

\[\boxed{\begin{aligned}&\text{the canonical infinite-region completion of } \mathrm{OI}_Q\\ &\qquad \cong\; \text{the unique quasilocal fixed-lattice } C^{*}\text{-system carrying}\\ &\qquad\quad\; \text{the substratum's local stages and OI-induced discrete dynamics}\end{aligned}}\]

where \(\mathrm{OI}_Q\) denotes the quantum-completed substratum of the
statements above --- the current OI substratum together with continuous
off-diagonal controllability --- and not bare observation
incompleteness. The statement removes the finite-region cutoff for the
algebra, the state space, and the OI-induced discrete dynamics; it is
not a characterization of all infinite-dimensional quantum instruments,
nor of all locality-preserving quantum dynamics.

The restriction to the substratum-induced dynamics is a theorem rather
than a convention. General locality-preserving automorphisms form a
strictly larger class: conjugation by a phase unitary at one site
preserves locality and is induced by no reversible finite-range
substratum bijection, multiplying a single-site matrix unit by \(i\)
where every transported observable carries entries that are sums of
entries of the original (kernel: \texttt{phase\_localityPreserving},
\texttt{phaseQ\_ne\_heisQ}).

Continuous time is separate. The discrete substratum dynamics does not
determine a continuous one-parameter interpolation: two Hermitian
generators give flows agreeing at every integer time and differing at
\(t = 1/2\) (kernel: \texttt{continuous\_extension\_not\_unique} in
\texttt{OIBridge/RegionLimit.lean}). What the substratum does determine
is a path. The second-order update, in its phase-space form on the pair
of the previous and current slices, factors exactly as a depth-two local
circuit, a shear layer followed by a swap layer, each a product of
commuting involutive gates, one per site, its reversibility a
consequence of the factorization (kernel:
\texttt{leap\_eq\_swap\_shear}, \texttt{gate\_comm},
\texttt{depth\_two\_circuit}, \texttt{leap\_leap\_symm} in
\texttt{OIBridge/SecondOrderCircuit.lean}). Each layer is the time-one
map of a strongly continuous one-parameter group of isometric
\(*\)-automorphisms of the quasilocal algebra, driven gate by gate
(kernel: \texttt{layerQ\_add\_time} in
\texttt{OIBridge/SecondOrderLayer.lean}, \texttt{swapQ\_add\_time} in
\texttt{OIBridge/SwapLayer.lean}, \texttt{layerQ\_one\_eq\_heisQ},
\texttt{swapQ\_one\_eq\_heisQ} in
\texttt{OIBridge/SecondOrderDrive.lean}), and the two flows, run in the
order the Heisenberg picture forces, the swap first, give a strongly
continuous path of isometric \(*\)-automorphisms from the identity to
the update's own Heisenberg action (kernel: \texttt{heisQ\_of\_comp},
\texttt{driveQ\_isContinuousPath}, \texttt{driveQ\_smul},
\texttt{driveQ\_one\_eq\_heisQ} in
\texttt{OIBridge/SecondOrderDrive.lean}). For the composite path no
one-parameter-group law is established and no generator is exhibited:
whether the path satisfies a one-parameter-group law, and whether one
time-independent finite-range interaction has the update as its
unit-time map, are open, and nothing here decides either. A
continuous-time Hamiltonian law is additional structure where that
formulation is wanted, and the equivalence above does not use it.

The lattice spacing stays fundamental. Nothing above asserts a spatial
continuum limit or a refinement of the spacing; the continuum
description used elsewhere in the framework is an effective
calculational approximation rather than the object completed here. The
uniqueness is among systems built from the substratum's finite-region
matrix algebras rather than among all \(C^{*}\)-algebras, and no
Hilbert-space representation and no superselection sector is selected at
the level of the laws, a sector choice being an input of the
initial-condition kind (kernel: \texttt{state\_selection\_audit} in
\texttt{OIBridge/RegionTower.lean}).

Every member of the residual family is invisible to the transition data
for a stated reason --- diagonal rephasing moves no modulus, the energy
origin is a global phase, the antiunitary reflection conjugates \(U(t)\)
entrywise --- and none carries a scale, so §3.2's determination of
\(\hbar\) does not wait on the theorem. Cross-interval transition
probabilities fix the relative gauge up to a single global phase; in the
continuum limit this becomes an unobservable smooth function
\(e^{i\theta(t)}\).

\textbf{Dimensional obstruction.} The unitary \(U(t)\) is dimensionless;
\(\hbar\) enters only via
\(\hat{H} = i\hbar \, \partial_t U \cdot U^\dagger\). No dimensionless
data can fix a dimensionful constant --- Step 4's thermal
self-consistency provides the necessary dimensionful input.

\textbf{Remark (phase recovery under discrete sampling).} Where the
substratum is modelled as a continuous Hamiltonian flow ({[}Main,
§2.3{]}), the Liouville marginalization of that flow (Lemma 3) gives
\(T_{ij}(t)\) as a continuous function of \(t\); the discrete update
itself determines no continuous interpolation (above). The conditional
D-gauge theorem, with H-BG and its other hypotheses, applies to the
continuous object without aliasing corrections. Aliasing would arise
only if the observer sampled \(T_{ij}\) at discrete intervals, in which
case trans-Planckian frequencies could be misidentified --- but these
lie outside the emergent QFT's domain.

\begin{center}\rule{0.5\linewidth}{0.5pt}\end{center}

\hypertarget{self-consistency-and-the-discreteness-scale}{%
\subsection{4. Self-Consistency and the Discreteness
Scale}\label{self-consistency-and-the-discreteness-scale}}

Rearranging \(\hbar = c^3 \epsilon^2 / (4G\sigma)\):

\[\epsilon^2 = \frac{4\sigma\hbar G}{c^3} = 4\sigma\,l_p^2\]

The framework yields \(\epsilon = 2\sqrt{\sigma}\,l_p\) exactly, which
is \(\epsilon = 2\,l_p\) at \(\sigma = 1\) (§3.1). This is algebraically
equivalent to the formula for \(\hbar\): the framework contains one free
geometric parameter (\(\epsilon\)), and thermal self-consistency
establishes its relationship to \(\hbar\). The result is a
self-consistency condition, not an independent determination, and with
\(\sigma\) undetermined it fixes only the ratio
\(\epsilon/\sqrt{\sigma} = 2\,l_p\). A Planck-scale \(\epsilon\)
additionally requires \(\sigma\) of order one. The bound
\(\sigma \leq \ln q\) limits \(\sigma\) from above for a \(q\)-state
mode, while the count bounds it from below only by \(\sigma > 0\):
\(\sigma = 10^{-12}\) gives \(\epsilon = 2 \times 10^{-6}\,l_p\). The
gap equation does not exclude discreteness scales far from \(l_p\);
whether they are excluded on independent grounds is open. The
framework's predictive content comes from the relationship
\(\hbar = c^3 \epsilon^2/(4G\sigma)\), the Bekenstein-Hawking formula,
and the consequences that follow.

\begin{center}\rule{0.5\linewidth}{0.5pt}\end{center}

\hypertarget{the-bekenstein-hawking-entropy}{%
\subsection{5. The Bekenstein-Hawking
Entropy}\label{the-bekenstein-hawking-entropy}}

With \(\epsilon=2\sqrt{\sigma}\,l_p\) from the stated calibration, the
boundary count gives:

\[N_{\text{modes}} = \frac{A}{\epsilon^2} = \frac{A}{4\sigma\,l_p^2} = \frac{S_{\text{dS}}}{\sigma}\]

where \(S_{\text{dS}} = A/(4\,l_p^2)\) is the Bekenstein-Hawking entropy
of the de Sitter horizon. The factor of \(1/4\) in the
Bekenstein-Hawking formula --- introduced by Bekenstein (1973) with a
coefficient he could estimate only heuristically, and fixed at \(1/4\)
by the Hawking temperature and the first law (Hawking 1975) --- is here
recovered, under the conditions of the §3 calibration: each area
\(\epsilon^2 = 4\sigma\,l_p^2\) of the horizon carries one boundary mode
of entropy \(\sigma\), and \(\sigma A/\epsilon^2 = A / (4\,l_p^2)\)
recovers the standard formula exactly, whatever \(\sigma\). The factor
of \(4\) in \(\epsilon^2 = 4\sigma l_p^2\) traces directly to §3 Step
4's thermal matching: the ratio \(2\pi/8\pi = 1/4\) between the
Euclidean KMS period (numerator: thermal periodicity of QFT on a Killing
horizon) and Jacobson's Einstein-equation prefactor (denominator:
gravity's coupling in \(dE = (c^2\kappa/8\pi G)dA\)). The
Bekenstein-Hawking \(1/4\) is therefore not an arbitrary constant but
the consistency condition between gravitational and quantum thermal
physics, conditional on H-slope (§3.2), H-count (§2.3), H-frame (§2.3),
the H-Hawking conditions of §8.5, local boost structure at the horizon
and the symmetric detailed-balance rate ratio (§8.5). For the
single-link lattice count, this application requires a sphere in the
lattice frame or an equal area-averaged link density (§3.1). Applying
the same calibration to a black-hole horizon also requires its stated
thermal and frame conditions.

\textbf{Classical area test (GW250114).} Under Hawking's classical area
theorem, the remnant horizon area is bounded below by the sum of the
progenitor areas. Inspiral and ringdown estimates test that area
comparison. GW250114 was observed on January 14, 2025 and reported in
September. The fiducial pre- and post-merger Kerr inference supports an
area increase at \(4.4\sigma\) {[}5, Fig. 5{]}. The significance depends
on the analysis cut and mode assumptions. This inference tests the
classical area theorem. The \(1/4\) entropy coefficient remains a
separate thermal calibration. The coefficient, fixed here as the exact
ratio \(2\pi/8\pi\) between Euclidean KMS periodicity (thermal physics)
and Jacobson's Einstein-equation prefactor (gravitational coupling), is
the ratio of the two normalizations that the calibration matches,
conditional on H-slope, H-count, the horizon and frame conditions, local
boost structure at the horizon and the symmetric detailed-balance rate
ratio (§8.5), with no direct empirical test at present; the area
comparison uses the Kerr inference of {[}5{]}. Untested is not
untestable: the coefficient is equivalent to the Hawking-temperature
normalization (\(T = \hbar\kappa/2\pi c k_B\) plus the first law fixes
it), so any measurement of a horizon's radiation would fix it directly.
No known astrophysical source is accessible --- a solar-mass black hole
radiates at \(\sim 60\) nK against the 2.7 K CMB --- but one physically
contingent channel exists: the final evaporation burst of a primordial
black hole (\(M_i \sim 5 \times 10^{14}\) g expiring today), whose burst
light curve would measure \(T(M)\) and hence the coefficient. Dedicated
all-sky burst searches are active --- Fermi-LAT {[}45{]}, HAWC {[}46{]},
and most recently LHAASO's all-sky search {[}47{]} --- with only
rate-density limits so far, and next-generation sensitivity projected
for CTAO {[}48{]}; detection is contingent on PBHs existing in the
relevant mass window. Analogue-gravity horizon experiments measure the
kinematic (\(2\pi\)/KMS) side of the ratio but not the gravitational
(\(8\pi G\)) side. The ongoing LVK observing runs will accumulate
additional events testing the area theorem across a range of progenitor
masses and spins.

\begin{center}\rule{0.5\linewidth}{0.5pt}\end{center}

\hypertarget{dissolution-of-the-cosmological-constant-problem}{%
\subsection{6. Dissolution of the Cosmological Constant
Problem}\label{dissolution-of-the-cosmological-constant-problem}}

\hypertarget{the-two-levels-of-description}{%
\subsubsection{6.1 The two levels of
description}\label{the-two-levels-of-description}}

\textbf{From finite-dimensional QM to QFT.} Part I delivers an exact
fixed-basis unitary/Born representation of the observer's finite-horizon
law on a finite-dimensional Hilbert space; whether an operational
theory's available instruments on a finite carrier and its composites
are the quantum instruments is a separate question, which §3.3
characterizes, within complex matrix kinematics taken as given, by its
completion conditions. For a visible lattice region with product
configuration space \(\mathcal C_V=\prod_{k\in V}\mathcal C_k\), the
site factors give \(\mathcal H_V=\bigotimes_{k\in V}\mathcal H_k\). Here
\(N=|V|\) counts visible sites. The area count \(A/\epsilon^2\) counts
boundary modes. The proposed lattice-regularized field description uses
the site spacing \(\epsilon_{\rm lat}=\sqrt{6\sigma}\,l_p\) on the
calibrated spherical single-link branch (§3.1), with local continuous
observer dynamics conditional on H-local-lift and the stated dynamical
inputs.

\textbf{Locality preservation --- conditional (H-local-lift).} A
radius-1 reversible discrete map does not by itself fix a local
continuous generator. Take four binary sites on a ring with the one-site
shift \(\varphi(x_1,x_2,x_3,x_4)=(x_4,x_1,x_2,x_3)\) --- maximally
local, each output site reading one nearest neighbour --- and restrict
\(U_\varphi\) to the one-particle orbit, where it is the \(4\)-cycle.
Writing \(U_\varphi=e^{-iH}\), every Hermitian branch has
\(\lambda_k=\pi k/2+2\pi n_k\), so
\(H_{0,2}=\tfrac14(\lambda_0-\lambda_1+\lambda_2-\lambda_3)=-\pi/4+(\pi/2)m\),
which never vanishes (\(m\) would have to be \(\tfrac12\)); the minimum
modulus is \(\pi/4\). But \(|1000\rangle\) and \(|0010\rangle\) differ
at opposite sites, and a Hamiltonian of onsite and nearest-neighbour
terms has zero matrix element between them. So no branch of the
logarithm is nearest-neighbour. The infinitesimal argument
\(T_{xx'}(dt)=|H_{x'x}|^2dt^2+\cdots\) presupposes observer-level
locality rather than deriving it from locality of \(\varphi\); {[}Main
§3.2{]} already records that a finite permutation does not determine a
continuous flow. The theorem below therefore holds under
\textbf{H-local-lift}: the selected continuous observer generator is
local or quasi-local. The normalized wave/Laplacian branch of {[}SM
§4.1{]} supplies a plausible such generator, so this is a missing
hypothesis, not an impossibility. Under H-local-lift: \emph{If the
classical Hamiltonian is spatially local (couples only neighbors), then
the emergent quantum Hamiltonian inherits spatial locality.}

\emph{Proof.} For infinitesimal \(dt\) and \(x' \neq x\):
\(T_{x,x'}(dt) = |\langle x'|\hat{H}|x\rangle|^2\,dt^2 + O(dt^3)\). If
\(x, x'\) differ at non-neighboring sites, H-local-lift gives
\(T_{x,x'}(dt) = 0\) at leading order, hence
\(|\langle x'|\hat{H}|x\rangle|^2 = 0\). The same leading-order
coefficient fixes this zero for every Hermitian generator realizing the
same continuous transition data, independently of the classification
hypothesis, so \(\langle x'|\hat{H}|x\rangle = 0\) in any gauge. The
emergent Hamiltonian has the form
\(\hat{H} = \sum_k \hat{h}_k + \sum_{\langle k,l\rangle} \hat{h}_{kl}\).
\(\square\)

\textbf{Scope of the emergent QFT.} A local lattice field theory assigns
a zero-point contribution \(\sim\hbar\omega/2\) per mode with a
wave-number cutoff of order \(\epsilon_{\rm lat}^{-1}\). The
normalization in §§2--5 uses the horizon first law, count and matched
slopes. Treating the zero-point sum as gravitationally decoupled
requires the separate coupling of §6.2.

\textbf{Supplied classical geometry.} On the stated horizon branch,
\(T_{\rm cl}\) of §3.1 equals the hidden-sector energy-shell temperature
under H-count. H-balance supplies the thermal ratio on the specified
transitions. The flat Friedmann equation gives
\(\rho_{\rm crit}=3c^2H^2/(8\pi G)\). The observed vacuum density is a
supplied fraction of this scale, with
\(\rho_{\rm vac}=\Omega_\Lambda\rho_{\rm crit}\) in the reference
cosmology.

\textbf{Emergent QFT} (what local quantum measurements probe):
Zero-point energy \(\frac{1}{2}\hbar\omega\) per mode summed to the
Planck cutoff gives \(\rho_{\text{QFT}} \sim 10^{113}\) J/m\(^3\) --- a
\(\sim 10^{122}\) discrepancy with the observed value {[}6, 7, 8{]}.

\hypertarget{why-gravity-sees-the-classical-value}{%
\subsubsection{6.2 Why gravity sees the classical
value}\label{why-gravity-sees-the-classical-value}}

The designated cosmological trace-out uses a supplied classical
spacetime metric to define its causal partition. The proposed
pre-trace-out coupling would obtain a local conserved source from the
complete classical configurations. Deriving Einstein dynamics and the
native stress or flux map belongs to G3 (§8.7).

\textbf{Scope of the causal ordering.} In this construction the causal
geometry enters before the designated trace-out. This is an ordering of
the ingredients of that construction. The finite-law representation in
{[}Main{]} applies on its own hypotheses.

\emph{(i) Definiteness.} The marginalization is defined for a specified
product partition \(\Gamma=\Gamma_V\times\Gamma_H\) {[}Main §1.4{]}. The
horizon application uses a definite causal cut in the metric supplied
for that purpose.

\emph{(ii) Non-circularity.} Null geodesics determine the causal cut in
this application. Deriving the metric solely from the output of that
same trace-out would leave its own input unspecified. A metric grounded
independently of that output can supply the cut.

\emph{(iii) Action-scale calibration.} The designated matching uses a
classical horizon geometry and the thermal slopes stated in §3.2. Area
and surface gravity cancel for the matched boundary conditions. This is
a condition on the calibration and the geometry it uses.

\emph{Scope: laboratory sources in superposition.} The designated
trace-out uses a metric at the partition boundary. That input leaves
open the gravitational field of a laboratory mass that the observer's
description places in a spatial superposition. Three completions are
compatible with them: (a) the field follows the substratum configuration
of the mass; (b) the field follows the hidden-sector average, which is
semiclassical gravity at laboratory scales, a coupling whose no-collapse
form was tested in {[}56{]}; (c) the near field belongs to the emergent
description and the classical metric of (i)--(iii) is confined to the
partition boundary, so that gravitationally mediated entanglement
between masses is admitted {[}57, 58{]}. The consequences of each
completion for the coherence of the superposition are not derived. Two
kinds of coherence loss are distinguished: environmental decoherence, in
which which-path records pass to unobserved degrees of freedom such as
emitted radiation while the total evolution remains unitary, and an
additional nonunitary effect that no environment accounts for. Any of
(a)--(c) may produce the first; whether any produces the second is open.
The selection among (a)--(c) and the derivation of its coherence
consequences are open, and together they fix the framework's prediction
for interferometric tests of gravitationally mediated entanglement and
of gravitational decoherence.

\textbf{Descriptive distinction established; dynamical decoupling
conditional on G3.} That the zero-point sum belongs to the emergent
description does not entail that it cannot source geometry: the two are
different claims. Quantum transition data are invariant under
\(\Omega_{\rm eff}\to\Omega_{\rm eff}+CI\) (a global phase), so no
ordinary observer statistic detects \(C\); yet two gravitational
completions, \(\mathcal G_{\mu\nu}=8\pi G T^{\rm cl}_{\mu\nu}\) and
\(\mathcal G_{\mu\nu}=8\pi G(T^{\rm cl}_{\mu\nu}-\rho_Cg_{\mu\nu})\),
can share the same substratum, partition, reduced transition law,
fixed-basis representation and action-scale calibration while differing
in exactly that offset. Invisibility to quantum mechanics is not
decoupling from gravity; the latter is a coupling theorem and is what G3
names as open. Accordingly the descriptive-level distinction is
established and the dissolution is conditional. The continuum
feasibility calculation compares complete, conserved stress tensors on
one metric with a fixed physical reference. A common offset then cancels
in their difference. The weak-field response retains the total energy
difference, including supports and reservoirs. The native stress or flux
map, the choice of physical reference and the response to physical
changes in vacuum remain open.
\href{../verification/programmes/gravity/drafts/g3-coupling-feasibility-20261118/G3-COUPLING-FEASIBILITY.md}{G3
coupling assessment}

\textbf{Comparisons on the designated branch.} The observed vacuum
density is a supplied fraction of the Friedmann scale. The action-scale
calibration and the entropy count are one normalization, with
\(S=\sigma A/\epsilon^2=A/(4l_p^2)\) under their stated conditions. The
Hubble-volume density comparison uses the same classical count and the
supplied Friedmann scale (§7.2). Interpreting the complement to the
baryon fraction as boundary gravity requires the source identification
and the open matter-like and vacuum-like split. The near-null Type II
reading and the Bertini et al.~comparison have the scope of §7.1. The
MOND and deep-acceleration Tully-Fisher formulas use the temperature and
heat-transfer inputs in §7.3. The cluster residual and the Milky Way and
high-redshift rotation-curve tensions remain empirical comparisons. The
Bullet Cluster displacement is a proposed lensing mechanism. In the
gauge-coupling comparison the three values at \(M_Z\) use three matched
parameters (\(\delta_0\), \(A\), \(B\)), with the derived
\(1/\alpha_0=23.25\) and threshold universality carrying the assumptions
of {[}SM §6{]}, including the open observer bundle and hypercharge
vertex.

\textbf{The classical vacuum energy scale.} The observed vacuum density
is a supplied fraction of the Friedmann scale for the reference
cosmology. Decoupling an arbitrary common offset would remove the
corresponding cancellation problem within the proposed coupling. The
value of the reference density and its dynamics remain open (§6.2).

\textbf{Quantum corrections to gravity.} The emergent Hamiltonian
reorganizes the visible transition data and is invariant in these data
under a uniform energy shift. The proposed coupling uses the conserved
classical source of the complete configuration. Its relation to emergent
energies must preserve relative-energy effects where the source varies,
with the native map and the curvature response open. Section 7.1
discusses a separate, conditional attribution of late-time running.

\hypertarget{the-structural-origin-of-the-10122-discrepancy}{%
\subsubsection{\texorpdfstring{6.3 The structural origin of the
\(10^{122}\)
discrepancy}{6.3 The structural origin of the 10\^{}\{122\} discrepancy}}\label{the-structural-origin-of-the-10122-discrepancy}}

\(\rho_{\text{QFT}} / \rho_{\text{classical}} \sim M_{\text{Pl}}^4 / (H^2/G) \sim S_{\text{dS}} = A/(4\,l_p^2)\)
--- a dimensional identity for any theory with a Planck cutoff and a
Hubble-scale density, holding up to a convention-dependent numerical
factor. With \(\rho_{\text{QFT}} = M_{\text{Pl}}^4\) and
\(\rho_{\text{classical}} = H^2/G\) (\(\hbar = c = 1\),
\(M_{\text{Pl}}^2 = 1/G\)) the ratio is exactly \(S_{\text{dS}}/\pi\);
the mode sum \(k^4/(16\pi^2)\) cut off at
\(1/\epsilon_{\rm lat} = M_{\text{Pl}}/\sqrt{6\sigma}\) (§3.1, §4)
against \(\rho_{\text{crit}} = 3H^2/(8\pi G)\) gives
\(S_{\text{dS}}/(216\pi^2\sigma^2)\); and the same mode sum cut off at
the reduced Planck mass \((8\pi G)^{-1/2}\) {[}7{]} against the observed
dark-energy density \(\Omega_\Lambda \rho_{\text{crit}}\) gives
\(S_{\text{dS}}/(384\pi^4\Omega_\Lambda) \approx 3.9 \times 10^{-5}\,S_{\text{dS}}\).
The framework reads this ratio as the Bekenstein-Hawking entropy in its
role as the information compression ratio of the trace-out; on that
reading, and conditional on the coupling step (§6.2), the ``worst
prediction in physics'' is the entropy of the observer's blind spot ---
a category error, not a fine-tuning failure. This cannot be resolved
from within, and not merely as a matter of present ignorance: any
adjudication procedure available to an embedded observer is a function
of its accessible data, and both descriptions are constructed from the
same accessible data, so no criterion grounded in the accessible
observations themselves can separate them --- non-adjudicability follows
from the equivalence itself, not from limited resources. (Rankings by
theoretical virtue are not blocked by this factoring; but they order
presentations of the same observational content rather than track
fundamentality, and the framework's own answer to the fundamentality
question is supplied at the substratum level, not adjudicated from
within either description.) A future observation could break the
equivalence, but that is falsification of the equivalence, not
adjudication between the descriptions. Nor is this factoring limitation
an artifact of the framework's particular observer architecture:
Wolpert's limits of inference {[}9{]} and T. Breuer's self-measurement
theorem {[}55{]} establish barriers of the same type for any embedded
inference device under far weaker assumptions, corroborating that the
limitation is device-general rather than model-specific (their results,
unlike the factoring argument, do assume free availability of the
device's setups). The observed vacuum density is the supplied fraction
\(\Omega_\Lambda\rho_{\rm crit}\) of the Friedmann scale. Its physical
value and dynamics remain open under the native coupling (§6.2).

\textbf{Relation to Maldacena's de Sitter observer cancellation.}
Maldacena {[}38{]} has recently shown that the unphysical phase of the
de Sitter sphere partition function --- long understood to obstruct a
clean dS thermodynamic interpretation --- cancels exactly when one
incorporates an observer with a clock. In the substratum-count
formulation the partition trace-out is the \emph{defining} operation
rather than an added ingredient, so the observer is present from the
outset; the sphere partition function and its phase are not computed
here, and whether the trace-out reproduces Maldacena's cancellation is
open. CLPW {[}39{]} establish that the von Neumann algebra of QFT
observables in a gravitational subregion promotes from type III to type
II once an observer is included. The present construction derives no
type II algebra. Its observable algebra on each finite region is the
full matrix algebra on the region's configuration space, and its
infinite-region completion (§3.3) is a quasilocal \(C^{*}\)-algebra with
no state or representation selected. No von Neumann closure of that
algebra, no relation to CLPW's observer-dressed algebra and no
identification of CLPW's entropy with a substratum quantity is derived;
whether the construction reproduces CLPW's algebra and its entropy is
open. Its horizon entropy is the entropy \(\sigma\) of one boundary mode
summed over the \(N_{\text{modes}}\) of §5, which recovers
\(S_{\text{dS}}\) under the §3 calibration. HUZ {[}40{]} argue that the
closed-universe Hilbert-space dimension is determined by the observer's
degrees of freedom; in the present framework the representation
dimensions are finite configuration counts fixed by the partition
({[}Main, §3.2{]}), and the horizon mode count
\(N_{\text{modes}} = S_{\text{dS}}/\sigma\) is fixed by the partition
geometry under the §3 calibration (§5). The post-2022
observer-essentiality literature thus shares the present construction's
starting point, the observer as part of the physics, while differing in
mechanism.

\begin{center}\rule{0.5\linewidth}{0.5pt}\end{center}

\hypertarget{predictions-and-tests}{%
\subsection{7. Predictions and Tests}\label{predictions-and-tests}}

\textbf{What the classical tests test.} A scope observation frames this
section. The classic precision tests of general relativity ---
perihelion advance, light deflection, Shapiro delay, gravitational
redshift, frame dragging --- test the vacuum Kerr and Schwarzschild
geometries: Ricci-flat \emph{solutions}, not the field equations that
select them. The tested solution set radically underdetermines the
equations: any theory whose vacuum sector contains those geometries in
the tested regimes passes identically, and theories with different field
equations exploit exactly this freedom. The radiative sector
(binary-pulsar decay, gravitational waveforms) probes the vacuum
\emph{dynamics} and discriminates more; the \emph{sourced} sector ---
cosmology, where the Friedmann geometry is nothing like Ricci-flat ---
is where the equations themselves, their coupling structure and their
response to sources, are tested directly. This framework's relationship
to that division is specific on both sides. If the emergent field
equations are the Einstein equations --- the covariant Level-G3 target,
open (§8.7) --- the framework inherits every vacuum-solution and
radiative confirmation exactly, its derivation route being invisible to
Ricci-flat tests; the established weak-field results reach Newtonian
gravity at the framework's own \(G\) (§8.7). And because the framework's
distinctive content is at the \emph{equation level} --- the
thermodynamic calibration of the horizon, the gauge status of the
uniform vacuum offset (§6), the running-vacuum channel below --- its
novel exposure concentrates precisely in the sourced sector, which is
where the predictions of this section live. The classical record's
silence about the equations is therefore not a loophole the framework
hides in but the reason its forward tests are cosmological.

\hypertarget{dark-energy-evolution-in-running-vacuum-form}{%
\subsubsection{7.1 Dark energy evolution in running vacuum
form}\label{dark-energy-evolution-in-running-vacuum-form}}

The calibrated action scale is independent of \(H\) at fixed counting
density (§3.2). The proposed vacuum-density reading uses the state
dependence of the boundary description. Any smooth density
\(\rho_\Lambda(H)\) admits the following local expansion near \(H_0>0\)
(with \(c=1\)):

\[\rho_\Lambda(H)=\rho_\Lambda(H_0)+\frac{3\nu}{8\pi G}(H^2-H_0^2)+\mathcal{O}\bigl((H^2-H_0^2)^2\bigr).\]

This is the local form used in the running-vacuum comparison {[}10,
11{]}. Setting \(X=H^2\) makes it a Taylor expansion of
\(\rho_\Lambda(\sqrt X)\) about \(X_0=H_0^2\). The counterfamily
\(\rho_0+A(H^2-H_0^2)+B(H^2-H_0^2)^2\) preserves the same linear term
and admits the Type II Bianchi bookkeeping with running \(G\) for
arbitrary \(B\). A prescribed global law or the absence of higher terms
would require a further dynamical result.

\textbf{Type II structural reading.} The implementation considered here
assumes conserved matter and transfers the Bianchi compensation for
\(\rho_\Lambda(H)\) to running \(G(H)\) {[}11{]}. The trace-out and the
boundary budget leave the covariant source and the absence of
vacuum--matter exchange open. The Type I implementations cited in
{[}11{]} instead specify fixed \(G=G_N\) and a vacuum--dark-matter
exchange such as \(Q=3\nu H\rho_{\rm dm}\).

\textbf{Coefficient and source.} The coefficient depends on the physical
attribution of the uniform part of the boundary budget and on its
variation with \(H\). The following relations have different inputs:

\begin{itemize}
\tightlist
\item
  Under the horizon first-law and count inputs, the Hubble-volume mean
  \(\rho_s(H)=3c^2H^2/(8\pi G)\) equals the supplied flat Friedmann
  density in three spatial dimensions (§7.2).
\item
  The identification of a baryonic fraction uses the supplied
  \(\Omega_B\) and a physical relation between matter and boundary
  entropy (§7.2).
\item
  The split of the remaining \(1-\Omega_B\) into vacuum-like and
  matter-like contributions and its \(H\) dependence remain open.
\end{itemize}

\textbf{Magnitude from the emergent field channel.} For the emergent
matter theory on a supplied curved semiclassical background, the
Shapiro--Solà one-loop running formula {[}33, 34{]} is

\[\nu = \frac{1}{48\pi^2}\sum_i a_i \frac{M_i^2}{M_{\rm Pl}^2}\]

with \(M_i\) the physical masses of the participating fields. For the
Standard Model spectrum (\(M_i\leq m_t\approx173\) GeV) it gives the
order-of-magnitude estimate

\[|\nu_{\rm QFT}| \sim (M_{\rm SM}/M_{\rm Pl})^2 \sim 10^{-32}\]

which is below the precision of the cosmological comparisons below.
Assigning this term to a gravitational source requires the source
reading and split specified here.

The boundary count alone labels modes by spatial location. For an
assumed two-dimensional linear dispersion, \(dN/d\omega\propto\omega\),
the fraction at horizon frequencies is of order \((H\epsilon/c)^2\) at
one nat per mode. A reported surface Green-function calculation for the
free linear half-space construction gives a projected \(\omega^4\)
measure (§7.4). The spectral measure of the interacting boundary and its
weighting in the source remain open.

\textbf{Remark (source attribution for the running).} If the emergent
curvature response feeds back into the source,
\(\nu=\nu_{\rm QFT}+\nu_{\rm sub}\) on the stated split. Applying {[}34,
49{]} to the light Standard Model spectrum and the stated curvature
couplings \(|\xi-1/6|\leq1/6\) gives estimates
\(\nu_{\rm QFT}\approx-(4\text{–}8)\times10^{-33}\). Heavy states at
\(M_X\) in the emergent field theory could dominate the mass-weighted
sum and change the sign. If only the pre-trace-out geometry sources
gravity, the emergent running contributes to the effective description
and \(\nu=\nu_{\rm sub}\) in the source. With the indicated local
spectral weighting both readings give deviations from \(w=-1\) below the
precision considered here. A detectable running would test this
near-null model and its attribution assumptions. The free spectral
result gives neither the interacting source nor a unique inference of
nonlocality or Lorentz violation.

\textbf{Observational status.} The reported \(2.8\)--\(4.2\sigma\)
preference in DESI DR2 depends on the dataset combination and compares
\(w_0w_a\)CDM with \(\Lambda\)CDM {[}12, 13{]}. The separate Type I
variant analysis tests vacuum--dark-matter exchange models {[}14{]}.
Testing the Type II reading requires a fit in that model class.

Bertini, Novaes, von Marttens, and Shapiro {[}15{]} fit a separate
renormalization-group-corrected \(\Lambda\)CDM model with the Type II
running structure at leading order in \(\nu\), with both \(G(\mu)\) and
\(\Lambda(\mu)\) running and matter conserved, against Planck 2018 +
DESI DR2 + DES-Y5, reporting

\[\nu^{\rm Bertini} = -(2.5 \pm 1.3) \times 10^{-4}\]

Zero is within about two quoted standard deviations. The earlier
Gómez-Valent and Solà fit {[}16{]} used a different dataset and
specification and obtained \(\nu=(1.34\pm0.38)\times10^{-3}\). The fit
of {[}15{]} is compatible with both a near-null Type II reading and
\(\Lambda\)CDM.

\textbf{Cosmological neutrino mass bound.} DESI DR2 BAO plus CMB gives
\(\Sigma m_\nu < 0.0642\) eV (95\% confidence) in baseline
\(\Lambda\)CDM with degenerate neutrino masses {[}17{]}. Including
oscillation information gives \(m_l < 0.023\) eV and a preference for
normal ordering. The approximate normal-ordering minimum of \(0.059\) eV
uses the measured splittings and a negligible lightest mass. It is
compatible with the baseline upper bound. The reported tension with the
oscillation lower limit ranges from \(2.7\) to \(4.0\sigma\) by
statistic. The weaker upper limit, \(0.163\) eV, comes from the
\(w_0w_a\)CDM fit in {[}17{]}. Bertini et al.~{[}15{]} provide a
separate constraint on \(\nu\). A joint neutrino-mass fit in the OI Type
II model remains open.

\textbf{Solar neutrino comparison.} For normal ordering, JUNO's 59.1-day
measurement gives \(\sin^2\theta_{12}=0.3092\pm0.0087\) and
\(\Delta m_{21}^2=(7.50\pm0.12)\times10^{-5}\) eV\(^2\) {[}18{]}. The
global fit including JUNO and updated SNO+ gives \(0.3085\pm0.0073\) and
\((7.48\pm0.10)\times10^{-5}\) eV\(^2\) {[}19{]}. The leading flavor
value \(1/3-1/(4\pi^2)\approx0.308003\) lies \(0.14\sigma\) and
\(0.07\sigma\) from these central values, respectively, using
overlapping data (Juno §§3.5, 4.1). The comparison is retrodictive.

\textbf{Tests of the specified late-time model.} With the near-null
coefficient and the source assumptions above, the Type II model
considered has \(w(z)\approx-1\). Future measurements can constrain a
fit of that model. A null running is compatible with it and with
\(\Lambda\)CDM. Detectable running at the \(10^{-4}\) scale or a
confirmed crossing outside its predicted variation would challenge the
stated near-null implementation. Establishing the dynamics or
attributing the signal to a particular boundary spectrum requires
additional evidence.

\hypertarget{the-dark-sector-corollary}{%
\subsubsection{7.2 The dark-sector
corollary}\label{the-dark-sector-corollary}}

The calculation below compares the boundary thermal energy with the
supplied flat Friedmann density. Its use as a physical source is a
separate question.

\textbf{Boundary-density normalization.} Under the de Sitter first-law
and counting assumptions, the Hubble-volume mean of the horizon thermal
energy equals the flat critical density in three spatial dimensions.

\emph{Calculation.}

\emph{(i) Total effective density.} The horizon carries
\(A/\epsilon^2 = S_{\text{dS}}/\sigma\) modes, of total entropy
\(S_{\text{dS}} = \sigma A/\epsilon^2\) (§5). The classical temperature
is \(T_{\text{cl}} = c^3 \epsilon^2 H / (8\pi G k_B \sigma)\) (§3.1),
computed entirely within the classical substratum with no reference to
\(\hbar\) or the emergent quantum description. The total thermal energy
is \(E_s = S_{\text{dS}} \cdot k_B T_{\text{cl}}\). Its mean over the
Hubble volume is obtained by dividing by \(V_H = 4\pi c^3/(3H^3)\):

\[\rho_s = \frac{E_s}{V_H} = \frac{\sigma A}{\epsilon^2} \cdot \frac{c^3 \epsilon^2 H}{8\pi G \sigma} \cdot \frac{3H^3}{4\pi c^3} = \frac{3c^2 H^2}{8\pi G} = \rho_{\text{crit}}\]

The factors \(\epsilon^2\) and \(\sigma\) cancel between the count and
temperature. The identity uses the horizon first-law normalization
\(dE=(c^2\kappa/8\pi G)\,dA\), the area-law entropy and the supplied
Friedmann scale.

The normalization relates a boundary thermal energy to a volume average.

\emph{Classical inputs.} The boundary entropy, temperature and Hubble
volume enter before the designated trace-out. The zero-point estimate
\(\rho_{\rm QFT}\sim M_{\rm Pl}^4\) belongs to the emergent description
(§6.1). The comparison in §6.3 uses their scales. The field equations
and offset response require the G3 coupling.

\emph{Physical locus.} The counted modes are at the boundary. Even in
Newtonian gravity a spherical shell has zero interior force, while a
uniform mass density produces an interior field proportional to radius.
The matching mean density leaves the interior field undetermined.

\emph{Physical source.} A bulk stress or flux associated with
deformations of the boundary requires the native coupling at G3.

\emph{General \(d\).} Using the same first-law normalization
\(dE/dA=c^2\kappa/(8\pi G_d)\) with boundary entropy
\(S=\sigma A/\epsilon^{d-1}\), the corresponding thermal gap equation
has \(\hbar=c^3\epsilon^{d-1}/(4G_d\sigma)\) on its stated calibration
conditions. For the classical density comparison, set \(c=1\), take
\(R=1/H\), \(\kappa=H\) and \(V_d=AR/d\). Then
\(S\,k_BT_{\rm cl}/V_d=dH^2/(8\pi G_d)\). The supplied flat Friedmann
density is \(d(d-1)H^2/(16\pi G_d)\). Their full coefficients give

\[\frac{\rho_s}{\rho_{\rm crit}} = \frac{2}{d-1}\]

which equals unity only at \(d=3\). Its use as a consistency filter on
physical sourcing requires the identification of the boundary energy
with the relevant gravitational source {[}SM §3.2{]}.

Within A4's assumed periodic-lattice class and the continuum conditions
of SM §3.1--§3.2, propagating gravity gives \(d\geq3\) and stable matter
gives \(d\leq3\), already selecting \(3\). Boundary-entropy concordance,
with its stated horizon assumptions, is redundant for the selection once
these two hold but not logically dependent on them; it gives a separate
conditional comparison with \(d\geq4\) (SM §3.2).

\textbf{Dark-sector interpretation.}

The boundary variables are integrated out of the designated reduced
description. Their possible effect on the visible record is encoded in
the resulting law.

A gravitational budget associated with those variables would require the
proposed coupling between boundary entropy and the bulk metric. It would
have to reproduce the effect of actual matter configurations and their
stresses (§6.2).

On that proposed identification, using the baryon fraction \(\Omega_B\)
leaves \(1-\Omega_B\approx0.95\) for the nonbaryonic contribution. The
fraction uses the supplied baryon budget.

The division of that contribution into matter-like and vacuum-like parts
remains open.

\textbf{Observational status of particle dark matter searches.} The LZ
417-day low-mass analysis {[}21{]} reports no signal from weakly
interacting massive particles. It sets limits on particle--nucleon
scattering cross sections for its tested interactions in the
\(3\)--\(9\,\mathrm{GeV}/c^2\) low-mass comparison. A detected particle
species that accounts for the structured dark sector would contradict
the proposed boundary-only interpretation.

\textbf{Memory effects and record persistence.} A redistribution of the
boundary entropy in response to matter is a proposed mechanism. A
relaxation time of order \(H^{-1}\), as in the physical C2 input of
§2.2, would give persistence on sub-Hubble timescales. The
matter-to-boundary map and the readback condition C4 remain separate
requirements.

Section 7.3 treats the conditional displacement law and the open
relation between the local and large-scale regimes.

\textbf{Falsifiability.} The corollary would be falsified by detection
of dark matter particles accounting for the full \textasciitilde27\%
structured dark sector, or by evidence that the dark sector's
gravitational effects are unrelated to the boundary geometry.

\hypertarget{dark-matter-from-entropy-displacement}{%
\subsubsection{7.3 Dark matter from entropy
displacement}\label{dark-matter-from-entropy-displacement}}

The boundary-density normalization of §7.2 motivates a proposal for the
structured component. The displacement calculation uses the additional
horizon and matter inputs below.

\textbf{Inherited conditions.} The chain below rests on the de Sitter
temperature of §3.2, which is conditional: \(T_{\rm dS}\) requires
H-count (the hidden sector's energy-dependent entropy is its boundary
count), H-slope (the infrared detailed-balance slope), H-frame (the
preferred foliation is the comoving one), the H-Hawking conditions,
local boost structure at the horizon and the symmetric detailed-balance
rate ratio (§8.5), and the action-scale calibration is stated on an
exact de Sitter horizon. Everything in §7.3 therefore carries those
conditions. A second assumption is local to this section: writing
\(Q=M_Bc^2\) takes the entire rest energy of a static configuration to
be delivered as heat to the boundary. That is a statement about how the
emergent sector couples to the boundary --- the same coupling map G3
names as open --- and does not follow from (C1)'s requirement of nonzero
cross-boundary interaction alone. A third input is imported: the
dimensional factor converting the Hubble acceleration \(cH\) into the
critical acceleration is Verlinde's, not derived here (Step 2).

\textbf{Step 1: Entropy displacement (conditional on the above).}
Baryonic matter \(M_B\) is coupled to the boundary entropy through
\(H_{VB}\) (condition C1). The boundary is a thermal reservoir at
\(T_{\text{dS}} = \hbar H / (2\pi k_B)\) (§3.2, conditional). For a
static matter configuration in gravitational equilibrium --- a
quasi-static process --- the Clausius relation gives the entropy
displacement:

\[\Delta S = \frac{M_B c^2}{k_B T_{\text{dS}}} = \frac{2\pi M_B c^2}{\hbar H}\]

This is the Clausius relation applied to the boundary reservoir: the
boundary's entropy reduction is \(Q/T\) for a reversible process, where
\(Q = M_B c^2\) and \(T = T_{\text{dS}}\). This step uses the reversible
static heat input.

\textbf{Step 2: The critical acceleration (conditional; imported).} The
displaced entropy redistributes over the volume between the matter and
the horizon. The proposed radial coupling uses the acceleration
conversion in the Verlinde branch. The conversion is dimensional: in
\(d\) spacetime dimensions, the factor relating the natural Hubble
acceleration scale \(cH\) to the empirical MOND acceleration is
\((d-3)/[(d-2)(d-1)]\). This derivation is due to Verlinde {[}20{]} (eq.
1.7), and follows from the interplay between volume-law entropy
distributed over \((d-1)\) spatial dimensions and area-law entropy
scaling with \((d-2)\) surface dimensions, together with the standard
surface-density-to-acceleration relation via Gauss's law. In our
universe (\(d = 4\)):

\[a_0 = \frac{(d-3)}{(d-2)(d-1)} \cdot cH = \frac{1}{2 \cdot 3} \cdot cH = \frac{cH}{6} = 1.09 \times 10^{-10} \; \text{m/s}^2 \quad (H_0 = 67.4\ \text{km s}^{-1}\text{Mpc}^{-1})\]

This is the MOND critical acceleration with no fitted parameter:
\(cH/6\) runs from \(1.09 \times 10^{-10}\) m/s² at the Planck value of
\(H_0\) to \(1.18 \times 10^{-10}\) m/s² at \(H_0 = 73\), against the
radial-acceleration-relation fit
\(g_\dagger = (1.20 \pm 0.02_{\rm rand} \pm 0.24_{\rm sys}) \times 10^{-10}\)
m/s² {[}22{]}, a ratio of \(0.91\)--\(0.98\) that lies inside the
stellar mass-to-light systematic. The factor 1/6 has dimensional origin
--- it is fixed by \(d = 4\) together with the volume-vs-area entropy
interplay --- and is taken over from Verlinde's analysis rather than
rederived from the substratum.

The OI framework's relation to Verlinde's derivation is that the
volume-law entropy mechanism is placed on the framework's named
conditions. In Verlinde's framework, the de Sitter states'
non-thermalization (``memory effect'') is hypothesized; in OI it is the
persistence condition C2 (\(\tau_B \gg \tau_S\)) used as a physical
input to the horizon application (§2.2). Whether this meets the
consistency objection of {[}37{]} depends on the uncomputed relaxation
kernel (§7.3); the published Verlinde--MOND consistency debate is
relevant to the elasticity assumption in {[}20{]} and, through that
kernel, to the present account.

\textbf{Step 3: The deep-MOND regime.} For baryonic acceleration
\(\ll a_0\), the assumed displacement law supplies the dominant radial
acceleration, giving \(g_{\text{total}} = \sqrt{g_B \cdot a_0}\) and
rotation velocity:

\[v^2 = r \cdot g_{\text{total}} = \sqrt{GM_B \cdot cH/6} = \text{constant} \qquad \Longrightarrow \qquad v^4 = GM_B \cdot cH/6\]

This is the deep-MOND baryonic Tully-Fisher relation on the stated
branch. Its normalization is compared with measured asymptotic rotation
velocities and baryonic masses.

\textbf{Step 4: Radial interpolation.} For a static radial law, an
entropy displacement potential is proposed with \(g_D=-\nabla\Psi_D\).
The following restrictions apply to the radial comparison:

\begin{enumerate}
\def\labelenumi{(\roman{enumi})}
\item
  \emph{No free parameters.} \(g_D\) depends only on \(g_B\) and
  \(a_0 = cH/6\) --- no halo-specific parameters, no fitting constants.
\item
  \emph{Conservative.} \(\nabla \times g_D = 0\), since \(g_D\) derives
  from the entropy potential.
\item
  \emph{Monotonic.} \(\partial g_D / \partial g_B > 0\) --- more
  baryonic mass produces more entropy displacement, hence more dark
  matter effect.
\item
  \emph{Limits.} \(g_D/g_B\to0\) for \(g_B/a_0\to\infty\) (Newtonian),
  and \(g_D\sim\sqrt{g_Ba_0}\) for \(g_B/a_0\to0\) (deep-MOND regime).
\item
  \emph{Analytic.} The entropy redistribution is a smooth function of
  the coupling structure --- no discontinuities or phase transitions in
  the interpolation.
\end{enumerate}

The endpoint and monotonicity conditions admit a range of smooth
interpolation functions. The simple interpolation is one radial example:

\[g_D = \frac{g_B}{2}\left(\sqrt{1 + \frac{4a_0}{g_B}} - 1\right)\]

It gives \(g_{\rm total}=g_B\nu(g_B/a_0)\) with
\(\nu(y)=(1+\sqrt{1+4/y})/2\). Rationalizing \(g_D\) gives
\(2a_0/(1+\sqrt{1+4a_0/g_B})\), so the absolute correction approaches
\(a_0\) and the fractional correction approaches zero at large \(g_B\).
The radial-acceleration data in {[}22{]} constrain the acceleration-law
comparison.

\textbf{Epistemic status.} Steps 1--3 are Tier 2 and conditional: the
displacement formula follows from the Clausius relation applied to the
boundary reservoir (§3.2) and therefore inherits the conditions on
\(T_{\rm dS}\) together with the local assumption \(Q = M_Bc^2\); the
dimensional factor is imported (Step 2); given these, \(a_0 = cH/6\) and
\(v^4 = GM_B \cdot cH/6\) contain no fitted parameter. The interpolation
constraints (i)--(v) are Tier 2 on the same conditions; the transition
steepness is Tier 3 (solution-specific).

\textbf{Redshift evolution.} On extending the conditional radial law to
the supplied flat cosmological background, \(a_0(z) = cH(z)/6\) and
\(H(z) = H_0\sqrt{\Omega_m(1+z)^3 + \Omega_\Lambda}\), so the critical
acceleration increases with redshift: \(a_0(z=1) = 1.79\,a_0(0)\),
\(a_0(z=2) = 3.03\,a_0(0)\), \(a_0(z=3) = 4.57\,a_0(0)\). The MOND
crossover radius \(r_M = \sqrt{6GM_B/(cH)}\) correspondingly shrinks ---
for a Milky-Way-like baryon budget
(\(M_B \approx 6 \times 10^{10}\,M_\odot\) at the stellar-plus-cold-gas
level), \(r_M = 8.8\) kpc at \(z = 0\) but \(5.0\) kpc at \(z = 2\) and
\(3.0\) kpc at \(z = 5\). Galaxies with effective radii comparable to or
exceeding \(r_M\) show flat rotation curves (the familiar local
behavior); galaxies with \(r_{\text{eff}} < r_M\) decline,
Keplerian-like, inside \(r_M\) and approach flat beyond it. At fixed
radius the baryonic share of the acceleration is
\(g_B/g_{\text{total}} = 1/\nu(g_B/a_0)\), and \(\nu\) is decreasing, so
the share falls as \(a_0\) rises. High-redshift galaxies are therefore
less baryon-dominated than local galaxies with the same baryon
distribution, and their outer curves flatten at smaller radii: for the
\(10^{10.8}\,M_\odot\) disk of test (i) below, \(g_B/g_{\text{total}}\)
at \(2r_{\text{eff}}\) is \(0.55\), \(0.46\) and \(0.38\) at \(z = 0\),
\(1\) and \(2\) (enclosed mass in the spherical approximation).

Two comparisons run opposite to this account. Genzel et al.~{[}23, 24{]}
report stacked rotation curves for massive star-forming galaxies at
\(z = 0.9\text{–}2.4\) with outer velocities declining at \(> 3\sigma\)
significance relative to local spirals, and read them as baryon
dominance together with pressure support; at a fixed baryonic profile
the law predicts flatter outer curves at those redshifts. In the Milky
Way, Jiao et al.~{[}25{]} detect a Keplerian decline from \(\sim 19\) to
\(\sim 26.5\) kpc using Gaia DR3 kinematics, with the flat rotation
curve hypothesis rejected at \(3\sigma\). That range lies beyond the
crossover for either baryon budget --- \(r_M = 8.8\) kpc for
\(M_B = 6 \times 10^{10}\,M_\odot\), and \(16.0\) kpc for
\(M_B \sim 2 \times 10^{11}\,M_\odot\) including the hot circumgalactic
gas halo --- and there the law gives
\(v(26.5\,\text{kpc})/v(19.5\,\text{kpc}) = 0.97\) and \(0.95\) for a
point mass, against \(0.86\) for a Keplerian decline. At fixed radius
the circular velocity increases with the enclosed baryonic mass, and the
point-mass ratio decreases with the mass, so in the spherical
approximation the ratio is at least \(0.95\) for any baryon distribution
with \(M_B \leq 2 \times 10^{11}\,M_\odot\). Both observations are
therefore in tension with the law. A quantitative comparison needs the
Milky Way baryon model and the circular velocities of Jiao et al.'s
Jeans analysis and, for the high-redshift sample, the galaxies' own
baryon profiles and the pressure-support correction, since real galaxies
at high redshift differ from local ones in both while the law's
prediction above is stated at a fixed profile; it is not made here.

\textbf{Tully-Fisher evolution.} McGaugh et al.~(2024), analyzing Nestor
Shachar et al.~(2023) JWST kinematics, present the \emph{baryonic}
Tully-Fisher relation out to \(z \sim 2.5\) and find no clear indication
of evolution, noting unresolved differences between samples in the
definition of baryonic mass and in the velocity measured. The baryonic
relation is the direct test of
\(v_{\text{flat}}^4 = GM_B \cdot cH(z)/6\): at fixed \(v_{\text{flat}}\)
the law predicts a normalization shift
\(\Delta\log M_B = \log(H_0/H(z)) = -0.48\) dex at \(z = 2\). The
reported absence of evolution is therefore in tension with the law. The
tension is not yet statistically quantified: that needs baryonic masses
with reliable ALMA gas masses, and velocities measured as the law's
\(v_{\text{flat}}\). Where only stellar mass is available, the stellar
relation shifts at least as far as the baryonic one, since
\(M_* = (1-f_{\text{gas}})\,M_B \leq M_B\); the large gas fractions at
high redshift (\(f_{\text{gas}} \sim 50\text{–}70\%\) at \(z \sim 2\);
Tacconi et al.~2018) enlarge the offset in \(M_*\), to
\(\Delta\log M_* = \log[(1-f_{\text{gas}})\,H_0/H(z)] = -0.78\) dex at
\(z = 2\) for \(f_{\text{gas}} = 0.5\) against a gas-poor local
relation.

\textbf{Observational comparisons.} The following tests use the stated
radial law and its inputs.

\begin{enumerate}
\def\labelenumi{(\roman{enumi})}
\item
  \emph{Rotation curve evolution.} The outer rotation curve evolves with
  \(z\) through \(a_0(z) = cH(z)/6\). For a \(10^{10.8}\,M_\odot\)
  exponential disk with \(r_{\text{eff}} = 5\) kpc, the outer slope
  \(\Delta v / v\) between \(r_{\text{turn}}\) and \(2r_{\text{turn}}\)
  is \(-0.08\) at \(z = 0\), evolving to \(-0.06\) at \(z = 2\) and
  \(-0.05\) at \(z = 5\) --- flattening as the disk increasingly sits in
  the deep-MOND regime, while the asymptotic
  \(v_{\text{flat}}(z) = (GM_B \cdot cH(z)/6)^{1/4}\) shifts upward by
  \(E(z)^{1/4}\) (\(+32\%\) at \(z = 2\)). Comparison with a
  particle-halo prediction requires a specified galaxy and halo
  evolution model.
\item
  \emph{Tully-Fisher evolution.} The baryonic Tully-Fisher relation
  evolves: \(v_{\text{flat}}^4 = GM_B \cdot cH(z)/6\). At fixed baryonic
  mass, \(v_{\text{flat}} \propto H(z)^{1/4}\): 7\% higher at
  \(z = 0.5\), 16\% higher at \(z = 1\), 32\% higher at \(z = 2\).
  Equivalently, at fixed \(v_{\text{flat}}\), the inferred baryonic mass
  at \(z = 2\) is \(H_0/H(z=2) = 0.33\times\) the local value while the
  dynamical mass at fixed \(v_{\text{flat}}\) and radius is unchanged,
  so the mass discrepancy \(M_{\text{dyn}}/M_B\) is \(3.0\times\) the
  local value --- galaxies appear to need more dark matter at high
  \(z\). Testable with JWST and ALMA kinematics combined with
  photometric baryonic mass estimates.
\item
  \emph{Radial acceleration comparison.} At fixed redshift the specified
  radial function gives a single \(g_{\rm total}\) for a given \(g_B\).
  The scatter of about \(0.13\) dex in {[}22{]} is measured for the
  radial acceleration relation. Its interpretation uses the
  uncertainties in the baryonic profiles, kinematics and other inputs to
  \(g_B\) and \(g_{\rm total}\).
\end{enumerate}

\textbf{Tests of the proposed radial law.} A species accounting for the
structured dark sector would contradict a boundary-only source. Measured
outer velocities and baryonic profiles can test the conditional redshift
trend, with pressure support included where applicable. Testing the
Tully-Fisher normalization requires measured \(v_{\rm flat}\) and
baryonic mass. The radial acceleration scatter is a separate observable.

\textbf{Cluster-scale status.} Galaxy clusters present the most
challenging regime for MOND-like theories. For the Coma cluster
(\(M_B \approx 1.4 \times 10^{14}\,M_\odot\), \(R_{200} \approx 1.9\)
Mpc), the acceleration at the virial radius is \(g_B/a_0 \approx 0.05\)
--- deep in the MOND regime. With the simple interpolation function
above, the predicted velocity is
\(v = \sqrt{g_B \cdot \nu(g_B/a_0) \cdot R_{200}} \approx 1260\) km/s;
the observed velocity dispersion implies \(\sim 1270\) km/s. For other
rich clusters (Perseus, A2029, A1689), the simple interpolation reduces
the standard MOND mass shortfall from a factor \(\sim 2\) to
\(\sim 1.0\)--\(1.5\); the residual is \textbf{not} closable by
undetected cluster baryons. Modern censuses find massive clusters
(\(M_{500} \gtrsim 2\times10^{14}\,M_\odot\)) baryon-complete to within
\(\sim 7\%\) of the cosmic fraction inside \(r_{500}\) (Planck
cosmology), and consistent with the cosmic value at the virial radius;
the late-time missing baryons localized by FRB dispersion measures
(Connor et al.~2025) reside in the diffuse IGM and cosmic filaments ---
\emph{between} clusters, not inside them. The quoted normal-ordering
minimum \(\Sigma m_\nu\approx0.059\) eV is a lower bound. The upper
limits in §7.1 use their respective cosmological models. Assessing a
neutrino contribution to the rich-cluster residual also needs the
cluster mass and dynamics model. The residual remains open. The proposed
connection between the local interpolation and a matter-like large-scale
source remains open: a single field cannot be MOND-like in cluster cores
and CDM-like in linear growth without a specified crossover, and
deriving that crossover from the non-local relaxation kernel is deferred
to the substratum-level analysis (the same uncomputed kernel as §7.3).
The deep-MOND limit \(v^4 = GM_B \cdot a_0\) (Tier 2) is identical for
all interpolation functions; only the intermediate-acceleration behavior
(Tier 3) differs. The radial-acceleration data have a separately
measured scatter of about \(0.13\) dex.

\textbf{Colliding clusters (Bullet Cluster).} The observed lensing peaks
lie near the galaxy concentrations, offset from the X-ray gas. A
proposed account assumes a pre-collision boundary displacement centered
on those concentrations and a relaxation time
\(\tau_{\rm relax}\sim H^{-1}\approx14\) Gyr at low redshift. A crossing
time \(720\,\text{kpc}/(4700\,\text{km/s})\approx0.15\) Gyr then gives a
ratio near \(0.01\). A displacement that relaxes on the assumed time
would change little during the collision. The pre-collision pattern, its
relaxation law and the native map from matter and displacement to the
measured lensing remain to be established.

\textbf{Acoustic peaks.} The proposed large-scale mechanism uses
boundary heat transfer averaged over an acoustic cycle. A purely
oscillatory heat transfer with zero mean would give no net displacement,
leaving a possible response to a supplied growing component.
Establishing this mechanism requires the coupled matter, metric and
radiation evolution near recombination.

A reported CLASS comparison from April 2026 used standard cold dark
matter in both runs and encoded the dark-energy variation with its CLP
fluid module at \(\nu=10^{-4}\). The reported shifts in the first four
temperature peak amplitudes were \(\lesssim0.02\%\), with a median
\(0.006\%\) and maximum \(0.08\%\) residual in \(D_\ell^{TT}\) over
\(2\leq\ell\leq2500\). The calculation concerns a dark-energy proxy with
the matter content supplied. Deriving a spectrum from the proposed
displacement would require the native source and its scale dependence.

\textbf{Illustrative relaxation equation.} Taking
\(\tau_{\rm relax}=1/H(z)\) and supplying a matter-era growing mode
\(\delta_m\propto a\), the toy equation
\(d\delta_E/dt=(\delta_m-\delta_E)/\tau_{\rm relax}\) becomes
\(d\delta_E/da=(\delta_m-\delta_E)/a\). It has
\(\delta_E=\delta_m/2+C/a\) with \(C\) fixed by the initial condition.
Thus it tracks a supplied matter-era growing mode up to a decaying
homogeneous term. With \(\Omega_m\simeq0.3\) in the expansion law above,
the fixed-timescale response rate relative to the local rate at
\(z=100\) would be \(H_0/H(100)\simeq1/560\). In the collision example
\(t_{\rm cross}/\tau_{\rm relax}\simeq0.15/12\simeq0.013\) if the
horizon time is \(12\) Gyr. These time comparisons leave the
perturbation source, the \(k\) dependence and the recombination-era
Einstein--Boltzmann evolution open.

\hypertarget{summary-of-gr-side-predictions}{%
\subsubsection{7.4 Summary of GR-side
predictions}\label{summary-of-gr-side-predictions}}

The table summarizes the gravitational comparisons. Each formula uses
the conditions of its derivation.

\begin{longtable}[]{@{}
  >{\raggedright\arraybackslash}p{(\columnwidth - 6\tabcolsep) * \real{0.2500}}
  >{\raggedright\arraybackslash}p{(\columnwidth - 6\tabcolsep) * \real{0.2500}}
  >{\raggedright\arraybackslash}p{(\columnwidth - 6\tabcolsep) * \real{0.2500}}
  >{\raggedright\arraybackslash}p{(\columnwidth - 6\tabcolsep) * \real{0.2500}}@{}}
\toprule\noalign{}
\begin{minipage}[b]{\linewidth}\raggedright
Prediction
\end{minipage} & \begin{minipage}[b]{\linewidth}\raggedright
Formula / value
\end{minipage} & \begin{minipage}[b]{\linewidth}\raggedright
Observed / status
\end{minipage} & \begin{minipage}[b]{\linewidth}\raggedright
Match
\end{minipage} \\
\midrule\noalign{}
\endhead
\bottomrule\noalign{}
\endlastfoot
Type II structural reading & Conserved matter and running \(G\) in the
stated implementation. Near-null estimate conditional on the source
split (§7.1) & Bertini et al.~\(\nu=-(2.5\pm1.3)\times10^{-4}\) &
compatible with near-null and \(\Lambda\)CDM \\
Dark sector fraction & \(1-\Omega_B\) if the proposed boundary source
accounts for the nonbaryonic budget &
\(\Omega_{\rm CDM}+\Omega_\Lambda\approx0.95\) & uses the supplied
baryon fraction \\
MOND acceleration scale &
\(a_0 = cH/6 = 1.09\)--\(1.18 \times 10^{-10}\) m/s²
(\(H_0 = 67.4\)--\(73\)) &
\(g_\dagger = (1.20 \pm 0.02 \pm 0.24) \times 10^{-10}\) m/s² & within
systematic; factor imported \\
Baryonic Tully-Fisher & \(v_{\rm flat}^4=GM_B cH/6\) under the horizon
and heat-input conditions & asymptotic velocities and baryonic masses &
normalization comparison open \\
Rotation curve crossover & \(r_M = \sqrt{6 G M_B/(cH)}\); nearly flat
beyond & Milky Way: \(9\)--\(16\) kpc depending on the baryon budget
(\(6\times10^{10}\)--\(2\times10^{11}\,M_\odot\)); Jiao et al.~2023
Keplerian decline at \(19\)--\(26.5\) kpc & in tension \\
High-\(z\) rotation curves & Less baryon-dominated, flatter with \(z\)
at a fixed baryon profile & Genzel et al.~\(z = 0.9\)--\(2.4\),
declining at \(> 3\sigma\) & in tension \\
\end{longtable}

The conditional horizon calibration fixes
\(\epsilon/\sqrt{\sigma}=2l_p\) through the action-scale relation and
the Planck-length definition. The Bekenstein-Hawking \(1/4\) coefficient
uses those same horizon, frame, counting and slope conditions (§5).
GW250114 tests the separate classical area theorem. The vacuum-energy
comparison uses the supplied Friedmann scale (§6).

The local \(H^2-H_0^2\) form is common to smooth densities near \(H_0\).
The Type II reading uses conserved matter and running \(G\) in the
considered model. The source identification and interacting spectral
weighting remain open. A near-null value under the assumptions of §7.1
and \(\Lambda\)CDM both admit a null running measurement.

\emph{Boundary spectral measures.} Under the local two-dimensional
dispersion used in §7.1, \(dN/d\omega\propto\omega\). A log-uniform
intrinsic measure would weight the infrared differently. A
stretched-horizon Clausius calculation in the de Sitter static patch
gives a log-uniform redshift band for a local spectrum, with the
contribution to the running suppressed by the near-wall weighting of the
density derivative. The reported free linear half-space calculation
gives \(dN/d\omega\propto\omega^4\), with the \(+1\) and \(+2\) cases as
estimator controls. These results concern the stated geometric and
free-linear settings. The interacting kernel and attribution of physical
running remain open. Preservation of locality is also an issue in the
Lorentz-sector calculation {[}SM §3.1{]}. A cross-sector implication
requires a map from the interacting boundary theory to those
observables.

\begin{center}\rule{0.5\linewidth}{0.5pt}\end{center}

\hypertarget{discussion}{%
\subsection{8. Discussion}\label{discussion}}

The foundational emergent QM result is treated as established in
{[}Main{]} and not re-derived here. The cubic-lattice realization of the
substratum --- which determines the Standard Model gauge structure via
the wave equation --- is independent of the present results and is
developed in {[}SM{]}. The reconstruction theorem, the substratum gauge
group, and the synthesis claim that quantum mechanics and general
relativity descend as projections of the same \((S, \varphi)\), with a
third projection bearing on the arrow of time, are developed in
{[}Substratum{]}.

\hypertarget{scope}{%
\subsubsection{8.1 Scope}\label{scope}}

The characterization theorem's equivalence is (ii) \(\iff\) (iii), with
(iii) carrying (C1), (C3) and (C4) --- both directions proved in the
foundational paper --- with the quantum representation following for the
source's finite-configuration class; (C2) lies outside the equivalence,
the slow-bath form of its necessity for physical hidden sectors
conditional on the hidden sector's mixing hypothesis (ETH is its
physical motivation; see {[}Main, §3.4{]}), its persistence form holding
without it. The C1, C3, and C4 necessity directions are unconditional
--- C4's is immediate from observed non-Markovianity. For the horizon
application the mixing timescale is an input, the constraint
correlations motivate C1, and C4 is not presently discharged at the
cosmological cut (§2.2). The theorem does not identify which physical
systems satisfy the conditions (empirical) nor derive \(\hbar\) for
specific systems (requires Part II). The physical memory-bearing
application requires the realization of its conditions on the chosen
causal partition.

\hypertarget{the-status-of-the-discreteness-scale}{%
\subsubsection{8.2 The status of the discreteness
scale}\label{the-status-of-the-discreteness-scale}}

Finite distinguishability supplies the finite visible alphabet of Lemma
1. A finite horizon area yields finitely many boundary modes given the
separate count with a positive area \(\epsilon^2\) per mode. The
physical identification of \(\epsilon\) with a lattice counting scale
uses the assumptions of §3.1. The calibration fixes
\(\epsilon/\sqrt{\sigma}=2l_p\) (§4).

\hypertarget{vacuum-energy-relative-effects-and-the-higgs-potential}{%
\subsubsection{8.3 Vacuum energy: relative effects and the Higgs
potential}\label{vacuum-energy-relative-effects-and-the-higgs-potential}}

\textbf{Relative effects.} Casimir forces and Lamb shifts compare
configurations within the emergent field theory. A gravitational
comparison of two preparations uses the stress of each complete system,
including supports and reservoirs. The proposed OI coupling to these
sources remains open (§6.2).

\textbf{The Higgs potential.} The shift in the effective Higgs potential
at electroweak breaking is conventionally of order
\((200\,\mathrm{GeV})^4\). Distinguishing a common additive offset from
a change in the physical preparation is essential to the gravitational
comparison.

\emph{Relative energies.} The test-source construction of §6.2 retains
the total energy difference between two complete preparations. Binding
and Casimir comparisons must include their material and work
contributions. The proposed native source would have to reproduce the
appropriate total differences.

\emph{Common offsets.} A constant term shared by two complete conserved
stress tensors cancels in their difference on one metric when the
physical reference is fixed. The value of that reference and the
response to a change in the vacuum state require further dynamics.

\emph{The reconstruction-gauge form of this claim.} An additive energy
shift preserves every visible transition probability exactly (§3.3): the
uniform offset of \(\hat H\) is a gauge parameter of the
marginalization, present in no visible transition statistics at any
precision. The embedded observer's quantum data therefore contain no
fact about the absolute zero-point offset. Whether the offset sources
curvature is not settled by this invariance; it is the coupling question
that §6.2 leaves open at Level G3. The invariance covers the uniform
offset. Comparing binding effects, Lamb shifts or masses requires the
native map from the corresponding preparations to stress or flux.

Zero-point, Higgs and condensate contributions that enter as the same
constant in both preparations are subject to this conditional
subtraction. The gravitational response to energy released in
electroweak or QCD transitions remains a question for the coupling.

\textbf{QCD binding energy and the quantum vacuum.} A comparison of
bound configurations uses their complete rest energies on the
common-metric weak-field branch (§6.2). To apply it to the proton, the
native construction must reproduce the stress associated with its QCD
dynamics and the rest-energy differences of the physical preparations.
The boundary entropy count does not supply this map.

\hypertarget{assumptions-limitations-and-falsifiability}{%
\subsubsection{8.4 Assumptions, limitations, and
falsifiability}\label{assumptions-limitations-and-falsifiability}}

\textbf{The theorem} requires Lemma 1 (finite distinguishability in Part
I, with the physical horizon count separately supplied in Part II), the
stochastic-quantum bridge ({[}Main, §3.1--§3.2{]}; two routes reach its
transition-statistics layer, neither reaching the operational instrument
algebra), and genericity conditions (the phase-locking hypotheses
(G1)--(G4) of {[}Main §3.1{]}: non-degenerate spectrum and gaps,
non-vanishing overlaps, triangle genericity), which hold for all but a
measure-zero set. The P-indivisibility proof uses total finiteness
({[}Main §4.5{]}, posit (ii)) via the recurrence argument
(\(\varphi^N = \text{id}\)), or the near-return condition (NR) of
{[}Main §2.3{]}, assumed in its place as a hypothesis and not adopted.
For the visible transition data on accessible timescales, the
deep-sector cardinality corollary (§3.2) establishes that on
\([0, n\,\Delta t]\) only \(\mathcal{C}_V\) and the causal-cone sector
\(\mathcal{C}_{B_n}\), a finite set of sites, need be finite, while the
deep sector may be infinite given the normalized prior and dynamics
stated there, and that at leading order under H-BD, for a finite deep
sector of any size, only \(\mathcal{C}_V\) and \(\mathcal{C}_B\) enter,
these two finite by holographic entropy bounds {[}26{]} independently of
the de Sitter finite-dimensionality conjecture; recurrence-scale
P-indivisibility is proved on a finite representative. The readback
lemma ({[}Main, §2.3{]}) provides accessible conditional memory
independently of \(|\mathcal{C}_H|\); it does not by itself make the
rooted family P-indivisible on accessible windows. The physical
slow-bath application assumes a mixing time of order \(10^{17}\) s and
requires the separate accessible readback condition. The finite-law
representation and the recurrence result have the scopes stated above.

\textbf{The sharp partition and constraint structure.} Lemma 2's product
decomposition \(\Gamma = \Gamma_V \times \Gamma_H\) is the kinematic
phase space. In GR, the Hamiltonian and momentum constraints restrict
the physical phase space to a submanifold
\(\Gamma_{\text{phys}} \subset \Gamma_V \times \Gamma_H\) on which
interior and exterior data are correlated. The recurrence-cycle
P-indivisibility of {[}Main §2.3{]} carries over to a finite
representative of the constraint surface, invariant under \(\varphi\),
with the hidden prior taken on each root's fiber: \(\varphi\) restricted
to \(\Gamma_{\text{phys}}\) is a permutation of finite order \(L\), so
the rooted visible maps \(K_t\), indexed by the visible values the
surface admits, satisfy \(K_{mL} = I\) for every \(m \geq 0\), and a
stochastic \(\Lambda\) with \(K_s \Lambda = K_{mL}\) for some \(mL > s\)
would be a stochastic right inverse of \(K_s\), which the
stochastic-inverse lemma of {[}Main §2.3{]} excludes unless \(K_s\) is a
permutation. Indivisibility across the recurrence cycle therefore needs
only that some rooted visible map is not a permutation; it does not
extend to accessible windows. At the cosmological horizon the boundary
Hamiltonian and constraint structure motivate a candidate coupling.
Accessible hidden-to-visible influence remains a condition {[}Main
§1.3{]}. For systems other than the cosmological horizon, the fidelity
of the product approximation determines the fidelity of the emergent
quantum description.

\textbf{Planck-scale ordering.} The classical-prior ontology may break
down at \(l_p\), where ``geometric'' and ``quantum'' lose operational
distinction. If a deeper theory unifies both, this framework describes
the effective macroscopic regime.

\textbf{Falsifiability.} The late-time model test uses the near-null
Type II implementation with the source and spectral assumptions of §7.1.
A measured \(w(z)\) variation or running outside its stated range would
challenge that implementation. A species accounting for the structured
dark sector would contradict the boundary-only interpretation of §7.2.

The Type II comparison constrains the specified late-time model. The
general local density expansion and the boundary normalization have
their respective assumptions in §§7.1--7.2.

\textbf{A conditional mechanism-side test.} The operational distinction
between classical and quantum memory in non-Markovian processes {[}27,
28{]} classifies a two-step process as classical-memory when its second
step can be written as channels conditioned on the classical outcome of
an instrument realizing the first step. The classification is a property
of the process maps on all density matrices, so it bears on the
hidden-sector ontology only through the coherent-preparation and
instrument bridge, which is open ({[}Main, §4.5{]}). On configuration
inputs every two-step process admits a classical-memory decomposition;
on coherent inputs the classification depends on the realization rather
than on the visible law, and a swap into an initialized hidden slot
admits none. A memory-witness experiment bears on the ontology only
given that bridge and an outcome-independent rule identifying the
\emph{fundamental} non-Markovian processes --- those arising from the
cosmological trace-out rather than from engineered quantum baths --- and
it falsifies the ontology only given, in addition, a constraint derived
from the completed model on what such processes must show, such as
classical memory for every eligible realization. No such constraint is
established ({[}Main, §4.5{]}).

\textbf{Black-hole horizons under partition-relativity.} For the
external observer considered here, the merger and the infalling matter
remain inside the cosmological visible region. The black-hole horizon is
a local causal boundary within that region. The designated cosmological
trace-out provides no echo or reflection mechanism, and no echo
amplitude is derived. Standard ringdown would follow on recovery of the
relevant gravitational dynamics at G3 (§7).

\textbf{Nested and time-dependent partitions.} The external observer's
designated partition here is the cosmological horizon (§2.1). When
additional causal boundaries are present --- most naturally, a black
hole horizon within the cosmological patch --- the trace-out machinery
extends to a nested formalism, addressed in Appendix A. The nested
formalism is a self-consistency exercise: it shows that nested
marginalization agrees under a shared dilation and initial state.
Appendix A also obtains a cycle-averaged Page curve, conditional on the
scrambling condition H-scramble and the decoupling condition H-decouple
and taking the energy shell of the black hole and its radiation to be
their full product space; H-scramble fails for every bijection acting on
an initial configuration with the bipartition aligned to the
configuration coordinates, so that derivation has no instance in the
setting it counts. Appendix A leaves the generalized second law open. It
is \emph{not} a claim that external observers' emergent quantum
descriptions acquire BH-horizon-specific features; under
partition-relativity (above), the cosmological horizon remains the
primary partition. Appendix A develops the formal nested-partition
machinery in three layers.

\emph{(i) Consistency of nested trace-outs (Appendix A, Theorems
A.1--A.4).} Taking the partial traces simultaneously or successively
gives the same channel for the same unitary dilation and product input
on the visible and hidden factors. The generators inferred from
transition data carry the reconstruction hypotheses of §3.3 and their
gauge and antiunitary alternatives. For an exactly unitary channel
family the generators differ only by a scalar shift. Agreement of the
calibrated \(\hbar\) also uses the conditions of §8.5(b) at each
horizon, with the same \(\epsilon\) and \(\sigma\).

\emph{(ii) The generalized second law (Appendix A, §A.6), open.}
Unitarity of \(\varphi\) and subadditivity bound the sum of the von
Neumann entropies of \(V\), \(B\) and \(D\) below by its value at a
product initial state, at every earlier and later time (Theorem A.6);
the bound is two-sided and does not make the sum monotone. That sum is
\(S_{\text{matter}} + S_{\text{BH}} + S_{\text{dS}}\) only if each
horizon's Bekenstein-Hawking entropy equals the von Neumann entropy of
its sector's reduced state, an identification the boundary-mode count
does not supply. A generalized second law would also require the missing
monotonicity.

\emph{(iii) The Page curve (Appendix A, Theorems A.7--A.9),
conditional.} Under the scrambling condition H-scramble and the
decoupling condition H-decouple of Appendix A, with the energy shell of
\(B \cup R\) the full product space
\(\mathcal{H}_B \otimes \mathcal{H}_R\), the entanglement entropy of the
radiation of an evaporating black hole, averaged over full cycles of the
dynamics at the bipartition of stage \(t\), follows
\(\overline{S}_{\text{rad}}(t) = \min(n_R(t), n_B(t)) \cdot \ln q\) up
to a correction of at most \(\tfrac12\), reached at \(n_R = n_B\) and
exponentially small in \(|n_R - n_B|\) away from it (Theorem A.9), where
\(n_B(t)\) is the shrinking boundary mode count and
\(n_R(t) = n_B(0) - n_B(t)\); whether the entropy along a single
trajectory follows the same curve is open. Unitarity is preserved
(bijectivity of \(\varphi\)); at \(\beta = 1\), the radiation gaining
one mode's state-space entropy per boundary mode lost (§A.7), the
turnover occurs at the Page time
\(t_P \approx 0.646 \, t_{\text{evap}}\) (in general
\(t_P/t_{\text{evap}} = 1 - [\beta/(1+\beta)]^{3/2}\); independent of
mass and of the emission coefficient when that coefficient is
mass-independent). Theorem A.8 derives the Page value as a cycle average
from H-scramble, which carries measure concentration {[}29, 61{]} to the
orbit, the cycle decomposition of the finite configuration space
supplying exact time averages. It counts initial configurations, with
the bipartition aligned to the configuration coordinates, and there
H-scramble fails for every bijection: every state on the orbit is a
product state and the radiation's entropy is zero, so Theorems A.8--A.9
as stated have no instance. Under a bijection the radiation's entropy at
an aligned bipartition never exceeds the entropy of the initial state's
distribution over configurations (§A.7), so a Page value requires an
initial superposition whose distribution over configurations has entropy
close to the Page value or above it, or a bipartition not aligned with
the configuration coordinates; neither is specified, and H-scramble is
not established (uniformity over classical microstates is not uniformity
over quantum pure states). H-decouple is not supplied by timescale
ordering. On a proper energy shell the value differs in general, and the
corresponding statement is open. The greybody factor is
solution-specific.

Laboratory tests of the characterization theorem may be possible in
analogue gravity systems where the hidden sector capacity is tunable.

\hypertarget{universality-class-character-of-the-partition-geometry-results}{%
\subsubsection{8.5 Universality-class character of the
partition-geometry
results}\label{universality-class-character-of-the-partition-geometry-results}}

The classical horizon temperature uses the supplied first law and mode
count. The two-level description distinguishes the emergent vacuum
energy from the supplied cosmological density. Matching the thermal
slopes gives the action scale \(\hbar\), the discreteness scale
\(\epsilon=2\sqrt{\sigma}\,l_p\) and the \(1/4\) coefficient under
H-count, H-slope, the horizon and frame conditions, local boost
structure and symmetric detailed-balance rates. The G3 coupling and the
source attribution of the Type II reading have separate requirements
(§§6.2, 7.1).

\textbf{The universality claim, in three parts.} Let \(\mathcal{S}\) be
any finite deterministic substratum (not necessarily OI's cubic lattice)
admitting a partition into a visible sector \(V\) and hidden sector
\(H\) at a cosmological-horizon-like causal boundary, satisfying
{[}Main{]}'s structural conditions C1 (non-trivial coupling), C2 (record
persistence --- slow bath \(\tau_S \ll \tau_B\) the realization), C3
(sufficient memory capacity --- \(N_H \gg N_V\) the comfortable regime),
and C4 (history readback).

\emph{(a) Classical normalization and descriptive distinction.} On the
stated first-law and boundary-count inputs:

(U1) \emph{Classical horizon temperature.} The classical temperature
\(k_B T_{\text{cl}} = c^2\epsilon^2\kappa/(8\pi G\sigma)\) of the
boundary count (§3.1). Its identification with the hidden sector's
energy-shell temperature \((\partial s_H/\partial E)^{-1}\), which (b)
uses, holds under H-count (§2.3).

(U2) \emph{Finite horizon counting.} The boundary entropy is
\(S = \sigma A/\epsilon^2\), for the stated finite mode count
\(A/\epsilon^2\), one mode per cell of area \(\epsilon^2\) (§5).

(U3) \emph{Two-level distinction.} The emergent quantum vacuum energy
(\(\sim M_{\rm Pl}^4\)) and the effective cosmological constant
(\(\sim H^2/G\)) refer to logically distinct objects across two levels
of description; the ratio is of the order of the de Sitter entropy on
the stated scale comparison (§6.3). Its interpretation as a
gravitational cancellation question requires the source law.

\emph{(b) Conditional universality.} If, in addition, H-count holds for
the hidden sector of \(\mathcal{S}\), H-slope, H-frame and the H-Hawking
conditions hold at its horizon, the emergent description has local boost
structure there (the boost-channel Lorentz-sector item of {[}Main
§3.5{]}, carried by the KMS step below) and the symmetric
detailed-balance rate ratio holds (obtained on the OI branch from
microreversibility under the time-reversal-even coarse-state condition
of §2.3; an additional condition for general non-bijective dynamics),
then with the same numerical content as for OI's substratum:

(C1\('\)) \emph{Emergent action scale.}
\(\hbar = c^3 \epsilon^2 / (4G\sigma)\), where \(\epsilon\) is the
partition's discreteness scale and \(\sigma\) the entropy per boundary
mode (§3.2).

(C2\('\)) \emph{Discreteness scale.} \(\epsilon = 2\sqrt{\sigma}\,l_p\)
(\(2\,l_p\) at \(\sigma = 1\)), fixed by the simultaneous solution of
the action-scale relation and the definition of the Planck length (§4).

(C3\('\)) \emph{Area law.} \(S = A/(4\,l_p^2)\), with the \(1/4\)
coefficient the ratio \(2\pi/8\pi\) between the Euclidean KMS period and
Jacobson's Einstein-equation prefactor (§5).

The conditions are statements about the horizon, the hidden sector's
density of states and the emergent dynamics, not about the cubic
lattice, so the conditional is itself a universality-class statement:
any substratum in the class that meets them obtains
(C1\('\))--(C3\('\)).

\emph{(c) Not universality-class results in the present state.} The
dynamical dissolution of the cosmological-constant problem is
conditional on the G3 coupling theorem, whose form S1--S4 do not fix
(§6.2). The Type II running-vacuum functional form is not derived; its
magnitude
\(|\nu| \approx |\nu_{\rm QFT}| \sim (M_{\rm SM}/M_{\rm Pl})^2 \sim 10^{-32}\)
is conditional on the §7.3 split (§7.1). The dark-sector phenomenology
is Tier 2: it inherits the conditions of §3.2 and the dimensional factor
imported in §7.3.

\textbf{What the derivations actually depend on.} Reviewing §3--§7:

\begin{itemize}
\tightlist
\item
  The \emph{\(\hbar\) derivation} (§3) uses the classical shell slope
  under H-count and the emergent frequency slope on the same operational
  transitions under H-slope. On exact de Sitter, the H-Hawking and frame
  conditions, local boost structure and symmetric detailed-balance ratio
  support the calibration \(\beta_E\hbar=\tau_K\). The boost
  prerequisite is tied to the symmetric-point matching under H-blind
  ({[}SM §4.4{]}) and the conditional custodial running, which rests on
  a reported one-loop calculation {[}SM §3.1{]}. The ultraviolet
  dispersion corrections are a separate issue. The spherical single-link
  count supplies the lattice density used here (§3.1). Section 2.2
  describes C1--C4 as conditions for the separate physical
  memory-bearing application. C4 is not a hypothesis of the calibration.
\end{itemize}

\emph{Free-level status, and the narrowed form of the conditional.}
Three results reduce what the boost prerequisite requires. First, the
free-level cone is not exact per channel on the normalized \(d=3\)
observer branch, and the direction dependence runs opposite to what an
axis-based reading suggests. With
\(\cos\omega=\tfrac13\sum_i\cos(ak_i)\) and the clock calibrated by
\(\Delta t=a/(c\sqrt3)\), \[
\frac{\Omega}{ck}=1+\frac{a^2k^2}{24}\Big(\tfrac13-\sum_i n_i^4\Big)+O(a^4k^4),\qquad n_i=k_i/k,
\] and since \(\tfrac13\le\sum_i n_i^4\le1\) the leading deviation lies
in \([-a^2k^2/36,\,0]\). It vanishes on the body diagonals, where all
three arguments coincide and \(\cos\omega=\cos(ak/\sqrt3)\) exactly, and
is maximal along the lattice axes. The \((1{+}1)\) reduction does not
rescue it: zeroing the transverse momenta leaves their two \(+1\)
contributions in the three-dimensional operator, giving
\(\cos\omega=(2+\cos ak)/3\). So the KMS prerequisite is not met exactly
along axes, and an exact all-momentum cone is available only on a
measure-zero set of directions --- not on a generic horizon normal.

What replaces it is dispersive robustness rather than exactness. Under
the Schützhold--Unruh result the Hawking temperature depends on the
product of the low-energy phase and group velocities; writing
\(\Omega=ck[1+q(\hat n)a^2k^2+\cdots]\) gives
\(v_{\rm ph}v_g/c^2=1+4q\,a^2k^2+O(a^4k^4)\) with
\(q=\tfrac1{24}(\tfrac13-\sum_i n_i^4)\), hence
\(-a^2k^2/9\le v_{\rm ph}v_g/c^2-1\le0\) and \[
\Big|\frac{\Delta T_H}{T_H}\Big|\;\lesssim\;\frac19\Big(\frac{\epsilon_{\rm lat}\kappa}{c^2}\Big)^2 .
\] This holds under named conditions, not universally --- call them
\textbf{H-Hawking}: (i) the preferred microscopic frame becomes the
local freely falling horizon frame; (ii) cutoff-scale outgoing modes are
in their freely falling ground state; (iii) cutoff dynamics is adiabatic
relative to \(\kappa\); (iv) the observer dispersion is the normalized
stable infrared branch; (v) the horizon geometry varies smoothly on
scales \(\gg\epsilon\). Conditions (iv) and (v) are discharged by §4.1
and the scale separation; (i) and (ii) are open. The dispersive
corrections vanish as \(k\to0\), so \(\hbar=c^3\epsilon^2/(4G\sigma)\)
and the \(1/4\) coefficient survive as an infrared thermal-calibration
statement, with Planck-suppressed departures at finite frequency --- not
as an exact microscopic KMS theorem. Second, the off-axis and
isotropization corrections belong to the
\(\mathcal{O}((\epsilon\kappa/c^2)^2)\) class already carried above;
numerically, at \(\sigma = 1\),
\((\epsilon\kappa/c^2)^2 \approx 6 \times 10^{-122}\) at the
cosmological horizon (\(4\sigma\) times the dimensionless ratio
\((H\,l_p)^2\) that constitutes the §6 discrepancy --- a scale identity,
recorded as such) and \(\approx 3 \times 10^{-77}\) for a solar-mass
horizon. Third, the conditional accordingly narrows to its interacting
core: radiative velocity-splitting between emergent species --- the
standard protection problem of emergent Lorentz invariance, in which
loop corrections on a discrete substrate generically split species
velocities at \(\mathcal{O}(\alpha)\) unless a mechanism forbids it.
This core is \emph{identical} to the open Lorentz-sector boost item
({[}Main §3.5{]}; {[}Substratum §6.4{]}) --- one open item, not two. The
protection has a matching part and a running part. At the symmetric
point the six carrier components are dynamically identical, so the
substratum dynamics and a component-blind trace-out commute with the
signed permutations of the components, whose six-dimensional
representation is irreducible; under H-blind the complete quadratic
response is custodially scalar and every massless block shares one cone
at the matching scale ({[}SM §4.4{]}, Proposition 4b). Conditional on
that matching, the running is the custodial normalization of the
Lorentz-sector accounting. The running part rests on a calculation
reported in the companion paper ({[}SM §3.1{]}) whose driver and
numerical output are not available for verification: a one-loop lattice
calculation of the boost-projected self-energy curvature reported that,
with the universal tree-level cone absorbed, the sector-differential
coefficient obeys a quadratic-with-logarithm decoupling law
\(\delta c \sim \Delta C_2\,(\alpha_0/4\pi)(M_X/M_{\text{Pl}})^2[\,c_L\ln(M_{\text{Pl}}/M_X)+c_0\,]\)
in the cross-charged condensate mass \(M_X\), rather than the
non-decoupling constant of the generic fine-tuning problem. What is
certified is the symmetry statement --- the signed-permutation symmetry
of the carrier components forbids a condensate-independent differential
counterterm at the regulator level --- not the reported form. What
remains of the conditional, under H-blind, is therefore not the
existence of a protection but its quantitative closure: the reported
decoupling form itself, pending an independent reproduction; the
vertex-weighted one-loop point value of the normalization (structurally
bounded, \(\Delta C_2 \in [0.58, 2.08]\)); the non-logarithmic constant
of that loop (the substratum's own nonlinearity is component-local and
contributes only to the universal cone, {[}SM §3.1{]}); and the
condensate-gap determination of \(M_X\). On the reported form and
constants the allowed window would be \(M_X \lesssim 10^{8}\) GeV on the
worst-case \(\Delta C_2\), \(\sim 4\times10^{7}\) GeV with the reported
kernel constants, against the tightest multi-species bounds, with a
phenomenological floor at \(10^{3}\)--\(10^{6}\) GeV set by leptoquark
bounds --- the cross-charged sector is baryon-safe at the gauge level,
so proton decay does not close the window from below ({[}SM §3.1{]}) ---
an illustrative one-and-a-half-to-four-and-a-half-decade window carrying
the provenance of that calculation. The \(\hbar\) chain's boost
conditional thus has the same two parts the Standard-Model sector
carries: a certified matching theorem and a reported running
calculation, the latter the one item shared across both papers and the
one still to be verified.

\begin{itemize}
\item
  The \emph{discreteness scale fixing} (§4) is purely algebraic given
  the \(\hbar\) relation: rearranging it gives
  \(\epsilon^2 = 4\sigma \hbar G/c^3 = 4\sigma l_p^2\). It inherits that
  relation's conditions and nothing else.
\item
  The \emph{area law} (§5) uses one counted mode per area
  \(\epsilon^2\), with entropy \(\sigma\) per mode. The \(1/4\) uses the
  thermal calibration, including its named conditions. The single-link
  lattice reading uses the spherical or equal-density scope.
\item
  The \emph{vacuum-energy comparison} (§6) distinguishes the emergent
  field-theory vacuum energy from the supplied cosmological density. The
  gravitational coupling of a vacuum offset and the native stress or
  flux require G3.
\item
  The \emph{Type II reading} (§7.1) uses the variation of the boundary
  count with \(H\) as motivation. The smooth local expansion is generic,
  while conserved matter and running \(G\) describe the chosen Type II
  implementation. Its source and the partition-dependent attribution of
  the coefficient remain open.
\end{itemize}

\textbf{Sufficient structural conditions.} The unconditional results
(U1)--(U3) require the following structural conditions on the
substratum:

(\emph{S1}) \emph{Cosmological-horizon partition.} The substratum admits
a partition \(V\) at a causal horizon with timescale
\(\tau_B \sim H^{-1}\) and information capacity
\(S_{\rm dS} \sim 10^{122}\).

(\emph{S2}) \emph{C1--C3 satisfied.} The partition has substantial
coupling (C1), slow-bath structure (C2), and high capacity (C3), per
{[}Main §3{]}.

(\emph{S3}) \emph{Boundary mode count.} The model supplies a positive
area \(\epsilon^2\) per counted mode. The finite distinguishability of
Lemma 1 gives a finite visible configuration space and does not fix this
physical boundary density. On the spherical single-link lattice branch,
\(\epsilon^2=\tfrac23\epsilon_{\rm lat}^2\) (§3.1). The physical scale
is obtained on the stated calibration branch.

(\emph{S4}) \emph{Stinespring/partial-trace structure.} The
visible-sector dynamics is generated by partial trace of
substratum-level dynamics over the hidden sector, with the trace
producing a CPTP channel admitting Stinespring dilation {[}Main §3.2{]}.

These conditions are weaker than OI's full structural conditions A1--A6
except at S3: A1 (finiteness) makes the number of configurations finite
but supplies neither the discreteness scale \(\epsilon\) of the boundary
count nor its relation to the substratum's lattice structure, which S3
requires; A6 (background independence) is required for the SM
gauge-group derivation but not for the partition-geometry results; A2-A5
(determinism, bounded coupling, center independence, linearity) are
required for OI's specific universality-class representative but not for
the universality-class-level claims. A substratum satisfying S1-S4 but
violating some of A2-A5 (e.g., a substratum with stochastic dynamics, or
with nonlinear evolution) would still produce (U1)--(U3). The
calibration \(\hbar = c^3 \epsilon^2/(4G\sigma)\),
\(\epsilon = 2\sqrt{\sigma}\,l_p\) and the area law with its \(1/4\)
coefficient use in addition, on the classical side, the symmetric
detailed-balance rate ratio used above --- obtained on the OI branch
from microreversibility under the stated time-reversal-even coarse-state
condition; for general non-bijective dynamics that ratio is an
additional condition, supplied neither by S1--S4 nor by H-slope, H-frame
and H-Hawking, and the calibration holds for such a substratum only
where it is established, with the partial-trace machinery and thermal
self-consistency.

\textbf{Theorem-level statement.} The universality claim can be stated
at the level of the framework's hierarchical structure (per {[}Structure
§2{]}):

\emph{Theorem (universality of the partition-geometry results).} For any
substratum \(\mathcal{S}\) satisfying S1--S4, (U1)--(U3) follow with the
same content as derived in this paper for OI's specific substratum; if
in addition H-count holds for its hidden sector, H-slope, H-frame and
the H-Hawking conditions hold at its horizon and the emergent
description has local boost structure there, together with the symmetric
detailed-balance rate ratio used above --- obtained on the OI branch
under the stated time-reversal-even coarse-state condition, and an
additional condition for general non-bijective dynamics --- so do
(C1\('\))--(C3\('\)) with the same numerical content.

The proof is the §3--§6 derivations themselves, with the observation
that each step a part uses depends only on S1--S4 and the named
conditions, not on OI's specific structural conditions A2--A5. Part (b)
uses the calibration \(a_*=\tau_K/\beta_E\) (§2.3) with the thermal
matching of §3.2, Step 4, and §4--§5; Steps 2--3 of §3.2, whose lemma
assumes a site-factorized bijective update and, for the one-layer
boundary, H-BD, are a consistency check that part (b) does not use. A
formal proof at the level of explicit theorem-style presentation has not
been written out in this paper; the present section identifies the
structural conditions required at each step, and the claim is
established at that level.

\textbf{Empirical implication.} The universality-class results have
limited direct empirical content. GW250114 supports the classical area
theorem through Kerr area inference (§5). The \(1/4\) thermal
coefficient has no direct measurement there. Conversely, the dark-sector
concordance (§7.2), the MOND-like phenomenology (§7.3), and the
rotation-curve crossover (§7.3) concern the specified branch, with the
source map and heat input open and the dimensional factor imported
(§§7.2--7.3).

The stratification clarifies what would falsify what: a measured
violation of the area theorem would challenge the classical horizon
dynamics the framework presupposes; a direct empirical test of the
\(1/4\) coefficient, if one becomes available, would test the
universality-class claim; a measured deviation from the dark-sector
phenomenology would falsify OI's specific representative without
necessarily refuting the universality class. The framework's empirical
content thus operates at multiple hierarchy levels, with falsifiability
stratified accordingly.

\begin{center}\rule{0.5\linewidth}{0.5pt}\end{center}

\hypertarget{relation-to-other-entropic-and-emergent-gravity-programs}{%
\subsubsection{8.6 Relation to other entropic and emergent-gravity
programs}\label{relation-to-other-entropic-and-emergent-gravity-programs}}

The thermodynamic derivation of §3--§5 places this framework within a
broader program that treats gravity as entropic or emergent rather than
fundamental. Jacobson's equation-of-state derivation {[}1{]} is the
direct antecedent used here; Verlinde's entropic-force account {[}50{]}
and Padmanabhan's thermodynamic-aspects program {[}51{]} reach
Einstein-level dynamics from horizon thermodynamics and holographic
equipartition, sharing with this framework both the entropic premise and
a holographic/boundary counting of degrees of freedom. More recently,
Bianconi's \emph{gravity from entropy} {[}52{]} derives modified
Einstein equations from an action built on the quantum relative entropy
between the spacetime metric and a matter-induced metric, reproducing
the low-energy Einstein equations and an emergent positive cosmological
constant \emph{without} invoking horizon entropy or holography. The
existence of this last route is worth stating plainly: ``gravity from
entropy'' does not uniquely select the horizon/holographic architecture
used here, so this framework's Jacobson--Bekenstein route is a specific
commitment among alternatives rather than the only entropic path to
Einstein dynamics. What this framework adds to the landscape is not the
entropic premise --- which it shares --- but the discrete substratum
beneath the thermodynamics ({[}Main{]}; {[}SM{]}) and the quantitative
consequences specific to it (the fixed \(1/4\) coefficient and
\(\epsilon = 2\sqrt{\sigma}\,l_p\), under the conditions of §8.5(b)).

There is a qualitative convergence: entropic and emergent-gravity
programs generically carry a non-constant dark-energy sector rather than
a bare constant \(\Lambda\) --- the Type II running-vacuum reading of
§7.1 here, and the always-positive G-field term of {[}52{]}. This
convergence is at the level of the qualitative form (dark energy as a
non-constant feature) only; it does not transfer magnitude, sign, or
mechanism between programs, which differ substantially. For the
running-vacuum coefficient, §7.1 gives the order-of-magnitude estimate
\(|\nu_{\rm QFT}|\sim10^{-32}\) for the emergent field channel,
conditional on the source reading and split stated there; the light
Standard Model spectrum gives a negative value, heavy states could
change its sign, and the substratum contribution \(\nu_{\rm sub}\)
depends on an open spectral weighting.

\hypertarget{kinematic-emergence-of-geometry-from-the-substratum}{%
\subsubsection{8.7 Kinematic emergence of geometry from the
substratum}\label{kinematic-emergence-of-geometry-from-the-substratum}}

The derivations of §3--§5 and the necessity routes of §3.2 (\emph{An
independent necessity statement}; gravitational closure, Lovelock,
spin-2 bootstrap) all take a smooth Lorentzian metric as given: the
former use its Einstein \emph{dynamics} through Jacobson's relation
(§3.1), and the latter would fix that dynamics from hypotheses the
framework has not established. A logically prior question is
\emph{kinematic}: does the discrete substratum ({[}SM{]}) actually yield
a smooth pseudo-Riemannian geometry, with the right dimension and
curvature, in the first place? §3.3 treats the continuum description as
an effective approximation. This subsection reports a research program
that makes that approximation concrete at the weak-field level, by
adapting standard discrete-geometry tools --- the spectral dimension
{[}53{]} and the Ollivier--Ricci curvature {[}54{]} --- to the
substratum graph. The results are consistency checks reproducing known
Newtonian physics; they add no new empirical content, and the program's
honest boundary is stated at the end.

\emph{Dimension and curvature tool.} The cubic substratum has spectral
dimension \(d_s \to 3\) (heat-kernel return probability), which with the
emergent time direction is the target \(3+1\). Encoding a chosen
emergent metric into the graph through the framework's own dispersion
relation (bond weights fixed by the local light-cone speed) and
computing curvature two independent ways --- combinatorially
(Ollivier--Ricci) and analytically (the Riemannian curvature of that
metric) --- the two agree (correlation \(\approx 0.97\);
\(\kappa \approx \tfrac12 K\) as expected in two dimensions). So the
discrete-curvature tool faithfully computes the framework's \emph{own}
emergent-metric geometry rather than an unrelated combinatorial
quantity.

\emph{Weak-field gravity at the fixed \(G\).} In the weak field the
Newtonian force is carried by the emergent-time-rate component
\(g_{00}\) (the local substratum clock rate), not by spatial curvature;
\(g_{00}\)-only geodesics reproduce closed Keplerian orbits obeying
Kepler's third law. The Newtonian field itself follows from holographic
equipartition using the framework's own constants
(\(\eta = \sigma/\epsilon^2\); \(\hbar = c^3\epsilon^2/(4G\sigma)\))
together with the classical horizon temperature of §3.1: one obtains
\(a = GM/r^2\) with the framework's \emph{fixed} \(G\) --- fixed through
\(\epsilon = 2\sqrt{\sigma}\,l_p\), so inheriting the conditions of §3.2
--- and the Poisson equation \(\nabla^2\Phi = 4\pi G\rho\) for extended
sources. The two ingredients that are heuristic in the usual
entropic-gravity argument are, here, consequences of the framework's
structure, the second on an assumed energy shell: equipartition is the
equipartition \emph{theorem} because the emergent boundary modes are
harmonic (the {[}SM{]} wave equation), and the entropic force
\(F = T\,dS/dx\) is a theorem of the reversible microcanonical dynamics
of {[}Main{]} Lemma 3 restricted to a conserved energy shell. The shell
is assumed, as in §2.3; it is not a consequence of Lemma 3, which
supplies determinism, injectivity and the selected measure but no
conserved energy. Notably the \emph{field} derivation requires no
entropy-displacement (Bekenstein) input; the force on a test mass is
field \(\times\) emergent inertia.

\emph{Honest boundary.} These are weak-field, consistency-level results:
they reproduce Newtonian gravity at the framework's own \(G\), not a new
prediction. Three items remain genuinely open. (i) The equipartition
degree-of-freedom normalization
\(N_{\rm dof} = A/l_p^2 = 4\,S_{\rm BH}\) --- the standard relation in
all entropic gravity --- is here \emph{adopted}, not derived from the
substratum's microscopic component count (\(K = 2d = 6\) per site;
boundary-crossing components at the link density of §3.1,
\(A/\epsilon^2 = S_{\rm BH}/\sigma\) modes); reconciling these counts is
unsolved. (ii) The full nonlinear \(G_{\mu\nu} = 8\pi G\,T_{\mu\nu}\),
and the Lorentzian/\(3{+}1\) generalization of the (Euclidean,
conformal) curvature check, are not established. (iii) For the optional
stronger claim that inertia itself is entropic, the
entropy/Compton-phase identification remains an assumption. The
program's value is thus (a) showing the discrete-curvature approach
targets the framework's real geometry and (b) turning §3.3's effective
continuum description into a computation that reaches the Newtonian
checkpoint with the remaining steps explicitly marked --- not a claim to
have derived general relativity from the substratum.

\hypertarget{conclusion}{%
\subsection{9. Conclusion}\label{conclusion}}

An exact de Sitter horizon supplies a candidate causal partition. Its
physical memory-bearing application has the conditions of §2.2. On
identified transitions, matching the H-count and H-slope thermal slopes
under the horizon, frame and detailed-balance conditions gives
\(\hbar=c^3\epsilon^2/(4G\sigma)\) and \(\epsilon=2\sqrt{\sigma}\,l_p\).
The calibrated count gives \(S=A/(4l_p^2)\) on its stated density scope.
The \(1/4\) is the ratio \(2\pi/8\pi\) of the thermal and first-law
factors. GW250114 supports the classical area theorem through
waveform-based area inference {[}5{]}. The descriptive vacuum-energy
distinction leaves Einstein dynamics and the physical source coupling
open at G3.

The Type II late-time reading uses matter conservation and running
\(G\). Its near-null estimate is conditional on the source split and
spectral assumptions of §7.1. The Bertini et al.~fit {[}15{]} is
compatible with a near-null result and with \(\Lambda\)CDM. Section 7.2
gives a horizon-energy mean equal to the supplied critical density, with
the matter coupling and the interpretation of \(1-\Omega_B\) open.
Direct-detection searches constrain the tested masses and interactions.
Under the stated de Sitter and heat-transfer inputs, with the factor
\(1/6\) imported from Verlinde, the displacement model gives
\(a_0=cH/6\) and the deep-MOND Tully-Fisher relation. The outer declines
reported for the Milky Way and high-redshift galaxies remain in tension
with the radial law. Quantitative comparisons require the respective
baryonic profiles and kinematics (§7.3).

The foundational framework developed in {[}Main{]} and the cubic-lattice
realization developed in {[}SM{]} are complementary to the present work:
{[}Main{]} establishes a finite-horizon correspondence between the
recordable laws of an embedded observer and fixed-basis quantum laws,
the present paper considers its conditional application at the
cosmological horizon, and the SM paper applies it to the microscopic
lattice substratum. The reconstruction theorem and the synthesis claim
--- that quantum mechanics and general relativity descend as projections
of the same deterministic \((S, \varphi)\), with a third projection
bearing on the arrow of time --- are developed in {[}Substratum{]}.
Together, this four-paper synthesis ({[}Main{]}, {[}SM{]}, {[}GR{]},
{[}Substratum{]}) proposes that the Standard Model and general
relativity descend from one finite deterministic construction, with the
covariant field equations (G3) and the selection of full operational
quantum mechanics open. The framework's hierarchical structural realism
and its relationship to other unification programs are developed in the
fifth core paper {[}Structure{]}.

\begin{center}\rule{0.5\linewidth}{0.5pt}\end{center}

\hypertarget{appendix-a-nested-partitions}{%
\subsection{Appendix A: Nested
Partitions}\label{appendix-a-nested-partitions}}

This appendix develops the nested-partition formalism for the framework,
establishing the self-consistency of multiple coexisting causal
boundaries (typically a black hole horizon inside a cosmological
horizon), deriving a cycle-averaged Page curve under the scrambling
condition H-scramble and the decoupling condition H-decouple when the
energy shell is the full product space (H-scramble fails for every
bijection acting on an initial configuration with the bipartition
aligned to the configuration coordinates, the setting the derivation
counts), and stating the entropy inequality the formalism supplies
toward the generalized second law.

\hypertarget{a.1-setup}{%
\subsubsection{A.1 Setup}\label{a.1-setup}}

Let \(\mathcal{C}_V\), \(\mathcal{C}_B\), \(\mathcal{C}_D\) be finite
configuration spaces with \(|\mathcal{C}_V|, |\mathcal{C}_B| \geq 2\),
and let
\(\varphi: \mathcal{C}_V \times \mathcal{C}_B \times \mathcal{C}_D \to \mathcal{C}_V \times \mathcal{C}_B \times \mathcal{C}_D\)
be a bijection. The physical picture: \(V\) is the region between a
black hole horizon and the cosmological horizon, \(B\) is the black hole
interior, \(D\) is the trans-cosmological region. Define
\(W = V \cup B\) (the cosmological interior).

Two procedures produce a quantum description on \(V\):

\emph{Direct:} Trace out \(B \cup D\) in one step, marginalizing over
\(\mathcal{C}_B \times \mathcal{C}_D\).

\emph{Sequential:} (Stage 1) Trace out \(D\), producing emergent QM on
\(W\). (Stage 2) Within the emergent QM on \(W\), trace out \(B\),
producing a further-reduced description on \(V\).

\hypertarget{a.2-classical-associativity}{%
\subsubsection{A.2 Classical
Associativity}\label{a.2-classical-associativity}}

\textbf{Theorem A.1.} \emph{The direct and sequential procedures produce
the same transition probabilities on \(\mathcal{C}_V\).}

\emph{Proof.} The direct transition matrix is:

\[T^{\text{dir}}_{ij}(t) = \frac{1}{|\mathcal{C}_B||\mathcal{C}_D|} \sum_{b \in \mathcal{C}_B} \sum_{d \in \mathcal{C}_D} \delta_{x_j}[\pi_V(\varphi_t(x_i, b, d))]\]

The sequential procedure marginalizes over \(D\) (Stage 1) then over
\(B\) (Stage 2). Composing:

\[T^{\text{seq}}_{ij}(t) = \frac{1}{|\mathcal{C}_B|} \sum_b \frac{1}{|\mathcal{C}_D|} \sum_d \delta_{x_j}[\pi_V(\varphi_t(x_i, b, d))] = T^{\text{dir}}_{ij}(t)\]

The sums commute (Fubini on a finite product). \(\square\)

\hypertarget{a.3-quantum-classical-consistency}{%
\subsubsection{A.3 Quantum-Classical
Consistency}\label{a.3-quantum-classical-consistency}}

\textbf{Theorem A.2} (Quantum-classical consistency of nested
trace-outs). \emph{Let \(\Phi_V^{\text{dir}}\) and
\(\Phi_V^{\text{seq}}\) be the CPTP channels on \(\mathcal{H}_V\)
produced by the direct and sequential procedures respectively. Then
\(\Phi_V^{\text{dir}} = \Phi_V^{\text{seq}}\).}

\emph{Proof.} By the Stinespring dilation theorem {[}30, 31{]},
\(\varphi\) defines a unitary \(U_\varphi\) on
\(\mathcal{H} = \mathcal{H}_V \otimes \mathcal{H}_B \otimes \mathcal{H}_D\).
The direct channel is
\(\Phi_V^{\text{dir}}(\rho_V) = \text{Tr}_{BD}[U_\varphi\,(\rho_V \otimes \rho_B \otimes \rho_D)\,U_\varphi^\dagger]\).
The sequential channel is
\(\Phi_V^{\text{seq}}(\rho_V) = \text{Tr}_B[\text{Tr}_D[U_\varphi\,(\rho_V \otimes \rho_B \otimes \rho_D)\,U_\varphi^\dagger]]\).
On finite-dimensional tensor products,
\(\text{Tr}_B \circ \text{Tr}_D = \text{Tr}_{BD}\), giving
\(\Phi_V^{\text{seq}} = \Phi_V^{\text{dir}}\). \(\square\)

\hypertarget{a.4-hamiltonian-and-hbar-consistency}{%
\subsubsection{\texorpdfstring{A.4 Hamiltonian and \(\hbar\)
Consistency}{A.4 Hamiltonian and \textbackslash hbar Consistency}}\label{a.4-hamiltonian-and-hbar-consistency}}

\textbf{Corollary A.3} (Hamiltonian consistency). \emph{The direct and
sequential procedures yield the same emergent dynamics in two regimes.
At the reconstruction level, conditional on H-BG and the §3.3
hypotheses: Theorem A.2 gives equal transition data, so under the §3.3
hypotheses
\(\hat{H}_V^{\text{dir}} = D\,\hat{H}_V^{\text{seq}}\,D^\dagger + E_0\)
or
\(\hat{H}_V^{\text{dir}} = -D\,\bar{\hat{H}}_V^{\text{seq}}\,D^\dagger + E_0\),
for a time-independent diagonal unitary \(D\) and a real \(E_0\). If
moreover the common channel family of Theorem A.2 is exactly unitary,
\(\Phi_t = \mathrm{Ad}_{U(t)}\) for both procedures, then
\(\mathrm{Ad}_{U_{\text{dir}}(t)} = \mathrm{Ad}_{U_{\text{seq}}(t)}\)
forces \(U_{\text{dir}}(t) = e^{i\theta(t)}\,U_{\text{seq}}(t)\) with
\(\theta(t) = -E_0 t\), and
\(\hat{H}_V^{\text{dir}} = \hat{H}_V^{\text{seq}} + E_0\) --- no
antiunitary branch and no nontrivial \(D\).}

\emph{Proof.} Theorem A.2 gives
\(\Phi_V^{\text{dir}} = \Phi_V^{\text{seq}}\), hence
\(T^{\text{dir}}_{ij}(t) = T^{\text{seq}}_{ij}(t)\) for all \(t\) (Born
rule applied to diagonal elements), and the conditional two-branch
theorem (§3.3), including H-BG, gives the reconstructed form. In the
exact-unitary regime, equality of the conjugation actions at each time
makes \(U_{\text{seq}}(t)^\dagger U_{\text{dir}}(t)\) commute with every
visible operator, hence a unimodular scalar (kernel:
\texttt{ad\_eq\_scalar}); the resulting identity of spectral phase
families \(e^{-iE_a t} = c(t)\,e^{-iE'_a t}\) forces the eigenvalue
differences \(E_a - E'_a\) to be mode-independent (kernel:
\texttt{phase\_families\_shift}), so \(\theta(t) = -E_0 t\) and the
generators differ by the scalar shift alone. Exact channel equality in
the same visible algebra carries the coherence information that
transition probabilities discard, which is why the antiunitary branch
and the diagonal gauge both disappear; a general CPTP reduced dynamics,
by contrast, has no uniquely defined Hamiltonian part, and the
reconstruction-level clause claims none. \(\square\)

\textbf{Corollary A.4} (\(\hbar\) universality). \emph{Under the
conditions of §8.5(b) at each horizon, with the same entropy \(\sigma\)
per boundary mode and the same discreteness scale \(\epsilon\) of the
boundary count at each horizon (under the single-link reading of §3.1,
the same for horizons that are spheres in the lattice frame), both
procedures yield the same emergent action scale
\(\hbar = c^3 \epsilon^2 / (4G\sigma)\).}

\emph{Proof.} At each horizon the calibration gives
\(\hbar=\tau_K/\beta_E=c^3\epsilon^2/(4G\sigma)\) (§2.3; §3.2, Step 4),
which contains no area and from which \(\kappa\) cancels, so the two
calibrated values agree at the matched \(\epsilon\) and \(\sigma\).
\(\square\)

\hypertarget{a.5-additive-dissipators}{%
\subsubsection{A.5 Additive
Dissipators}\label{a.5-additive-dissipators}}

\textbf{Theorem A.5} (Frozen deep sector: mixture of boundary channels).
\emph{Let \(\varphi\) be a bijection of
\(\mathcal{C}_V \times \mathcal{C}_B \times \mathcal{C}_D\) that leaves
every \(D\)-configuration fixed, \(\varphi(x, d) = (\varphi_d(x), d)\),
and let the initial state be \(\rho_V \otimes \rho_B \otimes \rho_D\)
with \(\rho_D = \sum_d p_d\,|d\rangle\langle d|\). Then each
\(\varphi_d\) is a bijection of \(\mathcal{C}_V \times \mathcal{C}_B\),
and at every time the reduced channel on \(\mathcal{H}_V\) is the
mixture \(\Phi_V = \sum_d p_d\,\Phi_V^{(d)}\), where \(\Phi_V^{(d)}\) is
the reduced channel of \(\varphi_d\) with \(B\) as the hidden sector.}

\emph{Proof.} \(\varphi\) maps each fiber
\(\mathcal{C}_V \times \mathcal{C}_B \times \{d\}\) into itself and is
injective, so its restriction \(\varphi_d\) is an injection of a finite
set into itself, hence a bijection; the same holds for every power of
\(\varphi\). \(U_\varphi\) is therefore block-diagonal,
\(U_\varphi = \sum_d U_{\varphi_d} \otimes |d\rangle\langle d|\), and
\(\text{Tr}_{BD}[U_\varphi(\rho_V \otimes \rho_B \otimes \rho_D)U_\varphi^\dagger] = \sum_d p_d\,\text{Tr}_B[U_{\varphi_d}(\rho_V \otimes \rho_B)U_{\varphi_d}^\dagger]\).
\(\square\)

Each \(\Phi_V^{(d)}\) is the Stinespring setup of {[}Main, §3.2{]} with
\(B\) as the hidden sector, whose departure \(\mathcal{D}_V^{(B)}\) from
its frozen-\(B\) description is of order \(\tau_S/\tau_B^B\), the
non-absorbable (dissipative) part entering at second order ({[}Main,
§3.2{]}). For a deep sector that moves, with \(D\)-sector timescale
\(\tau_B^D\) and coupling chain
\(V \leftrightarrow B \leftrightarrow D\), {[}Main, §3.2{]} applied with
\(B \cup D\) as the hidden sector gives a first-order absorbable
Hermitian drift and second-order dissipation. The additive form of that
correction on observation timescales \(t \ll \min(\tau_B^D, \tau_B^B)\),

\[\mathcal{D}_V = \mathcal{D}_V^{(B)} + \mathcal{D}_V^{(D)} + \mathcal{R}, \qquad \|\mathcal{D}_V^{(D)}\| = \mathcal{O}(\tau_S / \tau_B^D), \qquad \|\mathcal{R}\| = \mathcal{O}(\tau_S^2 / (\tau_B^B \cdot \tau_B^D)),\]

is open: no bound on the \(D\)-correction or on the cross term is
derived.

\textbf{Orders of magnitude.} For a \(30 M_\odot\) black hole inside the
cosmological horizon, \(\tau_S \sim 10^{-15}\) s,
\(\tau_B^B \sim \tau_{\text{scr}} \sim \beta \log S_{\text{BH}} \sim 10^{-3}\)
s and \(\tau_B^D \sim 1/H \sim 10^{17}\) s. The {[}Main, §3.2{]} order
of \(\mathcal{D}_V^{(B)}\) is then \(\tau_S/\tau_B^B \sim 10^{-12}\).
The additive form is stated for the chain
\(V \leftrightarrow B \leftrightarrow D\), not for this black hole,
which has the chain \(B \leftrightarrow V \leftrightarrow D\) (below).

\textbf{Scope of the coupling chain.} The chain
\(V \leftrightarrow B \leftrightarrow D\) is a hypothesis of the
additive form; Theorem A.5 uses no chain. It is not the geometry of
§A.1, in which \(V\) lies between the black-hole horizon and the
cosmological horizon, with \(B\) inside the first and \(D\) beyond the
second; spatial locality there gives the chain
\(B \leftrightarrow V \leftrightarrow D\), with \(D\) coupled to \(V\)
and not to \(B\). For the chain of §A.1, where \(D\)'s correction
reaches \(V\) directly, the decomposition is not derived here.

\hypertarget{a.6-the-generalized-second-law}{%
\subsubsection{A.6 The Generalized Second
Law}\label{a.6-the-generalized-second-law}}

The Bekenstein-Hawking entropy of a horizon is the information content
of the corresponding partition boundary (§5, under the conditions of the
§3 calibration): \(S_{\text{BH}} = A_{\text{BH}} / (4\,l_p^2)\) at the
black hole horizon and \(S_{\text{dS}} = A_{\text{dS}} / (4\,l_p^2)\) at
the cosmological horizon, each carried by \(S/\sigma\) boundary modes.
The matter entropy is the von Neumann entropy of the visible-sector
state: \(S_{\text{matter}} = -\text{Tr}[\rho_V \log \rho_V]\).

The generalized second law is the statement that
\(S_{\text{matter}} + S_{\text{BH}} + S_{\text{dS}}\) is non-decreasing.
The nested formalism supplies a bound relative to a product state at a
reference time.

\textbf{Theorem A.6} (Sector-entropy inequality). \emph{Let the state on
\(\mathcal{H}_V \otimes \mathcal{H}_B \otimes \mathcal{H}_D\) at
\(t = 0\) be a product \(\rho_V(0) \otimes \rho_B(0) \otimes \rho_D(0)\)
evolving under \(U_\varphi\). Then for every \(t\), and for every
grouping of the same tensor factors into three sectors \(V_t\), \(B_t\),
\(D_t\),}

\[S(\rho_{V_t}(t)) + S(\rho_{B_t}(t)) + S(\rho_{D_t}(t)) \geq S(\rho_V(0)) + S(\rho_B(0)) + S(\rho_D(0)).\]

\emph{Proof.} Subadditivity of von Neumann entropy {[}31{]} bounds the
left side below by the joint entropy \(S(\rho(t))\). Unitary evolution
preserves the joint entropy, \(S(\rho(t)) = S(\rho(0))\), and the
entropy of a product state is the sum of its factors' entropies.
\(\square\)

\textbf{The generalized second law is open.} Theorem A.6 bounds the sum
of the sectors' von Neumann entropies. A generalized second law requires
both the identification \(S_{\text{BH}}(t)=S(\rho_{B_t}(t))\),
\(S_{\text{dS}}(t)=S(\rho_{D_t}(t))\) (which the boundary count does not
supply) and monotonicity of the identified sum between successive times.
The reference-state bound allows a decrease after an increase, including
under a unitary involution taking a product state to an entangled state
and back. \(S_{\text{BH}} = n_B\,\sigma\) counts the horizon's modes,
whereas in the evaporation of §A.7, where \(B \cup R\) is effectively
pure under H-decouple, \(S(\rho_B) = S(\rho_R) \leq n_R \ln q\): it
vanishes at \(n_R = 0\), where \(S_{\text{BH}} = n_B(0)\,\sigma\).
Entropy increase for a single reduced channel has a different condition.
Relative entropy contracts under every CPTP map. For a channel whose
input and output spaces have equal dimension, the output entropy is at
least the input entropy for every input exactly when the channel is
unital: unitality and contraction of
\(D(\rho\,\|\,I/d) = \ln d - S(\rho)\) give one direction, and the input
\(I/d\) gives the converse. The marginal channel \(\Phi_V\) is unital
when the hidden-sector state is maximally mixed, and its entropy
non-decrease then bounds \(S_{\text{matter}}\) alone, not the horizon
terms. The horizon-entropy identification and dynamical monotonicity
remain open.

\hypertarget{a.7-the-page-curve}{%
\subsubsection{A.7 The Page Curve}\label{a.7-the-page-curve}}

\textbf{Setup.} A black hole with initial horizon area \(A_0\)
evaporates. The partition boundary at the black hole horizon moves:
modes transition from \(B\) (interior) to \(R\) (radiation, a subsystem
of \(V\)) as the horizon shrinks. Define
\(n_B(t) = A_{\text{BH}}(t)/\epsilon^2 = S_{\text{BH}}(t)/\sigma\) (the
number of boundary modes at the BH horizon, each carrying the entropy
\(\sigma\) of §3.1; \(A_{\text{BH}}(t)/(4\,l_p^2)\) at \(\sigma = 1\))
and \(n_R(t) = n_B(0) - n_B(t)\) (the number of emitted radiation
modes). Conservation:
\(n_B(t) + n_R(t) = n_B(0) = S_{\text{BH}}(0)/\sigma\).

The hypothesis of Theorem A.7 below freezes \(D\) only on timescales
\(t \ll \tau_B^D\), with \(\tau_B^D \sim 1/H \sim 10^{17}\) s (§A.5).
The evaporation is far longer: the trajectory below gives
\(t_{\text{evap}} \sim 10^{79}\) s for a \(30 M_\odot\) black hole at
\(\Gamma \sim 1\), so \(t_{\text{evap}}/\tau_B^D \sim 10^{62}\), and
\(t_{\text{evap}} > \tau_B^D\) for every black hole heavier than about
\(10^{11}\) kg. Theorem A.7 is stated for \(t \ll \tau_B^D\). Theorems
A.8--A.9 treat \(B \cup R\) as a closed system, on which \(\varphi\)
acts as a bijection of its configuration space
\(\mathcal{C}_B \times \mathcal{C}_R\) (in Theorem A.8 a single energy
shell, whose span is the full product space
\(\mathcal{H}_B \otimes \mathcal{H}_R\)), so that a pure joint state of
\(B\) and \(R\) stays pure; applied to the black hole and its radiation,
they therefore rely on the decoupling condition \textbf{H-decouple}:
\(B \cup R\) evolves with no coupling to the trans-cosmological sector
\(D\) or to the rest of the visible sector, \(V \setminus R\), and its
joint state is pure. The full-cycle averages of Theorem A.9 use
H-decouple over whole cycles of the dynamics, not only over the
evaporation. Timescale ordering does not supply it, and H-decouple is
not established here.

\textbf{Theorem A.7} (Effective pure state). \emph{Let the initial state
be \(|\psi_{VB}(0)\rangle \otimes |d_0\rangle\), with
\(d_0 \in \mathcal{C}_D\) a configuration, and assume that on timescales
\(t \ll \tau_B^D\) the \(D\)-configuration of \(\varphi_t(x, d_0)\)
equals \(d_0\) except for \(x\) in a set \(E_t\) of weight
\(w_t = \sum_{x \in E_t} |\psi(x)|^2 = \mathcal{O}(t / \tau_B^D)\). Then
on these timescales the joint state of \(V\) and \(B\) is effectively
pure:
\(\rho_{VB} = |\psi_{VB}(t)\rangle\langle\psi_{VB}(t)| + \mathcal{O}(t / \tau_B^D)\),
where \(|\psi_{VB}(t)\rangle \in \mathcal{H}_V \otimes \mathcal{H}_B\).}

\emph{Proof.} The total system evolves unitarily under \(U_\varphi\)
(§A.3). Write \(|\psi_{VB}(0)\rangle = \sum_x \psi(x)\,|x\rangle\) over
\(x \in \mathcal{C}_V \times \mathcal{C}_B\), so that
\(|\Psi(t)\rangle = \sum_x \psi(x)\,|\varphi_t(x, d_0)\rangle\). By
hypothesis the \(D\)-configuration of \(\varphi_t(x, d_0)\) equals
\(d_0\) except on \(E_t\). Since \(\varphi_t\) is injective, the terms
with \(x \notin E_t\) sum to \(|a(t)\rangle \otimes |d_0\rangle\) with
\(\||a(t)\rangle\|^2 = 1 - w_t\), and the remaining terms are orthogonal
to every state whose \(D\)-configuration is \(d_0\). Hence
\(\langle d_0|\rho_D(t)|d_0\rangle = 1 - w_t\), and the largest
eigenvalue of \(\rho_D(t)\) is at least \(1 - w_t\). For a pure total
state, \(\rho_{VB}\) and \(\rho_D\) have the same nonzero eigenvalues
(Schmidt decomposition across \(VB|D\)), so \(\rho_{VB}(t)\) lies within
trace distance \(2 w_t\) of
\(|\psi_{VB}(t)\rangle\langle\psi_{VB}(t)|\), with
\(|\psi_{VB}(t)\rangle\) its eigenvector of largest eigenvalue.
\(\square\)

\textbf{Remark (initial state of \(D\)).} The theorem uses that \(D\)
starts in a configuration; a pure but superposed initial state of \(D\)
does not suffice. When the dynamics of \(V \cup B\) depends on \(D\)'s
configuration, \(D\)'s coherences are lost while its configuration stays
frozen: with \(V\), \(B\) and \(D\) qubits and
\(\varphi(v, b, d) = (v, b \oplus d, d)\), every \(D\)-configuration is
fixed, and from \(|0\rangle_V |0\rangle_B |+\rangle_D\) one step gives a
\(\rho_{VB}\) of purity \(1/2\).

The restriction of \(|\psi_{VB}\rangle\) to \(B \cup R\), where
\(R \subset V\) is the radiation, is pure only when the rest of the
visible sector, \(V \setminus R\), is unentangled with \(B \cup R\);
Theorem A.7 does not establish this. Theorems A.8--A.9 take \(B \cup R\)
to be a closed system in a pure joint state on
\(\mathcal{H}_B(t) \otimes \mathcal{H}_R(t)\); that is the condition
H-decouple.

\textbf{Theorem A.8} (Cycle typicality, conditional on H-scramble).
\emph{Let \(\varphi\) act on \(B \cup R\) with its whole configuration
space as a single energy shell,
\(\Sigma_E = \mathcal{C}_B \times \mathcal{C}_R\), whose span is the
full product space \(\mathcal{H}_B \otimes \mathcal{H}_R\), and let
H-scramble hold over the pure states of
\(\mathcal{H}_B \otimes \mathcal{H}_R\). Then the time-averaged
entanglement entropy of the radiation equals the Page value {[}32{]}:}

\[\overline{S(\rho_R)} = S_{\text{Page}}(d_B, d_R) + \mathcal{O}(\epsilon)\]

\emph{for all initial conditions except a fraction
\(\leq (\ln d_{\min}/\epsilon)\,e^{-c \, d_{\min} \, \epsilon^2}\),
where \(d_B = |\mathcal{C}_B(t)|\), \(d_R = |\mathcal{C}_R(t)|\),
\(d_{\min} = \min(d_B, d_R)\), and:}

\[S_{\text{Page}}(d_B, d_R) = \ln(\min(d_B, d_R)) - \frac{\min(d_B, d_R)}{2\max(d_B, d_R)}\]

\emph{the large-dimension form of Page's mean entropy for Haar-random
pure states of \(\mathcal{H}_B \otimes \mathcal{H}_R\). The initial
conditions counted are the configurations of \(\Sigma_E\), and the
bipartition \(B|R\) is aligned with the configuration coordinates. In
that setting, for \(d_B, d_R > 1\) and \(\epsilon < S_{\text{Page}}\),
H-scramble in the form the proof uses (at most a fraction
\(e^{-c\,d_{\min}\epsilon^2}\) of the states on the orbit atypical)
fails for every bijection, and the time-averaged entropy is zero for
every initial condition, so the theorem as stated has no instance (scope
below).}

\emph{Proof.} By hypothesis the energy is constant on
\(\mathcal{C}_B \times \mathcal{C}_R\), so energy conservation places no
constraint on the orbit and \(\varphi\) acts as a bijection of
\(\Sigma_E = \mathcal{C}_B \times \mathcal{C}_R\). Since \(\Sigma_E\) is
finite, this restricted bijection decomposes into disjoint cycles
\(C_1, \ldots, C_k\) with lengths \(L_1, \ldots, L_k\) summing to
\(|\Sigma_E|\). On each cycle \(C_j\), the time average of any
observable \(f\) is exact:

\[\overline{f}_{C_j} = \frac{1}{L_j} \sum_{s \in C_j} f(s)\]

By measure concentration for Haar-random pure states of
\(\mathcal{H}_B \otimes \mathcal{H}_R\) (Lévy's lemma applied to the
entanglement entropy, as in Hayden-Leung-Winter {[}61{]};
Popescu-Short-Winter {[}29{]} give the corresponding bound for the
distance of the reduced state from its canonical form), the set of
atypical states --- those with
\(|S(\text{Tr}_B |\psi\rangle\langle\psi|) - S_{\text{Page}}| > \epsilon\)
--- has measure at most \(\eta = e^{-c \, d_{\min} \, \epsilon^2}\),
where \(d_{\min} = \min(d_B, d_R)\) and \(c > 0\) is a universal
constant.

\textbf{The measures are different, and the step below requires a
hypothesis (H-scramble).} {[}29{]} and {[}61{]} bound a Haar fraction on
the unit sphere of a Hilbert space or subspace. They do not bound the
fraction of an arbitrary discrete set of substratum configurations, nor
of a permutation orbit, so ``the total number of atypical states is at
most \(\eta|\Sigma_E|\)'' does not follow from {[}29, 61{]} alone. The
gap is not technical. Let the shell be the computational product basis
\(\{|b\rangle\otimes|r\rangle\}\) at a common energy and let \(\varphi\)
be one long cycle permuting it: \(\varphi\) is finite and bijective,
conserves the shell, its orbit can contain every basis configuration,
its counting measure is exactly uniform, and yet every visited state is
a product state with \(S(\rho_R)=0\), while the Page value is
\(\ln d_{\min}-d_{\min}/(2d_{\max})>0\) for \(d_B,d_R>1\) --- a gap of
\(0.886\) nats already at \(d_B=d_R=4\). The same holds for every
bijection: acting on a definite configuration, \(\varphi\) visits only
configurations, each a product state across a bipartition aligned with
the configuration coordinates, so \(S(\rho_R) = 0\) along every such
orbit. So finite bijectivity, energy conservation and uniform cycle
counting together do not give Page typicality.

What the argument needs is a separate condition, \textbf{H-scramble}:
the orbit ensemble equidistributes over quantum pure states on the shell
--- in Theorem A.8, over the pure states of
\(\mathcal{H}_B \otimes \mathcal{H}_R\) --- in the sense of an
approximate unitary design or comparable scrambling property. Under
H-scramble the counting step below goes through with \(\eta|\Sigma_E|\)
as written, and it uses H-scramble only through that bound: at most
\(\eta|\Sigma_E|\) of the states on the orbit are atypical. A full cycle
of configurations is an exact state 1-design, so a design property of
low order does not supply the bound, and the failure stated next is a
failure of the bound. By the preceding paragraph, H-scramble fails for
every bijection in the setting of this theorem's own counting, a
definite initial configuration with the bipartition aligned to the
configuration coordinates (scope below), where every configuration is
atypical once \(\epsilon < S_{\text{Page}}\): it requires an initial
superposition of configurations or a bipartition not aligned with them,
and which of these the substratum supplies is open. H-scramble is not
established here, and the lesson generalizes: uniformity over classical
microstates is not uniformity over quantum pure states. The same
distinction governs H-state in §2.

Call a cycle \emph{bad} if its time-averaged entropy deviates from
\(S_{\text{Page}}\) by more than \(2\epsilon\). Every state's deviation
is bounded by \(\ln d_{\min}\), so a cycle with atypical fraction
\(f_j\) satisfies
\(|\overline{S}_{C_j} - S_{\text{Page}}| \leq \epsilon + f_j \ln d_{\min}\);
a bad cycle therefore requires \(f_j \geq \epsilon/\ln d_{\min}\),
i.e.~\(L_j \leq (\ln d_{\min}/\epsilon)\,|\{s \in C_j : s \text{ atypical}\}|\).
Summing over all bad cycles:

\[\sum_{j \,\text{bad}} L_j \leq \frac{\ln d_{\min}}{\epsilon}\,\eta\,|\Sigma_E|\]

The fraction of initial conditions on \(\Sigma_E\) that land on a bad
cycle is therefore at most
\((\ln d_{\min}/\epsilon)\,\eta = (\ln d_{\min}/\epsilon)\,e^{-c\,d_{\min}\epsilon^2}\).
Under H-scramble the time average is within \(2\epsilon\) of the Page
value outside this bounded fraction. \(\square\)

\textbf{Scope of Theorem A.8.} The proof counts initial configurations
\(s \in \Sigma_E\), with \(B|R\) the factorization
\(\mathcal{C}_B \times \mathcal{C}_R\) of the configurations. The orbit
of a configuration consists of configurations, each a product state
across that bipartition, so \(S(\rho_R) = 0\) at every step and the time
average is zero for every counted initial condition: H-scramble fails
for every bijection, and the theorem holds in that setting only because
its hypothesis is never met. A Page value at a bipartition aligned with
the configuration coordinates needs entropy in the initial state. Let
\(|\psi_0\rangle = \sum_s \psi_0(s)\,|s\rangle\) be a pure initial state
of \(B \cup R\) and \(p_0(s) = |\psi_0(s)|^2\) its distribution over
configurations. For every bijection \(\varphi\), every \(t\) and every
bipartition aligned with the configuration coordinates,

\[S(\rho_R(t)) \leq H(p_0) \leq \ln|\operatorname{supp}\psi_0|.\]

\(U_\varphi^t\) permutes the amplitudes, so the distribution \(p_t\)
over configurations at time \(t\) is a relabelling of \(p_0\) and
\(H(p_t) = H(p_0)\). The diagonal of \(\rho_R(t)\) in the configuration
basis of \(R\) is the marginal \(p_{t,R}\), and the diagonal of a
density matrix is majorized by its spectrum, so
\(S(\rho_R(t)) \leq H(p_{t,R}) \leq H(p_t)\), the second inequality
bounding a marginal entropy by the joint entropy; a distribution on
\(|\operatorname{supp}\psi_0|\) points has entropy at most
\(\ln|\operatorname{supp}\psi_0|\). A cycle average within \(\delta\) of
\(S_{\text{Page}}(d_B, d_R)\) at an aligned bipartition therefore
requires \(H(p_0) \geq S_{\text{Page}} - \delta\), hence an initial
superposition over at least \(e^{S_{\text{Page}} - \delta}\)
configurations; otherwise it requires a bipartition not aligned with the
configuration coordinates. Which of these, if either, the substratum
supplies is open. For initial superpositions the counting above, which
runs over configurations, does not apply as written; a restatement of
Theorem A.8 over an ensemble of initial pure states, with the excluded
fraction counted in that ensemble's measure, is not given here.

\textbf{Proper energy shells.} Theorem A.8 is stated for the full
product space. Its hypothesis requires the energy to be constant on
\(\mathcal{C}_B \times \mathcal{C}_R\), so that energy conservation
places no constraint on the orbit of \(B \cup R\); it does not follow
from the conservation of a Hamiltonian that takes more than one value
there, whose energy shells all span proper subspaces. An evaporating
black hole whose radiation carries away its energy is in that
proper-shell case. When energy conservation confines the orbit to a
shell whose span is a proper subspace of
\(\mathcal{H}_B \otimes \mathcal{H}_R\), measure concentration on that
subspace concentrates the entropy about the subspace's own Haar mean,
which depends on the shell and in general differs from
\(S_{\text{Page}}(d_B, d_R)\). For the diagonal shell spanned by
\(|i\rangle_B \otimes |i\rangle_R\), \(i = 1, \ldots, m\) (an energy
shell for an energy depending only on \(b - r\) and taking its value at
\(b = r\) nowhere else), the Schmidt probabilities of a Haar-random
state are uniformly distributed on the simplex and the mean entanglement
entropy is exactly \(H_m - 1\), with \(H_m\) the \(m\)-th harmonic
number: \(1/2\) at \(m = 2\), against Page's mean \(1/3\) for
\(d_B = d_R = 2\). For shells of fixed Hamming weight the value can lie
below \(S_{\text{Page}}(d_B, d_R)\). Theorem A.8 on a proper energy
shell, with the shell's own Haar mean in place of
\(S_{\text{Page}}(d_B, d_R)\), and the Page curve of Theorem A.9 on such
a shell, are open.

\textbf{Theorem A.9} (The cycle-averaged Page curve; conditional on A.8,
hence on H-scramble, and on H-decouple). \emph{At each stage \(t\) of
the evaporation, take the bipartition of Theorem A.8 with
\(d_B = q^{n_B(t)}\) and \(d_R = q^{n_R(t)}\), held fixed, and the
energy shell of Theorem A.8, the full configuration space of
\(B \cup R\) at that stage. For all but the fraction of initial
conditions excluded in Theorem A.8, and up to its
\(\mathcal{O}(\epsilon)\), the full-cycle time average of the
radiation's entanglement entropy at that bipartition is:}

\[\overline{S}_{\text{rad}}(t) = \begin{cases} n_R(t) \cdot \ln q - \frac{q^{n_R(t)}}{2\,q^{n_B(t)}} & \text{for } n_R(t) < n_B(t) \quad \text{(early: before Page time)} \\ n_B(t) \cdot \ln q - \frac{q^{n_B(t)}}{2\,q^{n_R(t)}} & \text{for } n_R(t) > n_B(t) \quad \text{(late: after Page time)} \end{cases}\]

\emph{where \(q\) is the alphabet size, so that the state space of one
mode has dimension \(q\), \(n_B(t) + n_R(t) = n_B(0)\), and the Page
time \(t_P\) is defined by \(n_R(t_P) = n_B(t_P) = n_B(0)/2\).}

\emph{Proof.} Direct substitution of \(d_B = q^{n_B(t)}\) and
\(d_R = q^{n_R(t)}\) into the Page entropy formula (Theorem A.8), using
\(\ln(q^n) = n \ln q\) and the conservation law \(n_B + n_R = n_B(0)\).
The turnover at \(n_R = n_B\) follows from \(\min(d_B, d_R)\) switching
from \(d_R\) to \(d_B\). \(\square\)

Theorem A.9 concerns full-cycle averages at the bipartition of stage
\(t\), held fixed. The entropy \(S(\rho_R(t))\) along a single
trajectory, whose bipartition moves with \(t\), is a different quantity,
and whether it follows the same curve is open. Theorem A.9 inherits the
scope of Theorem A.8: from an initial configuration, with the
bipartition of each stage aligned with the configuration coordinates,
\(S(\rho_R) = 0\) at every step, and by the bound in that scope the
curve requires, at each stage, an initial state whose distribution over
configurations has entropy at least \(\overline{S}_{\text{rad}}(t)\) up
to the tolerance, or a bipartition not aligned with the configuration
coordinates.

\textbf{The evaporation trajectory.} The Hawking temperature (§3.2, Step
4) gives luminosity \(L \sim 1/M^2\), hence \(dM/dt = -\alpha/M^2\)
where \(\alpha = \hbar c^4 \Gamma / (15360 \pi G^2)\) and \(\Gamma\) is
the species-dependent greybody factor (solution-specific). Integrating:

\[M(t) = M_0\!\left(1 - t/t_{\text{evap}}\right)^{1/3}, \qquad t_{\text{evap}} = \frac{M_0^3}{3\alpha}\]

The boundary mode count
\(n_B(t) = 4\pi G M(t)^2 / (\hbar c \sigma) = n_B(0) \cdot (1 - t/t_{\text{evap}})^{2/3}\)
and \(n_R(t) = n_B(0) - n_B(t)\). The Page time \(t_P\) satisfies
\(n_B(t_P) = n_B(0)/2\):

\[t_P = t_{\text{evap}}\!\left(1 - 2^{-3/2}\right) \approx 0.646 \; t_{\text{evap}}\]

independent of \(M_0\) and of \(\alpha\), hence of the overall greybody
normalization \(\Gamma\).

\textbf{Dependence on the radiation's entropy.} The conservation law
\(n_B + n_R = n_B(0)\) assigns the radiation one mode of dimension \(q\)
per boundary mode lost, so \(\ln d_R\) grows by exactly the decrease of
\(\ln d_B\). If instead
\(\ln d_R(t) = \beta\,[\ln d_B(0) - \ln d_B(t)]\) with constant
\(\beta > 0\), the turnover \(d_R = d_B\) of Theorem A.9 falls at
\(n_B(t_P)/n_B(0) = \beta/(1+\beta)\), and the trajectory above gives,
as in {[}60, Eq. (8){]},

\[t_P = t_{\text{evap}}\left[1 - \left(\frac{\beta}{1+\beta}\right)^{3/2}\right],\]

with \(\alpha\) and \(M_0\) cancelling for every \(\beta\). With
\(\sigma = \ln q\), \(\beta\) is the radiation's coarse-grained entropy
per unit of Bekenstein-Hawking entropy lost. In the quasi-stationary
semiclassical regime --- a non-rotating, uncharged black hole emitting
free fields into empty space (incoming vacuum), outgoing modes in
independent thermal-form states with greybody factors
\(0 < \Gamma_\omega \leq 1\), and \(dS_{\text{BH}} = dE/T_H\) --- every
emitted mode carries more entropy than its energy divided by \(T_H\)
{[}59, Eqs. (5)--(10){]}, so \(\beta > 1\) in that regime for every
greybody profile. A different lower bound gives the same per-mode
inequality: a mode of energy \(\omega\) transmitted with probability
\(0 < \Gamma_\omega \leq 1\) leaves in a thermal-form state, independent
of the other modes, with mean occupation
\(N = \Gamma_\omega\,n(\omega)\) and entropy \(g(N)\), with \(n\) and
\(g\) the thermal occupation at \(T_H\) and the entropy at fixed
occupation for Bose or Fermi statistics; since \(g\) is concave with
\(g(0) = 0\), \(g(N)/N \geq g(n)/n\) for \(N \leq n\), and
\(g(n) - n\,\omega/T_H\) equals \(\ln(1+n)\) for bosons and
\(-\ln(1-n)\) for fermions, both positive, so \(g(N) > N\omega/T_H\).
Unfiltered blackbody emission of massless quanta in 3+1 dimensions
(\(\Gamma_\omega = 1\), mode density \(\propto \omega^2\)) gives
\(\beta = 4/3\) and \(t_P \approx 0.568\,t_{\text{evap}}\). For a large
Schwarzschild black hole emitting only photons and gravitons, at
constant \(\alpha\), Page's emission rates give \(\beta \approx 1.485\)
and \(t_P \approx 0.538\,t_{\text{evap}}\) {[}60, Eqs. (2), (8){]}; the
value for a black hole depends on the species content and the frequency
dependence of \(\Gamma_\omega\). The value
\(t_P \approx 0.646\,t_{\text{evap}}\) is the Page time at
\(\beta = 1\).

At \(\beta = 1\), the cycle-averaged entropy of the radiation at the
Page time is
\(\overline{S}_{\text{rad}}(t_P) \approx n_B(0) \ln q / 2\), half the
entropy \(n_B(0) \ln q\) of the initial boundary modes, up to the
finite-size correction of Theorem A.9 (\(-1/2\) at \(n_R = n_B\)). With
\(n_B(0) = S_{\text{BH}}(0)/\sigma\) this is
\(\overline{S}_{\text{rad}}(t_P) \approx S_{\text{BH}}(0) \ln q/(2\sigma)\):
approximately half the initial black-hole entropy when
\(\sigma = \ln q\), the entropy of a mode's state space, which gives
\(\epsilon = 2\sqrt{\ln q}\,l_p\) (§4); at \(\sigma = 1\), which the
bound \(\sigma \leq \ln q\) excludes for \(q = 2\), the ratio is
\(\ln q/2\).

\textbf{Remark (Comparison with QES + islands).} The quantum extremal
surface formula and the entanglement-island construction {[}41, 42, 43,
44{]} obtain the Page curve by a different route. In the QES approach,
the Page curve emerges through an extremization principle that selects
the dominant saddle in a replica-trick computation, with the
post-Page-time saddle including an ``island'' --- a region of bulk
spacetime that contributes to the radiation's entanglement entropy
despite being topologically disconnected. The replica wormhole
geometries that connect the replicas in the gravitational path integral
are essential to obtaining the saturating Page curve at late times.

The derivation here obtains a cycle-averaged Page curve of the same form
through the cycle-typicality argument of Theorem A.8, which applies
measure concentration {[}29, 61{]} to the deterministic substratum
dynamics on the full configuration space of \(B \cup R\) under
H-scramble, with no replica trick and no gravitational path integral;
H-scramble is an equidistribution hypothesis on the orbit over pure
states, which ergodicity over configurations does not supply. In the QES
computation the turnover occurs where the island saddle's generalized
entropy, at leading order the remaining Bekenstein-Hawking entropy,
falls below the semiclassical entropy of the radiation. The two Page
times coincide when that semiclassical entropy grows one-for-one with
the horizon entropy lost, the \(\beta = 1\) of the mode counting;
\(\beta > 1\) in the semiclassical regime defined above ({[}59{]}; for
the QES setting, {[}41{]}, §2), so agreement on the Page time and the
Page entropy is not established. The ``island'' in the QES formulation
is conjectured to correspond, in the present construction, to the
post-Page-time regime in which the BH boundary modes form the smaller
subsystem, and the ``replica wormhole'' geometries to the substratum's
cycle structure on the energy shell. This correspondence --- replica
wormholes ↔ substratum cycles --- is conjectural; making it precise
would require constructing the gravitational path integral as an
effective description of the substratum dynamics in the appropriate
limit. What the present derivation establishes, at the G1 register, is
conditional: under H-scramble and H-decouple, with the full product
space as the energy shell, a cycle-averaged Page curve with turnover at
\(d_R = d_B\) follows without the replica machinery. The Page value
enters through H-scramble, not through the finite deterministic
dynamics: from an initial configuration no bijection produces
entanglement across a bipartition aligned with the configuration
coordinates, and from any initial pure state it produces at most the
entropy of that state's distribution over configurations (§A.7).

Both approaches require \emph{something} to play the role of typicality
(measure concentration on the full product space, carried to the orbit
by H-scramble, here; replica saddles in QES + islands), and both place
the turnover where the radiation's entropy reaches the black hole's
remaining entropy. If the conjectured correspondence holds, the
typicality-providing ingredient need not be gravitational, and the
derivation here would be a microscopic realization of the
entanglement-island construction; on the substratum side the bijection
alone does not supply that ingredient, since the entanglement it can
produce across an aligned bipartition is bounded by the entropy of the
initial distribution over configurations (§A.7).

\begin{center}\rule{0.5\linewidth}{0.5pt}\end{center}

\hypertarget{appendix-b-the-deep-sector-corollary-for-a-deep-sector-of-any-cardinality}{%
\subsection{Appendix B: The deep-sector corollary for a deep sector of
any
cardinality}\label{appendix-b-the-deep-sector-corollary-for-a-deep-sector-of-any-cardinality}}

This appendix proves the statements of the corollary of §3.2
(deep-sector cardinality is gauge) that do not use H-BD, for a deep
sector of any cardinality.

\textbf{Setting.} The sites form a set \(\iota\) of any cardinality;
site \(i\) carries a nonempty alphabet \(Q_i\), and
\(\mathcal{C} = \prod_{i \in \iota} Q_i\). For \(A \subseteq \iota\),
\(\mathcal{C}_A = \prod_{i \in A} Q_i\), and \(s|_A\) is the restriction
of \(s \in \mathcal{C}\) to \(A\). The update is a map
\(\varphi : \mathcal{C} \to \mathcal{C}\) with components \(\varphi_i\).
The visible sites form \(V \subseteq \iota\), and
\(\mathcal{C}_H = \mathcal{C}_{\iota \setminus V}\). The dependence set
\(\mathrm{Dep}(i)\) is the set of sites \(j\) for which some
\(s, s' \in \mathcal{C}\) differing only at \(j\) have
\(\varphi_i(s) \neq \varphi_i(s')\); the coupling graph \(G_\varphi\)
joins \(i \neq j\) when \(j \in \mathrm{Dep}(i)\) or
\(i \in \mathrm{Dep}(j)\), as in {[}Main §3.3{]}, and
\(B_n = B_{G_\varphi}(V, n) \setminus V\) and
\(D_n = \iota \setminus B_{G_\varphi}(V, n)\), as in part (a) of the
lemma of §3.2. For a hidden prior \(\mu_H\), the same for every visible
initial configuration,
\(T_{ij}(m\,\Delta t) = \mu_H(\{h \in \mathcal{C}_H : \pi_V(\varphi^m(x_i, h)) = x_j\})\);
for the uniform prior on a finite \(\mathcal{C}_H\) this is the count of
§3.2. \(T^{(B_n)}_{ij}(m\,\Delta t)\) is the sum of part (a) of the
lemma with \(\varphi^m\) in place of \(\varphi^n\), at a given
\(d^* \in \mathcal{C}_{D_n}\).

\textbf{Hypotheses.}

\begin{itemize}
\tightlist
\item
  (D1) \(V\) is a finite set of sites.
\item
  (D2) \(\varphi\) is defined on every configuration, and each
  component's update is a function of finitely many components: for
  every site \(i\) there is a finite set \(R(i) \subseteq \iota\) with
  \(\varphi_i(s) = \varphi_i(s')\) whenever \(s|_{R(i)} = s'|_{R(i)}\).
\item
  (D3) \(G_\varphi\) is locally finite.
\item
  (D4) \(Q_i\) is finite for every \(i \in V \cup B_n\).
\item
  (D5) The hidden prior is \(\mu_H = u_{B_n} \otimes \nu\) on
  \(\mathcal{C}_H = \mathcal{C}_{B_n} \times \mathcal{C}_{D_n}\), where
  \(u_{B_n}\) is the uniform distribution on \(\mathcal{C}_{B_n}\) and
  \(\nu\) is a probability measure on a \(\sigma\)-algebra
  \(\mathcal{D}\) on \(\mathcal{C}_{D_n}\): a set
  \(E \subseteq \mathcal{C}_H\) is measurable when each section
  \(E_b = \{d : (b, d) \in E\}\) lies in \(\mathcal{D}\), and then
  \(\mu_H(E) = |\mathcal{C}_{B_n}|^{-1} \sum_{b} \nu(E_b)\).
\end{itemize}

On a finite site set (D1)--(D4) hold, and the uniform prior on
\(\mathcal{C}_H\) is (D5) with \(\nu\) uniform. (D2) and (D3) are
hypotheses on the dynamics, which hold on every finite site set; for the
second-order update of ledger item (xiii) of {[}Main §4.5{]} they follow
from that item's clauses that each \(F_i\) reads finitely many sites and
each site is read by finitely many \(F_i\), since the update at \(i\)
reads only the pair at \(i\) and the current values that \(F_i\) reads
(with Lemma B.1 for (D3)), and (D4) holds there because each site
carries a pair of values in a finite abelian group. (D1) is a condition
on the choice of the visible sector, which item (xiii) does not fix.
(D5) is a hypothesis on the prior of the model, not a statement about
nature's prior; \(\nu\) need not be a product over the sites of \(D_n\).

\textbf{The dependency patch.} With \(N_i = \mathrm{Dep}(i)\), put
\(N^0(V) = V\), \(N^{k+1}(V) = \bigcup_{i \in N^k(V)} N_i\),
\(U_n = \bigcup_{k \leq n} N^k(V)\) and \(W_n = U_n \setminus V\):
\(U_n\) is the set of sites from which a chain of at most \(n\)
dependences leads to \(V\). The cone \(B_n\) is defined by distance in
the symmetrized graph \(G_\varphi\), so it can also contain sites whose
updates depend on sites of \(V \cup B_{n-1}\) while no site of
\(V \cup B_{n-1}\) depends on them. By Lemma B.2, \(W_n \subseteq B_n\),
and the inclusion can be strict for a bijection: on bits at sites
\(v, b, k\) with \(V = \{v\}\) and
\(\varphi(v, b, k) = (b, v, k \oplus v)\), \(W_n = \{b\}\) while
\(B_n = \{b, k\}\) for every \(n \geq 1\). With \(n = 1\), deleting the
site \(k\) leaves the update at the site of \(V \cup B_0 = V\) unchanged
and changes \(B_1\) to \(\{b\}\), so \(B_n\) is not in general
determined by the update on \(V \cup B_{n-1}\). The proofs below use
\(W_n\).

\textbf{Lemma B.1} (dependence sets). \emph{Under (D2), for every site
\(i\): \(\mathrm{Dep}(i) \subseteq R(i)\), so \(\mathrm{Dep}(i)\) is
finite; \(\varphi_i(s) = \varphi_i(s')\) whenever
\(s|_{\mathrm{Dep}(i)} = s'|_{\mathrm{Dep}(i)}\), so that
\(\varphi_i(s) = \hat\varphi_i(s|_{\mathrm{Dep}(i)})\) for a unique
local rule \(\hat\varphi_i : \mathcal{C}_{\mathrm{Dep}(i)} \to Q_i\);
and every \(j \in \mathrm{Dep}(i)\) other than \(i\) is adjacent to
\(i\) in \(G_\varphi\).}

\emph{Proof.} If \(j \notin R(i)\), two configurations differing only at
\(j\) agree on \(R(i)\), so \(j \notin \mathrm{Dep}(i)\). Let \(s, s'\)
agree on \(\mathrm{Dep}(i)\), list the finite set
\(R(i) \setminus \mathrm{Dep}(i)\) as \(j_1, \ldots, j_r\), and let
\(s^0 = s\) and \(s^l\) be \(s^{l-1}\) with its \(j_l\)th component
replaced by \(s'_{j_l}\). Consecutive terms differ at most at
\(j_l \notin \mathrm{Dep}(i)\), so
\(\varphi_i(s^l) = \varphi_i(s^{l-1})\), and \(s^r\) agrees with \(s'\)
on \(R(i)\), so \(\varphi_i(s) = \varphi_i(s^r) = \varphi_i(s')\). Every
\(u \in \mathcal{C}_{\mathrm{Dep}(i)}\) extends to a configuration, the
alphabets being nonempty, which defines \(\hat\varphi_i\). The last
clause is the definition of \(G_\varphi\). \(\square\)

\textbf{Lemma B.2} (dependency patch). \emph{Assume (D1) and (D2), and
let \(n \geq 1\). (a) \(U_n\) is finite and
\(U_n \subseteq B_{G_\varphi}(V, n)\); hence \(W_n \subseteq B_n\) and
\(U_{n-1} \subseteq V \cup B_{n-1}\). Under (D3), \(B_n\) is finite. (b)
There is a map \(G_n : \mathcal{C}_{U_n} \to (\mathcal{C}_V)^{n+1}\)
with \(\pi_V(\varphi^m(s)) = G_n(s|_{U_n})_m\) for every
\(s \in \mathcal{C}\) and \(m \leq n\). (c) \(G_n\) is determined by the
set \(U_n\) with its alphabets and by \(\mathrm{Dep}(i)\) and
\(\hat\varphi_i\) for \(i \in U_{n-1}\). Let a second system
\((\iota', Q', \varphi')\) with the same \(V\) and visible alphabets
satisfy (D2), and let \(e : U_n \to \iota'\) be an injection that is the
identity on \(V\), with \(Q'_{e(k)} = Q_k\) for \(k \in U_n\), and with
\(\mathrm{Dep}'(e(i)) = e(\mathrm{Dep}(i))\) and
\(\hat\varphi'_{e(i)}(u \circ e^{-1}) = \hat\varphi_i(u)\) for
\(i \in U_{n-1}\) and \(u \in \mathcal{C}_{\mathrm{Dep}(i)}\). Then
\(e(U_n) = U'_n\), and \(\pi_V(\varphi'^m(s')) = G_n(s' \circ e)_m\) for
every \(s' \in \mathcal{C}'\) and \(m \leq n\), where
\((s' \circ e)_k = s'_{e(k)}\).}

\emph{Proof.} (a) By Lemma B.1, each \(\varphi_i\) is a function of the
components in the finite set \(N_i\), and every \(j \in N_i\) other than
\(i\) is adjacent to \(i\) in \(G_\varphi\). The coupling-graph causal
cone of {[}Main §3.3{]} for a site set of any cardinality, applied with
these reading sets, with \(R = V\), finite by (D1), and with
\(\Gamma = G_\varphi\), gives that each \(N^k(V)\) is finite and lies
within \(G_\varphi\)-distance \(k\) of \(V\); its argument uses no
finiteness of the alphabets. Hence \(U_n\) is finite and
\(U_n \subseteq B_{G_\varphi}(V, n)\), which gives the inclusions. Under
(D3), \(B_{G_\varphi}(V, k+1)\) is \(B_{G_\varphi}(V, k)\) together with
the finitely many neighbours of its sites, so it is finite. (b) By the
same result, for \(m \leq n\) two configurations that agree on
\(N^m(V) \subseteq U_n\) have \(m\)-step images that agree on \(V\);
since every element of \(\mathcal{C}_{U_n}\) extends to a configuration,
the alphabets being nonempty, this defines \(G_n\). (c) For
\(0 \leq k \leq n\) there are maps
\(g_k : \mathcal{C}_{U_n} \to \mathcal{C}_{U_{n-k}}\) with
\((\varphi^k s)|_{U_{n-k}} = g_k(s|_{U_n})\), obtained recursively from
the local rules. Take \(g_0\) the identity. Given \(g_k\) with
\(k < n\), let \(i \in U_{n-k-1}\); then \(i \in N^l(V)\) for some
\(l \leq n-k-1\), so
\(\mathrm{Dep}(i) \subseteq N^{l+1}(V) \subseteq U_{n-k}\), and
\((\varphi^{k+1} s)_i = \hat\varphi_i((\varphi^k s)|_{\mathrm{Dep}(i)}) = \hat\varphi_i(g_k(s|_{U_n})|_{\mathrm{Dep}(i)})\),
which defines \(g_{k+1}\). Since \(V \subseteq U_{n-k}\),
\(G_n = (g_m(\cdot)|_V)_{m \leq n}\), and the recursion uses only the
data listed, at the sites of \(U_{n-k-1} \subseteq U_{n-1}\). Write
\(N'^m(V)\) for the sets formed from \(\mathrm{Dep}'\) in the second
system. By induction on \(m\), \(e(N^m(V)) = N'^m(V)\) for \(m \leq n\):
for \(m < n\), \(N^m(V) \subseteq U_{n-1}\), so
\(N'^{m+1}(V) = \bigcup_{i \in N^m(V)} \mathrm{Dep}'(e(i)) = e(N^{m+1}(V))\).
Hence \(e(U_n) = U'_n\), and the recursion in the second system, carried
out on \(e(U_n)\) with the transported data, produces the maps of the
first system composed with \(s' \mapsto s' \circ e\). \(\square\)

\textbf{Lemma B.3} (marginal on a sub-patch). \emph{Let \(B\) be a
finite set of sites with finite alphabets, \(W \subseteq B\),
\((\mathcal{C}_D, \mathcal{D}, \nu)\) a probability space and
\(\mu = u_B \otimes \nu\) on \(\mathcal{C}_B \times \mathcal{C}_D\), as
in (D5). For every \(F \subseteq \mathcal{C}_W\) the set
\(\{(b, d) : b|_W \in F\}\) is measurable and has \(\mu\)-measure
\(|F|/|\mathcal{C}_W|\).}

\emph{Proof.} The set is \(E \times \mathcal{C}_D\) with
\(E = \{b \in \mathcal{C}_B : b|_W \in F\}\); each section is
\(\mathcal{C}_D\) or empty, and
\(\mu(E \times \mathcal{C}_D) = |E|\,\nu(\mathcal{C}_D)/|\mathcal{C}_B| = |E|/|\mathcal{C}_B|\).
Since \(|E| = |F|\,|\mathcal{C}_{B \setminus W}|\) and
\(|\mathcal{C}_B| = |\mathcal{C}_W|\,|\mathcal{C}_{B \setminus W}|\),
the measure is \(|F|/|\mathcal{C}_W|\). \(\square\)

\emph{Proof of the statements of the corollary that do not use H-BD.}
Assume (D1)--(D5) and \(n \geq 1\). By (D3), (D4) and Lemma B.2 (a),
\(B_n\) and \(W_n \subseteq B_n\) are finite sets of sites with finite
alphabets. Fix \(x_i, x_j \in \mathcal{C}_V\) and \(m \leq n\), and put
\(F = \{w \in \mathcal{C}_{W_n} : G_n(x_i, w)_m = x_j\}\). Since
\(U_n \subseteq V \cup B_n\), Lemma B.2 (b) gives
\(\pi_V(\varphi^m(x_i, b, d)) = G_n(x_i, b|_{W_n})_m\) for
\(b \in \mathcal{C}_{B_n}\) and \(d \in \mathcal{C}_{D_n}\), so the
event defining \(T_{ij}(m\,\Delta t)\) is
\(\{(b, d) : b|_{W_n} \in F\}\), and Lemma B.3 with \(B = B_n\) and
\(W = W_n\) gives

\[T_{ij}(m\,\Delta t) = \frac{|F|}{|\mathcal{C}_{W_n}|}.\]

For each \(d^* \in \mathcal{C}_{D_n}\) the number of
\(b \in \mathcal{C}_{B_n}\) with \(b|_{W_n} \in F\) is
\(|F|\,|\mathcal{C}_{B_n \setminus W_n}|\), so
\(T^{(B_n)}_{ij}(m\,\Delta t)\) has the same value. Neither
\(|\mathcal{C}_{D_n}|\) nor \(\nu\) enters, and the update enters only
through \(G_n\), which by Lemma B.2 (c) is fixed by the update on
\(U_{n-1} \subseteq V \cup B_{n-1}\). Let a second system satisfy
(D1)--(D5), with the same \(V\) and visible alphabets, and have the same
update on the sites of \(V \cup B_{n-1}\) in the sense of the corollary.
Since \(U_{n-1} \subseteq V \cup B_{n-1}\) and
\(U_n = V \cup \bigcup_{i \in U_{n-1}} \mathrm{Dep}(i)\), the
restriction of the identification to \(U_n\) is a map \(e\) as in Lemma
B.2 (c). Hence \(W'_n = e(W_n) \subseteq B'_n\), the second system's
event is \(\{(b', d') : (b' \circ e)|_{W_n} \in F\}\) with the same
\(F\), and Lemma B.3 in the second system, with \(B = B'_n\) and
\(W = W'_n\), gives
\(T'_{ij}(m\,\Delta t) = |F|/|\mathcal{C}_{W_n}| = T_{ij}(m\,\Delta t)\).
The two systems may differ in the cardinalities, alphabets and updates
of their hidden sites outside the patch, in their cones \(B_n\) and
\(B'_n\), and in their priors on \(\mathcal{C}_{D_n}\) and
\(\mathcal{C}'_{D_n}\). \(\square\)

\begin{center}\rule{0.5\linewidth}{0.5pt}\end{center}

\hypertarget{acknowledgments}{%
\subsection{Acknowledgments}\label{acknowledgments}}

AI assistance in preparing this work is disclosed at\\
\url{https://github.com/amaybaum/incompleteness\#ai-assistance}.

\begin{center}\rule{0.5\linewidth}{0.5pt}\end{center}

\hypertarget{references}{%
\subsection{References}\label{references}}

\textbf{Companion papers (cited inline by short name):}

{[}Main{]} A. Maybaum, ``The Incompleteness of Observation,'' (2026).

{[}SM{]} A. Maybaum, ``The Standard Model from a Cubic Lattice,''
(2026).

{[}Substratum{]} A. Maybaum, ``The Substratum Construction:
Reconstruction, the Substratum Gauge Group, and the QM-GR Synthesis,''
(2026).

\begin{center}\rule{0.5\linewidth}{0.5pt}\end{center}

{[}1{]} T. Jacobson, ``Thermodynamics of Spacetime: The Einstein
Equation of State,'' \emph{Phys. Rev.~Lett.} \textbf{75}, 1260 (1995).

{[}2{]} G. W. Gibbons and S. W. Hawking, ``Cosmological event horizons,
thermodynamics, and particle creation,'' \emph{Phys. Rev.~D}
\textbf{15}, 2738 (1977).

{[}3{]} W. G. Unruh, ``Sonic analogue of black holes and the effects of
high frequencies on black hole evaporation,'' \emph{Phys. Rev.~D}
\textbf{51}, 2827 (1995).

{[}4{]} K. Fredenhagen and R. Haag, ``On the derivation of Hawking
radiation associated with the formation of a black hole,'' \emph{Commun.
Math. Phys.} \textbf{127}, 273--284 (1990).

{[}5{]} LIGO Scientific Collaboration, Virgo Collaboration, and KAGRA
Collaboration: A. G. Abac et al., ``GW250114: Testing Hawking's Area Law
and the Kerr Nature of Black Holes,'' \emph{Phys. Rev.~Lett.}
\textbf{135}, 111403 (2025).

{[}6{]} J. Martin, ``Everything you always wanted to know about the
cosmological constant problem (but were afraid to ask),'' \emph{C. R.
Phys.} \textbf{13}, 566--665 (2012).

{[}7{]} S. Weinberg, ``The cosmological constant problem,''
\emph{Rev.~Mod. Phys.} \textbf{61}, 1 (1989).

{[}8{]} S. M. Carroll, ``The Cosmological Constant,'' \emph{Living
Rev.~Relativ.} \textbf{4}, 1 (2001).

{[}9{]} D. H. Wolpert, ``Physical limits of inference,'' \emph{Physica
D} \textbf{237}, 1257--1281 (2008).

{[}10{]} J. Solà Peracaula, ``The cosmological constant problem and
running vacuum in the expanding universe,'' \emph{Phil. Trans. R. Soc.
A} \textbf{380}, 20210182 (2022).

{[}11{]} J. Solà Peracaula, A. Gómez-Valent, J. de Cruz Pérez, and C.
Moreno-Pulido, ``Running Vacuum in the Universe: Phenomenological Status
in Light of the Latest Observations, and Its Impact on the σ₈ and H₀
Tensions,'' \emph{Universe} \textbf{9}, 262 (2023); arXiv:2304.11157.

{[}12{]} DESI Collaboration, ``DESI 2024 VI: Cosmological Constraints
from BAO,'' arXiv:2404.03002 (2024).

{[}13{]} DESI Collaboration, ``DESI DR2 Results II: Measurements of
Baryon Acoustic Oscillations and Cosmological Constraints,'' \emph{Phys.
Rev.~D} \textbf{112}, 083515 (2025); arXiv:2503.14738.

{[}14{]} J. de Cruz Pérez, A. Gómez-Valent, and J. Solà Peracaula,
``Dynamical Dark Energy models in light of the latest observations,''
arXiv:2512.20616 (2025).

{[}15{]} N. R. Bertini, M. H. Novaes, R. von Marttens, and I. L.
Shapiro, ``Renormalization group corrections to ΛCDM model and
observational consequences for H₀ tension,'' arXiv:2509.24026 (2025).

{[}16{]} A. Gómez-Valent and J. Solà Peracaula, ``Relaxing the
σ₈-tension through running vacuum in the Universe,'' \emph{Mon. Not.
Roy. Astron. Soc.} \textbf{478}, 126--145 (2018); arXiv:1801.08501.

{[}17{]} DESI Collaboration: W. Elbers et al., ``Constraints on Neutrino
Physics from DESI DR2 BAO and DR1 Full Shape,'' arXiv:2503.14744 (2025).

{[}18{]} JUNO Collaboration: A. Abusleme et al., ``First measurement of
reactor neutrino oscillations at JUNO,'' arXiv:2511.14593 (2025); Nature
654, 343 (2026).

{[}19{]} F. Capozzi et al., ``Updated bounds on the (1,2) neutrino
oscillation parameters after first JUNO results,'' arXiv:2511.21650
(2025).

{[}20{]} E. P. Verlinde, ``Emergent Gravity and the Dark Universe,''
\emph{SciPost Phys.} \textbf{2}, 016 (2017); arXiv:1611.02269.

{[}21{]} LZ Collaboration, 417-day low-mass search, arXiv:2512.08065
(2025).

{[}22{]} S. S. McGaugh, F. Lelli, and J. M. Schombert, ``Radial
acceleration relation in rotationally supported galaxies,'' \emph{Phys.
Rev.~Lett.} \textbf{117}, 201101 (2016).

{[}23{]} R. Genzel et al., ``Strongly baryon-dominated disk galaxies at
the peak of galaxy formation ten billion years ago,'' \emph{Nature}
\textbf{543}, 397--401 (2017).

{[}24{]} R. Genzel et al., ``Rotation curves in z \textasciitilde{} 1--2
star-forming disks: evidence for cored dark matter distributions,''
\emph{Astrophys. J.} \textbf{902}, 98 (2020).

{[}25{]} Y. Jiao et al., ``Detection of the Keplerian decline in the
Milky Way rotation curve,'' \emph{Astronomy \& Astrophysics}
\textbf{678}, A208 (2023).

\begin{center}\rule{0.5\linewidth}{0.5pt}\end{center}

{[}26{]} N. Bao, S. M. Carroll, and A. Singh, ``The Hilbert space of
quantum gravity is locally finite-dimensional,'' \emph{Int. J. Mod.
Phys. D} \textbf{26}, 1743013 (2017).

{[}27{]} C. Bäcker, K. Beyer, and W. T. Strunz, ``Local disclosure of
quantum memory in non-Markovian dynamics,'' \emph{Phys. Rev.~Lett.}
\textbf{132}, 060402 (2024); arXiv:2310.01205.

{[}28{]} A. Yosifov, A. Iyer, V. Vedral, and J. Sun, ``On the emergence
of quantum memory in non-Markovian dynamics,'' arXiv:2507.21907 (2025).

{[}29{]} S. Popescu, A. J. Short, and A. Winter, ``Entanglement and the
foundations of statistical mechanics,'' \emph{Nature Physics}
\textbf{2}, 754--758 (2006).

{[}30{]} W. F. Stinespring, ``Positive functions on C\emph{-algebras,''
}Proc. Amer. Math. Soc.* \textbf{6}, 211--216 (1955).

{[}31{]} M. A. Nielsen and I. L. Chuang, \emph{Quantum Computation and
Quantum Information} (Cambridge University Press, 10th anniversary
edition, 2010).

{[}32{]} D. N. Page, ``Average entropy of a subsystem,'' \emph{Phys.
Rev.~Lett.} \textbf{71}, 1291 (1993).

{[}33{]} I. L. Shapiro and J. Solà, ``Scaling behavior of the
cosmological constant: interface between quantum field theory and
cosmology,'' \emph{JHEP} \textbf{02}, 006 (2002); arXiv:hep-th/0012227.

{[}34{]} C. Moreno-Pulido and J. Solà Peracaula, ``Running vacuum in QFT
in FLRW spacetime: the dynamics of ρ\_vac(H) from the quantized matter
fields,'' \emph{Eur. Phys. J. C} \textbf{83}, 637 (2023);
arXiv:2302.04807.

{[}35{]} J. A. Barandes, ``The Stochastic-Quantum Correspondence,''
\emph{Philosophy of Physics} \textbf{3}(1), 8 (2025), DOI
10.31389/pop.186; arXiv:2302.10778.

{[}36{]} J. A. Barandes, ``The Stochastic-Quantum Theorem,''
arXiv:2309.03085 (2023).

{[}37{]} D.-C. Dai and D. Stojkovic, ``Inconsistencies in Verlinde's
emergent gravity,'' \emph{JHEP} \textbf{11}, 007 (2017);
arXiv:1710.00946.

{[}38{]} J. Maldacena, ``Real observers solving imaginary problems,''
arXiv:2412.14014 (2024).

{[}39{]} V. Chandrasekaran, R. Longo, G. Penington, E. Witten, ``An
algebra of observables for de Sitter space,'' \emph{JHEP} \textbf{02}
(2023) 082; arXiv:2206.10780.

{[}40{]} D. Harlow, M. Usatyuk, Y. Zhao, ``Quantum mechanics and
observers for gravity in a closed universe,'' \emph{JHEP} \textbf{02}
(2026) 108; arXiv:2501.02359.

{[}41{]} G. Penington, ``Entanglement Wedge Reconstruction and the
Information Paradox,'' \emph{JHEP} \textbf{09} (2020) 002;
arXiv:1905.08255.

{[}42{]} A. Almheiri, N. Engelhardt, D. Marolf, H. Maxfield, ``The
entropy of bulk quantum fields and the entanglement wedge of an
evaporating black hole,'' \emph{JHEP} \textbf{12} (2019) 063;
arXiv:1905.08762.

{[}43{]} G. Penington, S. H. Shenker, D. Stanford, Z. Yang, ``Replica
wormholes and the black hole interior,'' \emph{JHEP} \textbf{03} (2022)
205; arXiv:1911.11977.

{[}44{]} A. Almheiri, T. Hartman, J. Maldacena, E. Shaghoulian, A.
Tajdini, ``Replica wormholes and the entropy of Hawking radiation,''
\emph{JHEP} \textbf{05} (2020) 013; arXiv:1911.12333.

{[}45{]} Fermi-LAT Collaboration, ``Search for Gamma-Ray Emission from
Local Primordial Black Holes with the Fermi Large Area Telescope,''
\emph{Astrophys. J.} \textbf{857}, 49 (2018); arXiv:1802.00100.

{[}46{]} HAWC Collaboration, ``Constraining the Local Burst Rate Density
of Primordial Black Holes with HAWC,'' \emph{JCAP} \textbf{04} (2020)
026; arXiv:1911.04356.

{[}47{]} Z. Cao et al.~(LHAASO Collaboration), ``All-Sky Search for
Individual Primordial Black Hole Bursts with LHAASO,'' \emph{Phys.
Rev.~Lett.} \textbf{135}, 181005 (2025); arXiv:2505.24586.

{[}48{]} C. Yang, J.-D. Pan, X. Zhang, ``Prospective sensitivity of CTAO
on detection of evaporating primordial black holes,'' \emph{JCAP}
\textbf{09} (2025) 004; arXiv:2504.17478.

{[}49{]} C. Moreno-Pulido, J. Solà Peracaula and S. Cheraghchi,
``Running vacuum in QFT in FLRW spacetime: the dynamics of
\(\rho_{\rm vac}(H)\) from the quantized matter fields,''
arXiv:2301.05205.

{[}50{]} E. P. Verlinde, ``On the Origin of Gravity and the Laws of
Newton,'' \emph{JHEP} \textbf{04} (2011) 029; arXiv:1001.0785.

{[}51{]} T. Padmanabhan, ``Thermodynamical Aspects of Gravity: New
insights,'' \emph{Rep.~Prog. Phys.} \textbf{73} (2010) 046901;
arXiv:0911.5004.

{[}52{]} G. Bianconi, ``Gravity from entropy,'' \emph{Phys. Rev.~D}
\textbf{111}, 066001 (2025); arXiv:2408.14391.

{[}53{]} Y. Ollivier, ``Ricci curvature of Markov chains on metric
spaces,'' \emph{J. Funct. Anal.} \textbf{256} (2009) 810--864.

{[}54{]} J. Ambjørn, J. Jurkiewicz and R. Loll, ``Spectral dimension of
the universe,'' \emph{Phys. Rev.~Lett.} \textbf{95}, 171301 (2005);
arXiv:hep-th/0505113.

{[}55{]} T. Breuer, ``The impossibility of accurate state
self-measurements,'' \emph{Philosophy of Science} \textbf{62}, 197--214
(1995).

{[}56{]} D. N. Page and C. D. Geilker, ``Indirect evidence for quantum
gravity,'' \emph{Phys. Rev.~Lett.} \textbf{47}, 979--982 (1981).

{[}57{]} S. Bose, A. Mazumdar, G. W. Morley, H. Ulbricht, M. Toroš, M.
Paternostro, A. A. Geraci, P. F. Barker, M. S. Kim and G. Milburn,
``Spin entanglement witness for quantum gravity,'' \emph{Phys.
Rev.~Lett.} \textbf{119}, 240401 (2017); arXiv:1707.06050.

{[}58{]} C. Marletto and V. Vedral, ``Gravitationally induced
entanglement between two massive particles is sufficient evidence of
quantum effects in gravity,'' \emph{Phys. Rev.~Lett.} \textbf{119},
240402 (2017); arXiv:1707.06036.

{[}59{]} D. N. Page, ``Hawking radiation and black hole
thermodynamics,'' \emph{New J. Phys.} \textbf{7}, 203 (2005);
arXiv:hep-th/0409024.

{[}60{]} D. N. Page, ``Time dependence of Hawking radiation entropy,''
\emph{JCAP} \textbf{09} (2013) 028; arXiv:1301.4995.

{[}61{]} P. Hayden, D. W. Leung, and A. Winter, ``Aspects of generic
entanglement,'' \emph{Commun. Math. Phys.} \textbf{265}, 95--117 (2006);
arXiv:quant-ph/0407049.

\end{document}
% source-sha256: e612c2b3335b67b6639812bc4caf3aa75cb58af3fee895ce50ae7396ddd4e46a
